% Lesson 31 — Reading: Texts about environmental issues like climate change and global warming — Anglais Tle A — Cours signé OnBuch+
% Programme officiel MINESEC. Compiler avec : tectonic EN31-reading-texts-about-environmental-issues-like-climate-c.tex
\newcommand{\DOCMATIERE}{Anglais}
\newcommand{\DOCNIVEAU}{Tle A}
\newcommand{\DOCMODULE}{Chapter 3 : Environment, Well-being and Health}
\newcommand{\DOCLECON}{EN31}
\newcommand{\DOCTITRE}{Lesson 31 — Reading: Texts about environmental issues like climate change and global warming}
% OnBuch+ — préambule « cours signés » ENRICHI (compilé avec tectonic / XeLaTeX)
% Système visuel « Pop » d'OnBuch+ : fond crème (jamais blanc), bordures encre
% épaisses, ombres « dures » décalées, titres Archivo Black en capitales.
% Version MATHÉMATIQUES : graphes (pgfplots/tikz), boîtes pédagogiques
% (définition, propriété, méthode, attention, le savais-tu, exercice résolu,
% exercice, TP, bilan), page de garde et signature OnBuch+ ancrée en bas.
%
% Macros attendues dans main.tex AVANT \input{preamble} :
%   \DOCMATIERE  \DOCNIVEAU  \DOCMODULE  \DOCLECON (numéro)  \DOCTITRE
\documentclass[11pt]{article}

\usepackage{fontspec}
\usepackage[english,french]{babel} % français = langue principale ; l'anglais via \eng{..} / textebox
\usepackage[a4paper,margin=2cm,top=3.4cm,bottom=2.8cm,headheight=1.4cm,footskip=1.3cm]{geometry}
\usepackage{amsmath,amssymb,amsfonts}
\usepackage{mathtools} % \xmapsto, \coloneqq... souvent utilisés par les modèles
\usepackage{cancel} % simplifications barrées (bilans de forces)
% \abs{z} (module, valeur absolue) et \norm{u} (norme), utilisés par les modèles
\providecommand{\abs}{}\let\abs\relax\DeclarePairedDelimiter{\abs}{\lvert}{\rvert}
\providecommand{\norm}{}\let\norm\relax\DeclarePairedDelimiter{\norm}{\lVert}{\rVert}
\usepackage[table]{xcolor} % [table] : fournit aussi \rowcolors (alternance de lignes)
\usepackage{pagecolor}
\usepackage{tikz}
\usetikzlibrary{arrows.meta,positioning,calc,shapes.geometric,shapes.misc,shapes.arrows,decorations.pathmorphing,decorations.pathreplacing,decorations.markings,patterns,shadows,mindmap,trees,fit,backgrounds,matrix,intersections,angles,quotes}
\usepackage{pgfplots}
\usepackage[version=4]{mhchem} % équations nucléaires \ce{^{235}_{92}U}
\usepackage[european,siunitx]{circuitikz} % schémas électriques
\pgfplotsset{compat=1.18}
\usepgfplotslibrary{fillbetween,groupplots} % aires entre courbes (calcul intégral)
\usepackage{enumitem}
\usepackage{fancyhdr}
% [most] charge toutes les bibliothèques tcolorbox (listings, minted, raster,
% documentation...) et gonfle inutilement la mémoire TeX sur un document long
% (~300 boîtes) : on ne charge que ce qui est réellement utilisé.
\usepackage[skins,breakable]{tcolorbox}
\usepackage{graphicx}
\usepackage{colortbl}
\usepackage{booktabs}
\usepackage{tabularx}
\usepackage{diagbox} % case d'en-tête diagonale des tableaux à double entrée
% Colonnes extensibles centrée / gauche / droite (C, L, R), souvent utilisées par les modèles
\newcolumntype{C}{>{\centering\arraybackslash}X}
\newcolumntype{L}{>{\raggedright\arraybackslash}X}
\newcolumntype{R}{>{\raggedleft\arraybackslash}X}
\usepackage{array}
\usepackage{multicol}
\usepackage{multirow}
\usepackage{siunitx}
% parse-numbers=false : accepte aussi \SI{2,5 \times 10^{-3}}{...}
\sisetup{locale=FR,output-decimal-marker={,},group-minimum-digits=4,parse-numbers=false}
\DeclareSIUnit\ohm{\ensuremath{\symup{\Omega}}} % U+2126 absent de la police
\DeclareSIUnit\division{div} % réglages d'oscilloscope (ms/div, V/div)
\DeclareSIUnit\tr{tr} % tours (vitesse de rotation, ex. \SI{1500}{\tr\per\minute})
\DeclareSIUnit\molar{M} % concentration molaire (ex. \SI{3}{\molar}, \SI{1}{\micro\molar})
\DeclareSIUnit\an{an} % année (le modèle confond parfois avec \year, un compteur TeX) — ex. \SI{5}{\kilo\meter\per\an}
\DeclareSIUnit\FCFA{FCFA} % franc CFA (ex. \SI{500}{\FCFA\per\kilo\gram})
\DeclareSIUnit\degreeBrix{\textdegree Brix} % teneur en sucre d'un sirop/jus (réfractométrie)
\DeclareSIUnit\month{mois} % le modèle utilise parfois \month, un compteur TeX, comme unité de durée
\usepackage{qrcode}
% \degree : commande fréquemment utilisée par les modèles pour un angle ou une
% température en dehors de \SI{}{} (non fournie par les paquets chargés).
\providecommand{\degree}{^{\circ}}
% \qed (fin de démonstration), utilisé par les modèles sans amsthm
\providecommand{\qed}{\hfill$\square$}
% \barre : double barre des valeurs interdites dans un tableau de variations (tkz-tab)
\providecommand{\barre}{\ensuremath{\|}}
% environnement proof (amsthm non chargé) utilisé par les modèles
\ifdefined\proof\else\newenvironment{proof}[1][Démonstration]{\par\noindent\textit{#1.}\ }{\qed\par}\fi
\usepackage{lastpage}
\usepackage{hyperref}
\hypersetup{hidelinks}

% ============================================================
% Palette « Pop » OnBuch+ — mêmes hex que la classe P (plus_ui.dart)
% ============================================================
\definecolor{poporange}{HTML}{F59321}
\definecolor{popdark}{HTML}{E07A0C}
\definecolor{popcream}{HTML}{FFF6E8}
\definecolor{popink}{HTML}{1C1714}
\definecolor{poppurple}{HTML}{7A5AE0}
\definecolor{popgreen}{HTML}{2E9E6A}
\definecolor{poppink}{HTML}{E0457B}
\definecolor{popblue}{HTML}{2F7DE1}
\definecolor{popgold}{HTML}{C98A12}
\definecolor{popmuted}{HTML}{8A7F76}
\definecolor{popline}{HTML}{EADFD2}
\definecolor{popgray}{HTML}{8A7F76}
\colorlet{popgrey}{popgray}
% Alias tolérants (le modèle nomme parfois les couleurs différemment)
\colorlet{popwhite}{white}
\colorlet{popblack}{popink}
\colorlet{popred}{poppink}
\colorlet{popyellow}{popgold}
% Teintes claires (fonds de tableaux/encadrés) — le modèle nomme parfois
% les couleurs avec un suffixe L (ex. popblueL) sans les avoir définies.
\colorlet{poporangeL}{poporange!20}
\colorlet{popdarkL}{popdark!20}
\colorlet{poppurpleL}{poppurple!20}
\colorlet{popgreenL}{popgreen!20}
\colorlet{poppinkL}{poppink!20}
\colorlet{popblueL}{popblue!20}
\colorlet{popgoldL}{popgold!20}
\colorlet{popredL}{popred!20}
\colorlet{popgrayL}{popgray!20}
% Teintes claires pour fonds de tableaux / zones
\colorlet{popblueL}{popblue!10!white}
\colorlet{poporangeL}{poporange!14!white}
\colorlet{popgreenL}{popgreen!12!white}
\colorlet{poppurpleL}{poppurple!12!white}
\colorlet{poppinkL}{poppink!12!white}
\colorlet{popgoldL}{popgold!14!white}

\pagecolor{popcream}

% ============================================================
% Typographie
% ============================================================
\defaultfontfeatures{Ligatures=TeX}
\setmainfont{PlusJakartaSans-Medium.ttf}[Path=../../../fonts/,BoldFont=PlusJakartaSans-Bold.ttf]
% Maths en sans-serif (Fira Math) avec chiffres et lettres droites de Plus
% Jakarta, pour que les formules se fondent dans le texte.
\usepackage[math-style=ISO,bold-style=ISO]{unicode-math}
\setmathfont{FiraMath-Regular.otf}[Path=../../../fonts/]
\setmathfont{PlusJakartaSans-Medium.ttf}[Path=../../../fonts/,range={up/num,up/{Latin,latin},\mathup}]
\usepackage{icomma} % virgule décimale sans espace en mode math
% Caractères mathématiques Unicode absents des polices de texte (Plus Jakarta,
% Archivo Black) : rendus par leur équivalent mathématique, en texte comme en
% formule — sinon ils s'affichent en carré vide (ex. « Arithmétique dans ℤ »).
\usepackage{newunicodechar}
\newunicodechar{ℝ}{\ensuremath{\mathbb{R}}}
\newunicodechar{ℤ}{\ensuremath{\mathbb{Z}}}
\newunicodechar{ℂ}{\ensuremath{\mathbb{C}}}
\newunicodechar{ℕ}{\ensuremath{\mathbb{N}}}
\newunicodechar{ℚ}{\ensuremath{\mathbb{Q}}}
\newunicodechar{𝔻}{\ensuremath{\mathbb{D}}}
\newunicodechar{α}{\ensuremath{\alpha}}
\newunicodechar{β}{\ensuremath{\beta}}
\newunicodechar{γ}{\ensuremath{\gamma}}
\newunicodechar{δ}{\ensuremath{\delta}}
\newunicodechar{ε}{\ensuremath{\varepsilon}}
\newunicodechar{θ}{\ensuremath{\theta}}
\newunicodechar{λ}{\ensuremath{\lambda}}
\newunicodechar{μ}{\ensuremath{\mu}}
\newunicodechar{σ}{\ensuremath{\sigma}}
\newunicodechar{τ}{\ensuremath{\tau}}
\newunicodechar{φ}{\ensuremath{\varphi}}
\newunicodechar{ψ}{\ensuremath{\psi}}
\newunicodechar{ω}{\ensuremath{\omega}}
\newunicodechar{Σ}{\ensuremath{\Sigma}}
% majuscules produites par \MakeUppercase dans les titres (α-aminés -> Α)
\newunicodechar{Α}{\ensuremath{\alpha}}
\newunicodechar{Β}{\ensuremath{\beta}}
\newunicodechar{Γ}{\ensuremath{\Gamma}}
\newunicodechar{Δ}{\ensuremath{\Delta}}
\newunicodechar{Ε}{\ensuremath{\varepsilon}}
\newunicodechar{Θ}{\ensuremath{\Theta}}
\newunicodechar{Λ}{\ensuremath{\Lambda}}
\newunicodechar{Μ}{\ensuremath{\mu}}
\newunicodechar{Τ}{\ensuremath{\tau}}
\newunicodechar{Φ}{\ensuremath{\Phi}}
\newunicodechar{Ψ}{\ensuremath{\Psi}}
\newunicodechar{Ω}{\ensuremath{\Omega}}
\newunicodechar{↦}{\ensuremath{\mapsto}}
\newunicodechar{∈}{\ensuremath{\in}}
\newunicodechar{∉}{\ensuremath{\notin}}
\newunicodechar{∧}{\ensuremath{\wedge}}
\newunicodechar{⇔}{\ensuremath{\Leftrightarrow}}
\newunicodechar{⟺}{\ensuremath{\Longleftrightarrow}}
\newunicodechar{⇒}{\ensuremath{\Rightarrow}}
\newunicodechar{⟹}{\ensuremath{\Longrightarrow}}
\newunicodechar{∘}{\ensuremath{\circ}}
\newunicodechar{∪}{\ensuremath{\cup}}
\newunicodechar{∩}{\ensuremath{\cap}}
\newunicodechar{≡}{\ensuremath{\equiv}}
\newunicodechar{⊂}{\ensuremath{\subset}}
\newunicodechar{⊥}{\ensuremath{\perp}}
\newunicodechar{∃}{\ensuremath{\exists}}
\newunicodechar{∀}{\ensuremath{\forall}}
\newunicodechar{∛}{\ensuremath{\sqrt[3]{\,}}}
\newunicodechar{∜}{\ensuremath{\sqrt[4]{\,}}}
\newunicodechar{⃗}{\ensuremath{{}^{\to}}}
\newunicodechar{ⁿ}{\ifmmode^{n}\else\textsuperscript{\ensuremath{n}}\fi}
\newunicodechar{ˣ}{\ifmmode^{x}\else\textsuperscript{\ensuremath{x}}\fi}
\newunicodechar{ᵇ}{\ifmmode^{b}\else\textsuperscript{\ensuremath{b}}\fi}
\newunicodechar{ʳ}{\ifmmode^{r}\else\textsuperscript{\ensuremath{r}}\fi}
\newunicodechar{ᵘ}{\ifmmode^{u}\else\textsuperscript{\ensuremath{u}}\fi}
\newunicodechar{ʸ}{\ifmmode^{y}\else\textsuperscript{\ensuremath{y}}\fi}
\newunicodechar{ᵃ}{\ifmmode^{a}\else\textsuperscript{\ensuremath{a}}\fi}
\newunicodechar{ᵉ}{\ifmmode^{e}\else\textsuperscript{\ensuremath{e}}\fi}
\newunicodechar{ᵏ}{\ifmmode^{k}\else\textsuperscript{\ensuremath{k}}\fi}
\newunicodechar{ᵖ}{\ifmmode^{p}\else\textsuperscript{\ensuremath{p}}\fi}
\newunicodechar{ᴺ}{\ifmmode^{N}\else\textsuperscript{\ensuremath{N}}\fi}
\newunicodechar{ᶜ}{\ifmmode^{c}\else\textsuperscript{\ensuremath{c}}\fi}
\newunicodechar{ʰ}{\ifmmode^{h}\else\textsuperscript{\ensuremath{h}}\fi}
\newunicodechar{ᵠ}{\ifmmode^{q}\else\textsuperscript{\ensuremath{q}}\fi}
\newunicodechar{ₙ}{\ifmmode_{n}\else\textsubscript{\ensuremath{n}}\fi}
\newunicodechar{ₐ}{\ifmmode_{a}\else\textsubscript{\ensuremath{a}}\fi}
\newunicodechar{ᵢ}{\ifmmode_{i}\else\textsubscript{\ensuremath{i}}\fi}
\newunicodechar{ᵣ}{\ifmmode_{r}\else\textsubscript{\ensuremath{r}}\fi}
\newunicodechar{ₖ}{\ifmmode_{k}\else\textsubscript{\ensuremath{k}}\fi}
\newunicodechar{ᵦ}{\ifmmode_{\beta}\else\textsubscript{\ensuremath{\beta}}\fi}
\newunicodechar{ₓ}{\ifmmode_{x}\else\textsubscript{\ensuremath{x}}\fi}
\newfontfamily\popdisplay{ArchivoBlack-Regular.ttf}[Path=../../../fonts/]
\newfontfamily\poplogo{SpaceGrotesk-Bold.ttf}[Path=../../../fonts/]
\newfontfamily\popsemi{PlusJakartaSans-SemiBold.ttf}[Path=../../../fonts/]
\newfontfamily\popxbold{PlusJakartaSans-ExtraBold.ttf}[Path=../../../fonts/]
\newfontfamily\popxboldls{PlusJakartaSans-ExtraBold.ttf}[Path=../../../fonts/,LetterSpace=3.0]

\setlist[enumerate]{leftmargin=1.7em,itemsep=5pt,topsep=4pt}
\setlist[itemize]{leftmargin=1.7em,itemsep=4pt,topsep=4pt,label={\color{poporange}\textbullet}}
\setlength{\parindent}{0pt}
\setlength{\parskip}{5pt}
\renewcommand{\arraystretch}{1.25}

% pgfplots : style Pop par défaut
\pgfplotsset{
  every axis/.append style={axis line style={popink,line width=1pt},tick style={popink},
    grid style={popline,line width=0.5pt},label style={font=\small},tick label style={font=\footnotesize}},
  popcurve/.style={poporange,line width=1.8pt,no markers,smooth},
}
% Même style disponible dans un \draw TikZ simple (hors pgfplots)
\tikzset{popcurve/.style={poporange,line width=1.8pt}}
\tikzset{popgrid/.style={popline,very thin}}
\colorlet{popcurve}{poporange} % le nom du style est parfois utilisé comme couleur

% ============================================================
% Pill / squiggle
% ============================================================
\newcommand{\poppill}[3]{%
  \tikz[baseline=-0.55ex]\node[draw=popink,line width=1.2pt,rounded corners=2.6pt,%
    fill=#2,inner xsep=6.5pt,inner ysep=3pt]{%
    \popxboldls\fontsize{7}{7}\selectfont\color{#3}\MakeUppercase{#1}};%
}
\newcommand{\popsquiggle}[1]{%
  \tikz[baseline=-0.4ex]{
    \draw[#1,line width=2.6pt,line cap=round]
      plot [smooth] coordinates {(0,0.09) (0.34,0.01) (0.68,0.21) (1.02,0.01) (1.36,0.10)};
  }%
}
% Mot-clé surligné (marqueur orange sous le texte)
\newcommand{\cle}[1]{{\popxbold\color{popdark}#1}}

% ============================================================
% En-tête / pied
% ============================================================
\pagestyle{fancy}
\fancyhf{}
\renewcommand{\headrulewidth}{0pt}
\renewcommand{\footrulewidth}{0pt}
\lhead{\raisebox{-0.15ex}{\poplogo\fontsize{13}{13}\selectfont\color{popink}OnBuch\color{poporange}+}}
\rhead{\poppill{\DOCMATIERE\ \textbullet\ \DOCNIVEAU\ \textbullet\ Leçon \DOCLECON}{white}{popink}}
\lfoot{\popxbold\fontsize{7.6}{7.6}\selectfont\color{popmuted}\MakeUppercase{Cours signé OnBuch+}\ \textbullet\ {\popsemi\DOCTITRE}}
\rfoot{\tikz[baseline=-0.6ex]\node[rounded corners=8.5pt,fill=popink,minimum height=17pt,minimum width=17pt,inner xsep=5pt,inner sep=0]{\popxbold\fontsize{8}{8}\selectfont\color{popcream}\thepage\,/\,\pageref*{LastPage}};}

% ============================================================
% Titres
% ============================================================
\newcommand{\chaptitre}[2]{%
  \begin{tcolorbox}[enhanced,colback=poporange,colframe=popink,boxrule=2.6pt,arc=11pt,
      drop shadow=popink,left=15pt,right=15pt,top=12pt,bottom=11pt,before skip=3pt,after skip=18pt]
    \poppill{#2}{popink}{popcream}\par\vspace{9pt}
    {\raggedright\hyphenpenalty=10000\exhyphenpenalty=10000\popdisplay\fontsize{22}{24}\selectfont\color{popcream}\MakeUppercase{#1}\par}
  \end{tcolorbox}
}
% Section numérotée : pastille orange avec le numéro + titre Archivo
\newcounter{popsec}
\newcommand{\coursec}[1]{%
  \stepcounter{popsec}\setcounter{popsub}{0}%
  \par\vspace{16pt}\noindent
  \begin{minipage}{\linewidth}
  \tikz[baseline=-0.7ex]\node[draw=popink,line width=1.6pt,fill=poporange,rounded corners=4pt,minimum size=22pt,inner sep=0]{\popdisplay\fontsize{12}{12}\selectfont\color{popcream}\thepopsec};%
  \hspace{8pt}{\popdisplay\fontsize{15}{17}\selectfont\color{popink}\MakeUppercase{#1}}\par
  \vspace{4pt}\noindent{\color{poporange}\rule{36pt}{3.6pt}}
  \end{minipage}\par\nopagebreak\vspace{8pt}
}
\newcounter{popsub}
\newcommand{\courssub}[1]{%
  \stepcounter{popsub}%
  \par\vspace{9pt}\noindent{\popxbold\fontsize{11.5}{13}\selectfont\color{popink}{\color{poporange}\thepopsec.\thepopsub}\hspace{6pt}#1}\par\nopagebreak\vspace{4pt}
}

% ============================================================
% Boîtes pédagogiques — même recette : fond blanc, bordure épaisse colorée,
% ombre dure encre, badge plein. Titre optionnel [#1].
% ============================================================
\tcbset{popbase/.style={breakable,enhanced,colback=white,boxrule=2.3pt,arc=9pt,
  shadow={3pt}{-3pt}{0pt}{popink},left=14pt,right=14pt,top=11pt,bottom=11pt,before skip=11pt,after skip=15pt,
  fontupper=\popsemi\fontsize{10.6}{14.5}\selectfont\color{popink},
  fontlower=\popsemi\fontsize{10.6}{14.5}\selectfont\color{popink}}}
% Coche dessinée (le glyphe ✓ n'existe pas dans les polices du document)
\newcommand{\popcheck}[1]{\tikz[baseline=-0.5ex]\draw[#1,line width=1.6pt,line cap=round,line join=round](-0.13,0.0)--(-0.04,-0.09)--(0.13,0.1);}
\providecommand{\coche}{\popcheck{popgreen}}
\providecommand{\resultat}[1]{\par\smallskip\noindent\fcolorbox{popgreen}{popgreenL}{\parbox{\dimexpr\linewidth-2\fboxsep-2\fboxrule\relax}{\textbf{Résultat.} #1}}\par\smallskip}
\newcommand{\popboxhead}[3]{\poppill{#1}{#2}{white}\ifx\relax#3\relax\else\hspace{6pt}{\popxbold\fontsize{10.6}{12}\selectfont\color{popink}#3}\fi\par\vspace{7pt}}

\newtcolorbox{aretenir}[1][]{popbase,colframe=popgreen,before upper={\popboxhead{À retenir}{popgreen}{#1}}}
% Texte en anglais (document de lecture, dialogue, extrait) : cadre
% d'étude. Un titre optionnel (source, auteur) s'affiche dans l'en-tête.
\newtcolorbox{textebox}[1][]{popbase,colframe=popblue,before upper={\popboxhead{Text}{popblue}{#1}\begin{otherlanguage}{english}},after upper={\end{otherlanguage}}}
% Marques ✓ / ✗ (forme correcte / fautive) souvent utilisées par les modèles
\providecommand{\cross}{\textcolor{poppink}{\ensuremath{\times}}}
\providecommand{\xmark}{\cross}\providecommand{\crossmark}{\cross}\providecommand{\cmark}{\ensuremath{\checkmark}}
% Ligne de réponse / justification à compléter (inventée par les modèles)
\providecommand{\justification}[1]{\par\noindent\textit{Justification~:} \dotfill\par}
% \textipa (package tipa non chargé) : la police ne couvre pas l'API -> simple passage
\providecommand{\textipa}[1]{#1}
% Soulignement (ulem non chargé) : \uline, \ul
\providecommand{\uline}[1]{\underline{#1}}\providecommand{\ul}[1]{\underline{#1}}
% \glossaire{\item ...} inventée par les modèles : liste de glossaire
\providecommand{\glossaire}[1]{\begin{itemize}#1\end{itemize}}
% \alert (beamer) inventée par les modèles : texte mis en évidence
\providecommand{\alert}[1]{\textbf{\textcolor{poppink}{#1}}}
% variantes ASCII d'ordinaux écrites en macro
\providecommand{\iere}{\textsuperscript{re}}\providecommand{\ieme}{\textsuperscript{e}}\providecommand{\ere}{\textsuperscript{re}}\providecommand{\eme}{\textsuperscript{e}}\providecommand{\er}{\textsuperscript{er}}\providecommand{\ie}{\textsuperscript{e}}
% Ordinaux français écrits en macro par les modèles : 1\ière, 3\ième
\newcommand{\ière}{\textsuperscript{re}}\newcommand{\ième}{\textsuperscript{e}}
% \exemple (inventée par les modèles) : étiquette « Exemple : »
\providecommand{\exemple}{\textbf{Exemple~:}\ }
% \square utilisable aussi hors mode math (cases à cocher)
\AtBeginDocument{\let\squareMath\square\renewcommand{\square}{\ensuremath{\squareMath}}}
% Mot ou phrase en anglais à l'intérieur d'un paragraphe français
\newcommand{\eng}[1]{\ifmmode\text{\textit{#1}}\else\begingroup\selectlanguage{english}\itshape #1\endgroup\fi}
\let\esp\eng % alias toléré
\newtcolorbox{exemplebox}[1][]{popbase,colframe=poporange,before upper={\popboxhead{Exemple}{poporange}{#1}}}
\newtcolorbox{definition}[1][]{popbase,colframe=popblue,before upper={\popboxhead{Définition}{popblue}{#1}}}
\newtcolorbox{propriete}[1][]{popbase,colframe=poppurple,before upper={\popboxhead{Propriété}{poppurple}{#1}}}
\newtcolorbox{methode}[1][]{popbase,colframe=popgold,before upper={\popboxhead{Méthode}{popgold}{#1}}}
\newtcolorbox{attention}[1][]{popbase,colframe=poppink,before upper={\popboxhead{Attention}{poppink}{#1}}}
\newtcolorbox{experience}[1][]{popbase,colframe=popblue,colback=popblueL,before upper={\popboxhead{Expérience}{popblue}{#1}}}
% « Le savais-tu ? » : carte violette pleine (contexte, culture, Cameroun)
\newtcolorbox{savaistu}[1][]{popbase,colback=poppurple,colframe=popink,
  fontupper=\popsemi\fontsize{10.4}{14.2}\selectfont\color{white},
  before upper={\poppill{Le savais-tu\,?}{white}{poppurple}\ifx\relax#1\relax\else\hspace{6pt}{\popxbold\color{white}#1}\fi\par\vspace{7pt}}}
% Exercice résolu : énoncé en haut, \tcblower puis la solution sur fond crème
\newcounter{popexo}\newcounter{popexr}
\newtcolorbox{exoresolu}[1][]{popbase,colframe=poporange,colbacklower=popcream,
  segmentation style={popink,line width=1.2pt,dashed},
  before upper={\stepcounter{popexr}\popboxhead{Exercice résolu \thepopexr}{poporange}{#1}},
  before lower={\poppill{Solution}{popink}{popcream}\par\vspace{6pt}}}
% Exercice (non résolu en place) — la correction est dans la partie Corrigés
\newtcolorbox{exercice}[1][]{popbase,colframe=popink,
  before upper={\stepcounter{popexo}\popboxhead{Exercice \thepopexo}{popink}{#1}}}
% Corrigé
\newtcolorbox{corrige}[1][]{popbase,colframe=popgreen,colback=popgreenL,
  before upper={\popboxhead{Corrigé}{popgreen}{#1}}}
% Objectifs / prérequis / bilan
\newtcolorbox{objectifs}{popbase,colframe=popink,colback=poporangeL,
  before upper={\popboxhead{Objectifs de la leçon}{popink}{}}}
\newtcolorbox{prerequis}{popbase,colframe=popink,colback=popblueL,
  before upper={\popboxhead{Prérequis}{popblue}{}}}
\newtcolorbox{bilan}{popbase,colframe=popink,colback=white,boxrule=2.8pt,
  before upper={\popboxhead{Fiche bilan}{poporange}{L'essentiel à connaître pour l'examen}}}
% Figure encadrée avec légende
\newtcolorbox{popfigure}[1][]{enhanced,breakable=false,colback=white,colframe=popline,boxrule=1.4pt,arc=8pt,
  left=10pt,right=10pt,top=10pt,bottom=8pt,before skip=10pt,after skip=12pt,halign=center,
  lower separated=false,halign lower=center,
  fontlower=\popsemi\fontsize{9}{11.5}\selectfont\color{popmuted},
  #1}
\newcommand{\legende}[1]{\par\vspace{6pt}{\popsemi\fontsize{9}{11.5}\selectfont\color{popmuted}#1}}

% ============================================================
% Signature OnBuch+
% ============================================================
\newcommand{\poplogoinline}[1]{{\poplogo\fontsize{#1}{#1}\selectfont\color{popink}OnBuch\color{poporange}+}}
\newcommand{\signatureonbuch}{%
  \begin{tcolorbox}[enhanced,colback=popink,colframe=popink,boxrule=2.2pt,arc=10pt,
      drop shadow=poporange,left=14pt,right=14pt,top=10pt,bottom=10pt,before skip=0pt,after skip=0pt]
    \tikz[baseline=-0.7ex]\node[circle,fill=popgreen,draw=popcream,line width=1.3pt,minimum size=22pt,inner sep=0]{\popcheck{white}};%
    \hspace{9pt}\begin{minipage}[c]{0.62\linewidth}
      {\popxbold\fontsize{12}{13}\selectfont\color{popcream}Signé\ \textbullet\ {\poplogo OnBuch\color{poporange}+}}\par\vspace{2pt}
      {\raggedright\popsemi\fontsize{8}{10.5}\selectfont\color{popline}\MakeUppercase{Équipe pédagogique OnBuch+}\\\MakeUppercase{Conforme au programme officiel MINESEC}\par}
    \end{minipage}\hfill
    {\poplogo\fontsize{22}{22}\selectfont\color{popcream}OnBuch\color{poporange}+}
  \end{tcolorbox}%
}

% ============================================================
% Page de garde
% #1 titre · #2 matière · #3 niveau · #4 module · #5 n° leçon · #6 durée ·
% #7 description / objectifs (texte ou itemize)
% ============================================================
\newcommand{\pagegarde}[7]{%
  \thispagestyle{empty}
  \newgeometry{a4paper,margin=1.7cm,top=1.6cm,bottom=1.4cm}%
  \begin{tikzpicture}[remember picture,overlay]
    \fill[poporange!18] ([xshift=-3.2cm,yshift=-3.6cm]current page.north east) circle (4.6cm);
    \fill[poppurple!14] ([xshift=2.4cm,yshift=7.6cm]current page.south west) circle (3.8cm);
  \end{tikzpicture}%
  {\poplogo\fontsize{30}{30}\selectfont\color{popink}OnBuch\color{poporange}+}\hfill
  \raisebox{0.6ex}{\poppill{Cours signé \textbullet\ Premium}{popink}{popcream}}\par
  \vspace{1pt}\popsquiggle{poppurple}\par
  \vspace{8pt}
  \begin{tcolorbox}[enhanced,colback=poporange,colframe=popink,boxrule=2.8pt,arc=13pt,
      drop shadow=popink,left=18pt,right=18pt,top=12pt,bottom=13pt,after skip=10pt]
    \poppill{#2 \textbullet\ #3}{popink}{popcream}\hspace{5pt}\poppill{#4}{white}{popink}\par\vspace{8pt}
    {\popdisplay\fontsize{14}{14}\selectfont\color{popink}LEÇON #5}\par\vspace{4pt}
    {\raggedright\hyphenpenalty=10000\exhyphenpenalty=10000\popdisplay\fontsize{25}{27}\selectfont\color{popcream}\MakeUppercase{#1}\par}
    \vspace{8pt}
    {\popxbold\fontsize{9.5}{9.5}\selectfont\color{popink}\MakeUppercase{Durée conseillée : #6}}
  \end{tcolorbox}
  \begin{tcolorbox}[enhanced,colback=white,colframe=popink,boxrule=2.2pt,arc=10pt,
      drop shadow=popink,left=16pt,right=16pt,top=10pt,bottom=10pt,after skip=8pt]
    \poppill{Dans cette leçon}{poppurple}{white}\par\vspace{5pt}
    \popsemi\fontsize{9.3}{12.6}\selectfont\color{popink}#7
  \end{tcolorbox}
  \noindent\begin{minipage}[t]{0.487\linewidth}
    \begin{tcolorbox}[enhanced,colback=popgreen,colframe=popink,boxrule=2.2pt,arc=10pt,
        drop shadow=popink,left=11pt,right=11pt,top=10pt,bottom=10pt,height=3.1cm,valign=center]
      \tcbox[colback=white,colframe=popink,boxrule=1.3pt,arc=5pt,boxsep=0pt,nobeforeafter,
        left=3pt,right=3pt,top=3pt,bottom=3pt,valign=center]{\qrcode[height=1.9cm]{https://play.google.com/store/apps/details?id=com.luvvix.onbuch&hl=fr}}%
      \hspace{8pt}\begin{minipage}[c]{0.5\linewidth}
        {\popdisplay\fontsize{11}{12.5}\selectfont\color{white}\MakeUppercase{Télécharge OnBuch}}\par\vspace{4pt}
        {\raggedright\popsemi\fontsize{8.4}{11}\selectfont\color{white}Gratuit \textbullet\ Google Play\\Tuteur IA, annales, concours, résultats\par}
      \end{minipage}
    \end{tcolorbox}
  \end{minipage}\hfill
  \begin{minipage}[t]{0.487\linewidth}
    \begin{tcolorbox}[enhanced,colback=poppurple,colframe=popink,boxrule=2.2pt,arc=10pt,
        drop shadow=popink,left=11pt,right=11pt,top=10pt,bottom=10pt,height=3.1cm,valign=center]
      \tikz[baseline=-0.7ex]\node[circle,fill=white,draw=popink,line width=1.3pt,minimum size=1.6cm,inner sep=0]{\popdisplay\fontsize{24}{24}\selectfont\color{poppurple}+};%
      \hspace{8pt}\begin{minipage}[c]{0.6\linewidth}
        {\popdisplay\fontsize{11}{12.5}\selectfont\color{white}\MakeUppercase{Passe à OnBuch+}}\par\vspace{4pt}
        {\raggedright\popsemi\fontsize{8.4}{11}\selectfont\color{white}Tous les cours signés, Tuteur IA illimité, fiches PDF — onglet OnBuch+ de l'app\par}
      \end{minipage}
    \end{tcolorbox}
  \end{minipage}
  \vspace{8pt}
  \popcontenu
  \vfill
  \signatureonbuch
  \restoregeometry
  \newpage
}

% Rangée « Ce cours contient » (page de garde)
\newcommand{\poptile}[3]{% #1 couleur #2 gros texte #3 légende
  \begin{tcolorbox}[enhanced,colback=#1,colframe=popink,boxrule=2pt,arc=9pt,shadow={2.5pt}{-2.5pt}{0pt}{popink},
      left=6pt,right=6pt,top=9pt,bottom=9pt,halign=center,valign=center,height=1.75cm,width=\linewidth,nobeforeafter]
    {\popdisplay\fontsize{12.5}{13}\selectfont\color{white}\MakeUppercase{#2}}\par\vspace{3pt}
    {\centering\popsemi\fontsize{7.8}{9.5}\selectfont\color{white}#3\par}
  \end{tcolorbox}}
\newcommand{\popcontenu}{%
  \par\noindent{\popxboldls\fontsize{7.5}{7.5}\selectfont\color{popmuted}CE COURS SIGNÉ CONTIENT}\par\vspace{7pt}
  \noindent\begin{minipage}[t]{0.235\linewidth}\poptile{popblue}{Cours}{complet, illustré, pas à pas}\end{minipage}\hfill
  \begin{minipage}[t]{0.235\linewidth}\poptile{poporange}{Méthodes}{et exercices résolus}\end{minipage}\hfill
  \begin{minipage}[t]{0.235\linewidth}\poptile{poppink}{Exercices}{type examen + corrigés}\end{minipage}\hfill
  \begin{minipage}[t]{0.235\linewidth}\poptile{popgreen}{Fiche bilan}{l'essentiel à retenir}\end{minipage}\par}

% Sommaire
\newtcolorbox{sommaire}{enhanced,colback=white,colframe=popink,boxrule=2.2pt,arc=10pt,
  drop shadow=popink,left=18pt,right=18pt,top=13pt,bottom=13pt,before skip=6pt,after skip=20pt,
  before upper={\poppill{Sommaire}{poppurple}{white}\par\vspace{9pt}\popsemi\fontsize{10.6}{14.5}\selectfont\color{popink}}}

% Signature de fin de cours — poussée tout en bas de la dernière page
\newcommand{\findecours}{%
  \par\vfill\nopagebreak
  \begin{center}\popsquiggle{poporange}\end{center}\vspace{4pt}
  \signatureonbuch
}
\AtBeginDocument{\def\llbracket{\mathopen{[\mkern-3.5mu[}}\def\rrbracket{\mathclose{]\mkern-3.5mu]}}} % intervalles d'entiers (glyphe ⟦ absent de la police)
% Glyphes absents de Fira Math : redéfinis à partir de glyphes disponibles
\AtBeginDocument{%
  \def\setminus{\mathbin{\backslash}}\let\smallsetminus\setminus
  \def\vdots{\mathord{\vbox{\baselineskip3.2pt\lineskiplimit0pt\kern4pt\hbox{.}\hbox{.}\hbox{.}}}}%
  \def\ddots{\mathinner{\mkern1mu\raise6.5pt\hbox{.}\mkern2mu\raise3.6pt\hbox{.}\mkern2mu\raise0.7pt\hbox{.}\mkern1mu}}%
  \def\triangle{\mathord{\scalebox{0.9}{$\symup{\Delta}$}}}%
  \def\nabla{\mathord{\raisebox{\depth}{\scalebox{1}[-1]{$\symup{\Delta}$}}}}%
  \def\checkmark{\popcheck{popdark}}%
  \def\star{\mathbin{\ast}}\def\bigstar{\mathord{\ast}}%
  \def\diamond{\mathbin{\scalebox{0.6}{$\Diamond$}}}%
  \def\bigcup{\mathop{\mathchoice{\raisebox{-0.3em}{\scalebox{1.7}{$\cup$}}}{\scalebox{1.25}{$\cup$}}{\cup}{\cup}}\displaylimits}%
  \def\bigcap{\mathop{\mathchoice{\raisebox{-0.3em}{\scalebox{1.7}{$\cap$}}}{\scalebox{1.25}{$\cap$}}{\cap}{\cap}}\displaylimits}%
  \def\lesseqqgtr{\mathrel{\lessgtr}}\def\bigoplus{\mathop{\oplus}\displaylimits}%
}
\providecommand{\ssi}{si et seulement si} % abréviation fréquente
\AtBeginDocument{\providecommand{\wideparen}{\overparen}} % arc de cercle

\begin{document}

\pagegarde{Lesson 31 — Reading: Texts about environmental issues like climate change and global warming}{Anglais}{Tle A}{Chapter 3 : Environment, Well-being and Health}{EN31}{2 h}{Cette leçon de lecture propose un article original sur le changement climatique et le réchauffement planétaire. Elle entraîne les élèves aux stratégies de skimming et scanning, enrichit leur lexique de l'environnement et les prépare aux questions de compréhension et d'interprétation du Bac.
\vspace{4pt}
\begin{multicols}{2}\small
\begin{itemize}
\item À la fin de cette leçon, je sais utiliser le skimming pour dégager l'idée générale d'un texte sur l'environnement
\item À la fin de cette leçon, je sais utiliser le scanning pour repérer rapidement des informations précises (chiffres, dates, noms)
\item À la fin de cette leçon, je sais déduire le sens d'un mot inconnu à partir du contexte
\item À la fin de cette leçon, je sais répondre à des questions de compréhension de types variés (vrai/faux, QCM, réponses courtes)
\item À la fin de cette leçon, je sais employer le lexique essentiel du changement climatique en anglais
\item À la fin de cette leçon, je sais résumer un texte environnemental en quelques phrases
\end{itemize}\end{multicols}}

\begin{sommaire}
\begin{multicols}{2}
\begin{enumerate}
\item Warm-up et découverte du thème
\item Texte support : article de lecture
\item Compréhension du texte
\item Lexique thématique : l'environnement et le climat
\item Méthode de lecture : skimming, scanning, déduction
\item Production guidée : résumé et opinion
\item Activité d'intégration
\item Méthodes et exercices résolus
\item Exercices
\item Corrigés des exercices
\item Fiche bilan
\end{enumerate}
\end{multicols}
\end{sommaire}

\begin{objectifs}
\begin{itemize}
\item À la fin de cette leçon, je sais utiliser le skimming pour dégager l'idée générale d'un texte sur l'environnement
\item À la fin de cette leçon, je sais utiliser le scanning pour repérer rapidement des informations précises (chiffres, dates, noms)
\item À la fin de cette leçon, je sais déduire le sens d'un mot inconnu à partir du contexte
\item À la fin de cette leçon, je sais répondre à des questions de compréhension de types variés (vrai/faux, QCM, réponses courtes)
\item À la fin de cette leçon, je sais employer le lexique essentiel du changement climatique en anglais
\item À la fin de cette leçon, je sais résumer un texte environnemental en quelques phrases
\item À la fin de cette leçon, je sais exprimer mon opinion sur un problème environnemental à l'écrit
\item À la fin de cette leçon, je sais distinguer les faits des opinions dans un article
\end{itemize}
\end{objectifs}

\begin{prerequis}
\begin{itemize}
\item Le vocabulaire de base de la nature et de la météo vu en Première (rain, drought, flood, season)
\item Le présent simple et le present continuous pour décrire des faits et des tendances
\item Les comparatifs et superlatifs (warmer, the hottest)
\item Les verbes modaux de nécessité et d'obligation (must, should, have to)
\item La lecture de textes courts avec questions de compréhension (Première, Chapter Environment)
\end{itemize}
\end{prerequis}

\newpage

\coursec{Warm-up et découverte du thème}

\courssub{Situation d'accroche : et si le climat changeait sous nos yeux ?}

Imagine you are in Bamenda during the rainy season and the floods destroy houses in your neighbourhood, or in Maroua where the heat reaches 45°C and the rains arrive later every year. Your grandmother says: “When I was young, the seasons were regular.” What has changed? Scientists all over the world, from London to Lagos, are warning us: our planet is getting warmer. Today, we are going to read an article about climate change and global warming to understand what is happening to our environment and what we can do about it.

Cette situation n'est pas de la science-fiction : elle décrit la réalité quotidienne de milliers de Camerounais. À Douala, les rues se transforment en rivières après chaque forte pluie ; dans l'Extrême-Nord, les bergers doivent marcher de plus en plus loin pour trouver de l'eau ; dans l'Ouest, les éboulements de terrain emportent des maisons entières après les pluies diluviennes. Tous ces phénomènes ont un point commun : ils sont liés à l'\cle{environnement} et à son évolution.

\textbf{Problématique de la leçon :} \emph{What is happening to our environment, and how can reading English texts help us understand and talk about it?} (Que se passe-t-il pour notre environnement, et comment la lecture de textes en anglais peut-elle nous aider à le comprendre et à en parler ?)

\courssub{Brainstorming : What environmental problems do you know?}

Avant de lire un texte, un bon lecteur active toujours ses connaissances sur le sujet. En anglais, cette étape s'appelle \cle{brainstorming} (littéralement « tempête de cerveaux », c'est-à-dire remue-méninges). Posons-nous la question : \eng{What environmental problems do you know?} (Quels problèmes environnementaux connais-tu ?)

Voici les réponses qu'une classe de Terminale à Yaoundé pourrait donner :

\begin{itemize}
\item \eng{Pollution} — la pollution (de l'air, de l'eau, des sols) ;
\item \eng{Deforestation} — la déforestation, la destruction des forêts ;
\item \eng{Floods} — les inondations ;
\item \eng{Drought} — la sécheresse ;
\item \eng{Global warming} — le réchauffement climatique ;
\item \eng{Climate change} — le changement climatique ;
\item \eng{Waste} — les déchets (plastique, ordures) ;
\item \eng{Landslides} — les glissements de terrain, éboulements ;
\item \eng{Desertification} — la désertification ;
\item \eng{Extinction of species} — l'extinction des espèces animales et végétales.
\end{itemize}

\begin{attention}[Prononciation et orthographe]
Attention à quelques mots piégeux : \eng{environment} se prononce « in-vaï-ron-ment » (le \emph{n} du milieu existe : envi\textbf{r}onment, ne l'oubliez pas à l'écrit !). \eng{Pollution} se prononce « po-lou-chonne ». \eng{Drought} se prononce « draoute » (les lettres \emph{-gh-} sont muettes). \eng{Flood} se prononce « fleud » et non « floude ».
\end{attention}

Organisons maintenant ce vocabulaire en carte mentale : c'est un excellent outil pour mémoriser le lexique d'un thème.

\begin{popfigure}
\begin{tikzpicture}[mindmap, grow cyclic, every node/.style={concept, align=center, font=\small}, concept color=popblue!40, level 1/.append style={level distance=3.4cm, sibling angle=90}]
\node[concept color=popblue!60, font=\small\bfseries] {ENVIRONMENTAL ISSUES}
  child[concept color=popgreen!50] { node[align=center]{climate\\change} }
  child[concept color=poporange!50] { node {pollution} }
  child[concept color=popgold!60] { node {deforestation} }
  child[concept color=poppink!50] { node[align=center]{global\\warming} };
\end{tikzpicture}
\legende{Carte mentale du vocabulaire : les grands problèmes environnementaux. Vous pouvez enrichir chaque branche avec des sous-idées (floods, drought, waste, landslides...).}
\end{popfigure}

\courssub{Observer avant de lire : la prédiction}

Un bon lecteur ne plonge jamais dans un texte « les yeux fermés ». Avant de lire, il observe les \cle{indices paratextuels} : le titre, la source, les images, les mots en gras. Cette stratégie s'appelle \cle{predicting} (prédire). Elle prépare le cerveau et facilite la compréhension.

Regardez le titre de l'article que nous allons étudier dans la prochaine section :

\begin{center}
\textbf{“A Warmer Planet: How Climate Change Is Transforming Life in Cameroon”}
\end{center}

\begin{methode}[Prédire le contenu d'un texte à partir de son titre]
\begin{enumerate}
\item \textbf{Repérez les mots-clés du titre.} Ici : \eng{warmer planet} (planète plus chaude), \eng{climate change} (changement climatique), \eng{transforming life} (transformer la vie), \eng{Cameroon} (Cameroun).
\item \textbf{Posez-vous des questions.} \eng{What is changing? Who is affected? What can we do?}
\item \textbf{Formulez une hypothèse.} Exemple : \eng{I think the article will describe the effects of climate change in Cameroon and maybe give some solutions.}
\item \textbf{Vérifiez votre prédiction pendant la lecture.} C'est la lecture elle-même qui confirmera ou corrigera votre hypothèse.
\end{enumerate}
\end{methode}

\begin{exemplebox}[Prédictions possibles à partir du titre]
\eng{The article will probably talk about rising temperatures in Cameroon.} — L'article parlera probablement de la hausse des températures au Cameroun.

\eng{It may mention floods in Douala and drought in the Far North.} — Il mentionnera peut-être les inondations à Douala et la sécheresse dans l'Extrême-Nord.

\eng{It will probably explain the causes of global warming.} — Il expliquera probablement les causes du réchauffement climatique.

Remarquez l'usage de \eng{will probably} et \eng{may} pour exprimer une prédiction incertaine : c'est un outil linguistique précieux à l'oral comme à l'écrit.
\end{exemplebox}

\courssub{Weather ou climate ? La distinction fondamentale}

Observons d'abord ces deux phrases prononcées par deux élèves :

\eng{The weather is hot today: 38°C in Garoua!} — Il fait chaud aujourd'hui : 38 °C à Garoua !

\eng{The climate of the Far North is hot and dry most of the year.} — Le climat de l'Extrême-Nord est chaud et sec la majeure partie de l'année.

Quelle différence voyez-vous ? Dans la première phrase, on parle d'\textbf{aujourd'hui}, d'un moment précis. Dans la seconde, on parle d'une \textbf{tendance générale} observée sur de longues périodes. C'est exactement la différence entre \cle{weather} et \cle{climate}.

\begin{definition}[Weather]
\eng{Weather} désigne \textbf{le temps qu'il fait} à un moment et à un endroit précis : pluie, soleil, vent, température du jour. C'est un phénomène de \textbf{court terme} (quelques heures, quelques jours).
Exemple : \eng{It rained heavily in Douala yesterday.} — Il a beaucoup plu à Douala hier.
\end{definition}

\begin{definition}[Climate]
\eng{Climate} désigne \textbf{le climat}, c'est-à-dire l'ensemble des conditions atmosphériques habituelles d'une région, observées sur une \textbf{longue période} (au moins trente ans pour les climatologues).
Exemple : \eng{Buea has a cool and rainy climate.} — Buea a un climat frais et pluvieux.
\end{definition}

\begin{attention}[Le faux ami des francophones : weather ≠ climate]
En français, le mot « temps » sert pour les deux notions : « le temps qu'il fait » et « le temps qu'il fait dans cette région ». En anglais, il faut \textbf{deux mots différents} :
\begin{itemize}
\item \eng{weather} = le temps (aujourd'hui, demain, cette semaine) ;
\item \eng{climate} = le climat (sur des années, des décennies).
\end{itemize}
Ne dites jamais \eng{The weather of Cameroon is tropical} : c'est faux, car un climat ne se décrit pas avec \eng{weather}. Dites : \eng{The climate of Cameroon is tropical.} De même, ne dites pas \eng{What is the climate like today?} mais \eng{What is the weather like today?} (Quel temps fait-il aujourd'hui ?).
\end{attention}

\begin{center}
\begin{tabularx}{\linewidth}{|X|X|X|}
\hline
\rowcolor{popblueL} Français & English & Remarque \\
\hline
Quel temps fait-il aujourd'hui ? & What is the weather like today? & court terme : \eng{weather} \\
\hline
Il pleut beaucoup cette semaine. & The weather is very rainy this week. & court terme : \eng{weather} \\
\hline
Le climat du Nord est sahélien. & The climate of the North is Sahelian. & long terme : \eng{climate} \\
\hline
Le climat change depuis des décennies. & The climate has been changing for decades. & long terme : \eng{climate} \\
\hline
Les prévisions météo & the weather forecast & « météo » se dit \eng{weather forecast} \\
\hline
\end{tabularx}
\end{center}

\begin{savaistu}[Cameroon, “Africa in miniature”]
Le Cameroun est parfois appelé \eng{Africa in miniature} (« l'Afrique en miniature »). Pourquoi ? Parce qu'il réunit presque tous les climats et tous les paysages du continent africain sur un seul territoire : le climat sahélien et les savanes sèches de l'Extrême-Nord (Maroua), les hautes terres fraîches de l'Ouest et du Nord-Ouest (Bamenda, Dschang), la forêt équatoriale humide du Sud et du Centre (Yaoundé, Ebolowa), et le littoral chaud et très pluvieux (Douala, Limbé, Kribi). C'est aussi l'un des rares pays d'Afrique où l'on trouve à la fois des plages, des montagnes, des forêts denses et des zones presque désertiques. Voilà pourquoi le Cameroun est un « laboratoire » idéal pour observer les effets du changement climatique !
\end{savaistu}

\courssub{Deux notions clés : climate change et global warming}

Ces deux expressions reviennent constamment dans les médias anglophones. Elles sont liées, mais elles ne sont \textbf{pas synonymes}. Observons d'abord deux phrases :

\eng{Global warming is the rise of the Earth's average temperature.} — Le réchauffement climatique est l'augmentation de la température moyenne de la Terre.

\eng{Climate change includes rising temperatures, but also stronger storms, floods and droughts.} — Le changement climatique comprend la hausse des températures, mais aussi des tempêtes plus violentes, des inondations et des sécheresses.

\begin{definition}[Climate change]
\eng{Climate change} means \textbf{long-term changes in temperatures and weather patterns on Earth} — des changements durables des températures et des régimes climatiques (saisons décalées, pluies irrégulières, événements extrêmes plus fréquents). C'est la notion la plus \textbf{large}.
\end{definition}

\begin{definition}[Global warming]
\eng{Global warming} means \textbf{the gradual increase of the Earth's average temperature, mainly caused by human activities} — l'augmentation progressive de la température moyenne de la planète, provoquée surtout par les activités humaines (combustibles fossiles, déforestation, industrie). C'est une \textbf{composante} du changement climatique.
\end{definition}

\begin{aretenir}[La relation entre les deux notions]
\eng{Global warming} (le thermomètre qui monte) est une \textbf{cause} et une \textbf{partie} de \eng{climate change} (toutes les transformations qui en résultent). Autrement dit : la planète se réchauffe (\eng{global warming}), et ce réchauffement provoque des changements plus vastes dans les climats du monde (\eng{climate change}) : fonte des glaciers, montée du niveau des mers, sécheresses plus longues, pluies plus violentes.
\end{aretenir}

\courssub{Un premier texte de découverte}

Lisons maintenant un court texte qui met en scène ces deux notions dans un contexte camerounais. C'est un avant-goût de l'article que nous étudierons en détail dans la section suivante.

\begin{textebox}[Reading for understanding — “Grandmother Aissatou's Question”]
Grandmother Aissatou lives in Maroua, in the Far North of Cameroon. Every morning, she listens to the weather forecast on the radio. Last week, the temperature reached 44 degrees Celsius, and her grandson Moussa could not sleep at night because of the heat.

“When I was a young girl,” she says, “the rainy season always started in May. The farmers knew exactly when to plant their millet. Today, nobody knows when the rains will come. Sometimes they arrive in July, sometimes they come with terrible storms that destroy the fields.”

Moussa learned about climate change at school. He explains to his grandmother: “Grandma, scientists say the Earth is getting warmer. This is called global warming. Because of this warming, the climate is changing everywhere: the seasons are not regular any more, droughts are longer here in the North, and floods are more frequent in Douala.”

Grandmother Aissatou thinks for a moment and asks: “But who is warming the planet, my child?” Moussa answers: “We all are, Grandma. Cars, factories, bush fires and the cutting of trees produce gases that trap the heat around the Earth, like a blanket.”

The old woman smiles. “Then we must all help to remove that blanket,” she says. “Tell me what we can do.”
\end{textebox}

\textbf{Glossaire :}

\begin{center}
\begin{tabularx}{\linewidth}{|X|l|X|}
\hline
\rowcolor{popblueL} Mot anglais & Nature & Traduction \\
\hline
weather forecast & nom & bulletin météo, prévisions météo \\
\hline
heat & nom & chaleur \\
\hline
rainy season & nom & saison des pluies \\
\hline
millet & nom & mil (céréale) \\
\hline
drought & nom & sécheresse \\
\hline
flood & nom & inondation \\
\hline
bush fire & nom & feu de brousse \\
\hline
to trap & verbe & piéger, retenir \\
\hline
blanket & nom & couverture \\
\hline
\end{tabularx}
\end{center}

\textbf{Questions de compréhension rapide (oral) :}

\begin{enumerate}
\item \eng{Where does Grandmother Aissatou live?}
\item \eng{What has changed about the rainy season, according to her?}
\item \eng{What is the difference between global warming and climate change in Moussa's explanation?}
\item \eng{What human activities warm the planet, according to Moussa?}
\item \eng{What image does Moussa use to explain the greenhouse effect?}
\end{enumerate}

\courssub{Activer ses connaissances : le cas du Cameroun}

Pour terminer cette mise en route, relions le thème à notre expérience directe. Voici trois situations réelles que vous pouvez décrire en anglais :

\begin{itemize}
\item \eng{In Douala, heavy rains often cause floods. Streets and houses are full of water, and people cannot go to work or to school.} — À Douala, les fortes pluies provoquent souvent des inondations. Les rues et les maisons sont pleines d'eau et les gens ne peuvent pas aller travailler ni à l'école.
\item \eng{In the Far North, the drought is getting worse. The rainy season is shorter, and farmers cannot grow enough food.} — Dans l'Extrême-Nord, la sécheresse s'aggrave. La saison des pluies est plus courte et les agriculteurs ne peuvent pas produire assez de nourriture.
\item \eng{In the West Region, landslides sometimes destroy houses after heavy rains, especially on the hills around Dschang and Bafoussam.} — Dans la région de l'Ouest, des éboulements détruisent parfois des maisons après de fortes pluies, surtout sur les collines autour de Dschang et de Bafoussam.
\end{itemize}

\begin{experience}[Class discussion — “Have you noticed?”]
En groupes de quatre, discutez en anglais pendant cinq minutes à partir de ces questions :
\begin{enumerate}
\item \eng{Have you noticed changes in the weather in your town or village? Give examples.}
\item \eng{Do the seasons start at the same time as ten years ago?}
\item \eng{What environmental problem is the most serious in your region? Why?}
\end{enumerate}
Chaque groupe choisit un rapporteur qui présente les réponses à la classe en commençant par : \eng{In our group, we think that...}
\end{experience}

\courssub{Exercice résolu et entraînement}

\begin{exoresolu}[Weather or climate?]
Complétez chaque phrase avec \eng{weather} ou \eng{climate}.

1. \eng{The \dots\ in Buea is cool and humid all year round.}

2. \eng{What terrible \dots\ ! It has been raining since this morning.}

3. \eng{The \dots\ of the Far North is becoming hotter and drier.}

4. \eng{Listen to the \dots\ forecast: it will be sunny tomorrow.}

5. \eng{Scientists study \dots\ change over long periods.}
\tcblower
\textbf{Solution détaillée :}

1. \eng{climate} — on décrit les conditions habituelles d'une région (long terme).

2. \eng{weather} — il s'agit du temps d'aujourd'hui (court terme).

3. \eng{climate} — une évolution sur des années, donc long terme.

4. \eng{weather} — l'expression figée est \eng{weather forecast} (prévisions météo).

5. \eng{climate} — l'expression consacrée est \eng{climate change} (changement climatique).

\textbf{Astuce :} demandez-vous toujours « aujourd'hui ou sur des années ? ». Si la réponse est « aujourd'hui / cette semaine », choisissez \eng{weather} ; si c'est « sur des années / en général », choisissez \eng{climate}.
\end{exoresolu}

\begin{exercice}[À vous de jouer !]
1. Traduisez en anglais :
a) Le climat du Cameroun est très varié.
b) Il fait très chaud aujourd'hui à Maroua.
c) Le réchauffement climatique est causé par les activités humaines.
d) Les inondations sont de plus en plus fréquentes à Douala.

2. Répondez en deux ou trois phrases en anglais : \eng{What environmental problem affects your region the most? What can young people do about it?}
\end{exercice}

\begin{corrige}[Exercice « À vous de jouer ! »]
1. a) \eng{The climate of Cameroon is very varied.} (et non \eng{weather} : il s'agit d'une réalité permanente.)
b) \eng{It is very hot today in Maroua.} ou \eng{The weather is very hot today in Maroua.}
c) \eng{Global warming is caused by human activities.} (voix passive : \emph{is caused} ; attention à ne pas écrire \eng{by the humans activities} : pas d'article, pas de pluriel à \eng{human} employé comme adjectif.)
d) \eng{Floods are more and more frequent in Douala.} ou \eng{Floods are becoming more and more frequent in Douala.}

2. Exemple de réponse attendue : \eng{In my region, the most serious problem is deforestation, because people cut trees for firewood and farms. Young people can plant trees, sensitise their families and avoid bush fires.} Toute réponse correcte sur le plan grammatical et pertinente sur le fond est acceptée.
\end{corrige}

\begin{aretenir}[Ce qu'il faut retenir de cette mise en route]
\begin{itemize}
\item Avant de lire, on \textbf{active ses connaissances} (brainstorming) et on \textbf{prédit} le contenu à partir du titre et des images.
\item \eng{Weather} = le temps qu'il fait (court terme) ; \eng{climate} = le climat d'une région (long terme). Ne jamais les confondre !
\item \eng{Global warming} = la hausse progressive de la température moyenne de la Terre ; \eng{climate change} = l'ensemble des changements durables du climat qui en découlent.
\item Le Cameroun, « Afrique en miniature », est directement concerné : inondations à Douala, sécheresse dans l'Extrême-Nord, éboulements dans l'Ouest.
\end{itemize}
\end{aretenir}

\coursec{Texte support : article de lecture}

Dans la section précédente, nous avons activé nos connaissances sur l'environnement : vous avez observé des images, discuté du temps qu'il fait à Yaoundé et à Maroua, et émis des hypothèses sur le contenu de l'article à partir de son titre. Il est maintenant temps de \cle{vérifier ces prédictions} en lisant le texte lui-même. Nous allons procéder en trois temps : une première lecture globale (\eng{skimming}), une deuxième lecture ciblée (\eng{scanning}), puis une relecture collective avec le glossaire.

\courssub{Lecture silencieuse de l'article}

\begin{textebox}[Our Planet is Heating Up — article de lecture]
The Earth is getting warmer. According to scientists, the average global temperature has risen by about 1.1°C since 1880. This phenomenon, known as global warming, is mainly caused by human activities. When we burn fossil fuels such as coal, oil and gas to produce energy, we release greenhouse gases like carbon dioxide into the atmosphere. These gases trap the sun's heat, just like the glass of a greenhouse.

The consequences are already visible. Glaciers are melting, sea levels are rising, and extreme weather events such as floods, droughts and heatwaves are becoming more frequent. In Africa, farmers suffer from irregular rainfall: in northern Cameroon, for example, the rainy season is shorter than before, which threatens harvests and food security.

However, it is not too late to act. Governments can invest in renewable energy like solar and wind power. As individuals, we can plant trees, save electricity, use public transport and recycle our waste. Every small action counts. As the saying goes, we do not inherit the Earth from our ancestors; we borrow it from our children.
\end{textebox}

\textbf{Consigne de lecture silencieuse (3 minutes).} Lisez le texte une première fois \emph{sans dictionnaire}, sans vous arrêter sur les mots inconnus. Soulignez mentalement les mots que vous comprenez et entourez ceux que vous devrez deviner. L'objectif n'est pas de tout comprendre, mais de saisir \cle{l'idée générale}.

\courssub{Première lecture globale : le skimming}

\begin{definition}[Le skimming]
Le \cle{skimming} (lecture rapide, « survol ») consiste à parcourir un texte rapidement pour en dégager \cle{l'idée générale} (\eng{the gist}) et \cle{l'attitude de l'auteur} (\eng{the writer's attitude}), sans chercher à comprendre chaque mot. On lit surtout le titre, la première phrase de chaque paragraphe et la conclusion.
\end{definition}

Appliquons cette stratégie à notre article :

\begin{itemize}
\item \textbf{Le titre} — \eng{“Our Planet is Heating Up”} : le mot \eng{planet} annonce un sujet mondial ; \eng{heating up} suggère un problème de température. Prédiction : le texte parlera du réchauffement climatique.
\item \textbf{La première phrase} — \eng{“The Earth is getting warmer.”} : le thème général est confirmé dès la première ligne.
\item \textbf{La dernière phrase} — \eng{“we borrow it from our children”} : l'auteur termine sur un appel à la responsabilité, donc sur une note d'\emph{espoir} et non de désespoir.
\end{itemize}

\begin{exemplebox}[Résultat du skimming]
\textbf{General theme:} global warming and climate change.\\
\textbf{Writer's attitude:} concerned but optimistic — he presents the problem seriously, but he believes solutions exist (\eng{“it is not too late to act”}).
\end{exemplebox}

\begin{attention}[Ne confondez pas thème et attitude]
Beaucoup d'élèves répondent à la question \eng{“What is the writer's attitude?”} en répétant le thème. L'\cle{attitude} est un sentiment ou une position : \eng{concerned} (inquiet), \eng{optimistic} (optimiste), \eng{critical} (critique), \eng{neutral} (neutre). Ici, la bonne réponse est : \eng{“The writer is concerned about global warming but optimistic about the solutions.”}
\end{attention}

\courssub{Deuxième lecture ciblée : le scanning}

\begin{definition}[Le scanning]
Le \cle{scanning} (lecture ciblée) consiste à chercher dans le texte une \cle{information précise} (un chiffre, une date, un nom, une cause, une conséquence) en laissant l'œil glisser rapidement sur les lignes, sans tout relire. On s'arrête uniquement sur le passage qui contient l'information recherchée.
\end{definition}

Cherchons maintenant dans le texte les quatre catégories d'informations typiques d'un article sur le climat :

\begin{center}
\begin{tabularx}{\linewidth}{|l|X|}
\hline
\rowcolor{popblueL} \textbf{Information} & \textbf{Ce que dit le texte} \\
\hline
Chiffre clé & \eng{the average global temperature has risen by about 1.1°C since 1880} \\
\hline
Causes & \eng{burning fossil fuels} (coal, oil, gas) $\rightarrow$ \eng{greenhouse gases} trap the sun's heat \\
\hline
Conséquences & \eng{glaciers are melting, sea levels are rising, floods, droughts, heatwaves} ; en Afrique : \eng{irregular rainfall}, récoltes menacées \\
\hline
Solutions & \eng{renewable energy} (solar, wind) ; \eng{plant trees, save electricity, use public transport, recycle waste} \\
\hline
\end{tabularx}
\end{center}

\begin{exoresolu}[Repérer une information précise]
Énoncé — \eng{“Scan the text and answer: Why is food security threatened in northern Cameroon?”}
\tcblower
Solution détaillée — Étape 1 : je repère dans le texte les mots-clés de la question : \eng{northern Cameroon}, \eng{food security}. Étape 2 : je localise le passage : \eng{“in northern Cameroon, for example, the rainy season is shorter than before, which threatens harvests and food security”}. Étape 3 : je reformule la réponse avec mes propres mots : \eng{“Food security is threatened because the rainy season is shorter than before, so harvests are at risk.”} (« La sécurité alimentaire est menacée parce que la saison des pluies est plus courte qu'avant, ce qui met les récoltes en danger. »)
\end{exoresolu}

\courssub{La structure du texte}

Un bon article d'information suit presque toujours la même architecture. La reconnaître vous aide à prédire où se trouve chaque information avant même de lire.

\begin{popfigure}
\begin{center}
\begin{tikzpicture}[
  box/.style={draw=popblue, fill=popblueL, rounded corners=3pt, align=center, minimum width=3.1cm, minimum height=1cm, font=\small},
  arr/.style={-{Stealth[length=2.5mm]}, thick, poporange}
]
\node[box, align=center] (intro) at (0,0){\textbf{Introduction}\\ the problem:\\ global warming};
\node[box, align=center] (causes) at (3.9,0){\textbf{Causes}\\ fossil fuels,\\ greenhouse gases};
\node[box, align=center] (effects) at (7.8,0){\textbf{Effects}\\ melting glaciers,\\ droughts, heatwaves};
\node[box, align=center] (solutions) at (3.9,-2.2){\textbf{Solutions}\\ renewable energy,\\ small actions};
\node[box, fill=popgreenL, draw=popgreen, align=center] (concl) at (0,-2.2){\textbf{Conclusion}\\ hope: it is not\\ too late to act};
\draw[arr] (intro) -- (causes);
\draw[arr] (causes) -- (effects);
\draw[arr] (effects.south) |- (solutions.east);
\draw[arr] (solutions) -- (concl);
\end{tikzpicture}
\end{center}
\legende{Structure de l'article : problème, causes, effets, solutions, conclusion pleine d'espoir.}
\end{popfigure}

\begin{aretenir}[À retenir]
Un article sur l'environnement suit généralement le plan : \cle{problem} $\rightarrow$ \cle{causes} $\rightarrow$ \cle{effects} $\rightarrow$ \cle{solutions} $\rightarrow$ \cle{conclusion}. Pour trouver une cause, cherchez dans le premier tiers du texte ; pour une solution, cherchez vers la fin.
\end{aretenir}

\courssub{Glossaire des mots difficiles}

Avant la relecture collective, levons les obstacles lexicaux. Observez comment chaque mot est employé dans sa phrase, puis mémorisez le tableau.

\begin{center}
\begin{tabularx}{\linewidth}{|l|l|X|}
\hline
\rowcolor{popblueL} \textbf{English} & \textbf{Français} & \textbf{Exemple tiré du texte} \\
\hline
to heat up & se réchauffer & \eng{Our planet is heating up.} \\
\hline
fossil fuels & combustibles fossiles & \eng{We burn fossil fuels such as coal, oil and gas.} \\
\hline
to release & rejeter, libérer & \eng{We release greenhouse gases into the atmosphere.} \\
\hline
greenhouse gases & gaz à effet de serre & \eng{These gases trap the sun's heat.} \\
\hline
to trap & piéger, retenir & \eng{These gases trap the sun's heat.} \\
\hline
glacier & glacier & \eng{Glaciers are melting.} \\
\hline
sea level & niveau de la mer & \eng{Sea levels are rising.} \\
\hline
drought & sécheresse & \eng{Droughts are becoming more frequent.} \\
\hline
heatwave & vague de chaleur & \eng{Floods, droughts and heatwaves.} \\
\hline
rainfall & précipitations, pluies & \eng{Farmers suffer from irregular rainfall.} \\
\hline
harvest & récolte & \eng{It threatens harvests and food security.} \\
\hline
renewable energy & énergie renouvelable & \eng{Governments can invest in renewable energy.} \\
\hline
waste & déchets & \eng{We can recycle our waste.} \\
\hline
to inherit & hériter & \eng{We do not inherit the Earth from our ancestors.} \\
\hline
ancestor & ancêtre & \eng{We borrow it from our children.} \\
\hline
\end{tabularx}
\end{center}

\begin{attention}[Faux amis et pièges de prononciation]
\begin{itemize}
\item \eng{glacier} se prononce « gla-sieur » en anglais (/ˈɡleɪʃə(r)/), pas comme le français \emph{glacier} (marchand de glaces, /ɡlasje/) !
\item \eng{waste} signifie « déchets » (nom) ou « gaspiller » (verbe : \eng{Don't waste water!} = « Ne gaspille pas l'eau ! »). Ce n'est jamais « la taille » (qui se dit \eng{waist}).
\item \eng{sensible} et \eng{sensitive} : attention, \eng{sensitive} = « sensible » (fragile, réactif) ; « sensé, raisonnable » se dit \eng{sensible}. On parlera d'une \eng{environmentally sensitive area} (zone écologiquement fragile).
\item \eng{to inherit} : on hérite \emph{de} quelqu'un, mais en anglais on dit \eng{to inherit something FROM somebody}, jamais \emph{inherit of}.
\end{itemize}
\end{attention}

\courssub{Relecture collective et vérification des hypothèses}

Relisons maintenant le texte à voix haute, paragraphe par paragraphe, et confrontons nos prédictions du warm-up à la réalité du texte.

\begin{exemplebox}[Vérification des prédictions]
\begin{itemize}
\item Prédiction 1 : \eng{“The text is about pollution.”} $\rightarrow$ Partiellement vrai : le texte parle du réchauffement climatique, dont la pollution par les gaz est une cause.
\item Prédiction 2 : \eng{“The text mentions Africa or Cameroon.”} $\rightarrow$ Vrai : \eng{“in northern Cameroon, the rainy season is shorter than before”}.
\item Prédiction 3 : \eng{“The text is pessimistic.”} $\rightarrow$ Faux : la conclusion est pleine d'espoir (\eng{“it is not too late to act”}).
\end{itemize}
\end{exemplebox}

\begin{methode}[Comment lire un article au Bac — les 5 étapes]
\begin{enumerate}
\item \textbf{Lire le titre et prédire} : avant de lire, formulez deux ou trois hypothèses sur le contenu (\eng{“This text is probably about...”}).
\item \textbf{Skimming} : lisez rapidement le texte (titre, premières phrases, conclusion) pour dégager l'idée générale et l'attitude de l'auteur.
\item \textbf{Lire les questions} : avant la deuxième lecture, lisez toutes les questions de compréhension et soulignez leurs mots-clés.
\item \textbf{Scanning} : cherchez dans le texte les mots-clés des questions pour localiser précisément les réponses.
\item \textbf{Vérifier} : relisez le passage concerné pour confirmer chaque réponse, puis rédigez-la avec vos propres mots, sans recopier de longues phrases du texte.
\end{enumerate}
\end{methode}

\begin{exercice}[À vous de jouer]
Répondez en phrases complètes, en anglais :
\begin{enumerate}
\item \eng{What is the main cause of global warming according to the text?}
\item \eng{ Mention two consequences of climate change visible in the world today.}
\item \eng{ What can individuals do to fight global warming? Give two examples.}
\item \eng{ In your opinion, why does the writer end the text with a saying about our children?}
\end{enumerate}
\end{exercice}

\begin{corrige}[Exercice « À vous de jouer »]
\begin{enumerate}
\item \eng{Global warming is mainly caused by human activities, especially the burning of fossil fuels which releases greenhouse gases into the atmosphere.}
\item \eng{Glaciers are melting and sea levels are rising. Extreme weather events such as floods, droughts and heatwaves are also becoming more frequent.}
\item \eng{As individuals, we can plant trees and save electricity. We can also use public transport and recycle our waste.}
\item \eng{The writer ends with this saying to make readers feel responsible: the Earth does not belong to us, we must protect it for future generations.}
\end{enumerate}
\end{corrige}

\begin{savaistu}[D'où vient l'expression « greenhouse effect » ?]
L'expression \eng{greenhouse effect} (« effet de serre ») vient de la comparaison avec une serre de jardinier : le verre laisse entrer la lumière du soleil mais empêche la chaleur de ressortir. Les gaz comme le \eng{carbon dioxide} (CO\textsubscript{2}) jouent exactement ce rôle de « vitre » autour de la planète. Au Cameroun anglophone, les élèves de Buea et Bamenda étudient ce phénomène dans le cadre du programme de \eng{Environmental Education}, et de nombreuses écoles organisent des campagnes de \eng{tree planting} (plantation d'arbres) pour lutter contre la déforestation.
\end{savaistu}

Nous avons désormais lu le texte en profondeur, vérifié nos prédictions et enrichi notre vocabulaire. La section suivante vous proposera des questions de compréhension et d'interprétation plus approfondies pour tester votre maîtrise de l'article.

\coursec{Compréhension du texte}

Dans la section précédente, nous avons découvert le texte support, un article de presse intitulé \eng{“The Heat Is On: Cameroon Faces a Changing Climate”}. Nous l'avons lu une première fois pour en saisir l'idée générale (\eng{skimming}). Il est maintenant temps de passer à la \cle{lecture approfondie} : répondre aux questions de compréhension, comme on vous le demandera au baccalauréat. Cette section vous donne les stratégies précises pour chaque type de question, avec des exemples corrigés.

\courssub{Le texte de référence (rappel)}

Pour travailler sereinement, relisons l'article. Les paragraphes sont numérotés : au baccalauréat, on vous demandera souvent de justifier vos réponses en citant le numéro du paragraphe ou la ligne précise.

\begin{textebox}[The Heat Is On: Cameroon Faces a Changing Climate]
(1) In the village of Maroua, in the Far North of Cameroon, farmers say the rains no longer come when they should. “Twenty years ago, we knew exactly when to plant our millet,” explains Amadou, a sixty-year-old farmer. “Now the seasons are confused, and our harvests are smaller every year.”

(2) Amadou's experience is not unique. Across the world, climate change is transforming the lives of millions of people. Scientists agree that global warming — the slow rise of the Earth's average temperature — is mainly caused by human activities. When we burn petrol in cars and generators, or when forests are cut down, gases such as carbon dioxide accumulate in the atmosphere. These gases act like a blanket: they keep the sun's heat close to the Earth, and the planet becomes warmer.

(3) The consequences are already visible. In the North, droughts last longer and the desert is advancing. On the coast, around Douala and Kribi, the sea level is rising and floods destroy houses and roads. Farmers lose their crops, fishermen catch fewer fish, and diseases such as malaria spread to new regions.

(4) However, it is not too late to act. Simple solutions exist: planting trees, using cleaner energy, and educating young people about the environment. In Yaoundé, several schools have created “green clubs” where students plant trees and recycle plastic. “The Earth is our only home,” says one student. “If we destroy it, we will have nowhere else to go.”
\end{textebox}

\textbf{Glossaire}

\begin{center}
\begin{tabularx}{\linewidth}{|l|l|X|}
\hline
\rowcolor{popblueL} \textbf{Mot anglais} & \textbf{Nature} & \textbf{Traduction} \\
\hline
millet & nom & le mil (céréale) \\
\hline
harvest & nom & la récolte \\
\hline
to burn & verbe & brûler \\
\hline
blanket & nom & la couverture \\
\hline
drought & nom & la sécheresse \\
\hline
flood & nom & l'inondation \\
\hline
crop & nom & la culture, la récolte \\
\hline
to spread & verbe & se propager \\
\hline
\end{tabularx}
\end{center}

\courssub{Les questions True or False : toujours justifier}

\begin{methode}[Répondre à une question True or False]
\begin{enumerate}
\item Lisez l'affirmation et soulignez les mots-clés.
\item Repérez dans le texte la phrase qui parle de la même idée (\eng{scanning}).
\item Comparez mot à mot : attention aux petits mots qui changent tout (\eng{only}, \eng{never}, \eng{all}, \eng{always}).
\item Écrivez \eng{True} ou \eng{False}, puis \textbf{citez le passage exact du texte} entre guillemets anglais “ ”, en indiquant le numéro du paragraphe.
\end{enumerate}
\end{methode}

\begin{exemplebox}[Exemples traités]
\textbf{Q1.} \eng{Global warming is caused only by natural factors. True or False?}\par
\textbf{Réponse :} \eng{False. The text says global warming is “mainly caused by human activities” (paragraph 2).}\par
Le mot \eng{only} rend l'affirmation fausse : le texte dit \eng{mainly} (principalement), pas \eng{only} (uniquement).\par\medskip
\textbf{Q2.} \eng{The sea level is rising near Douala. True or False?}\par
\textbf{Réponse :} \eng{True. The text says “On the coast, around Douala and Kribi, the sea level is rising” (paragraph 3).}
\end{exemplebox}

\begin{attention}[Une réponse sans justification perd des points]
Au baccalauréat, un \eng{True} ou un \eng{False} seul, sans citation du texte, ne rapporte généralement \textbf{aucun point}, même si la réponse est correcte. La justification doit être une \textbf{citation exacte} du texte entre guillemets anglais “ ”, pas une phrase inventée de votre cru. Écrivez : \eng{False. The text says “...” (paragraph 2).} et non une paraphrase vague.
\end{attention}

\courssub{Les QCM : idées principales et détails}

Les QCM (\eng{multiple choice questions}) mélangent des questions globales (\eng{What is the text about?}) et des questions de détail. La stratégie : éliminez d'abord les options clairement fausses, puis vérifiez les deux restantes dans le texte.

\begin{exemplebox}[Exemples traités]
\textbf{Q1.} \eng{The text is mainly about:}\par
a) \eng{the history of Maroua} \quad b) \eng{climate change and its effects in Cameroon} \quad c) \eng{how to plant millet}\par
\textbf{Réponse : b.} Tout le texte parle du changement climatique, de ses causes (paragraphe 2) et de ses conséquences (paragraphe 3).\par\medskip
\textbf{Q2.} \eng{According to the text, greenhouse gases act like:}\par
a) \eng{a mirror} \quad b) \eng{a blanket} \quad c) \eng{a wall}\par
\textbf{Réponse : b.} Paragraphe 2 : \eng{“These gases act like a blanket.”}
\end{exemplebox}

\courssub{Les questions ouvertes (Wh- questions)}

Pour les questions commençant par \eng{What}, \eng{Why}, \eng{How}, répondez par une \textbf{phrase complète} en anglais, en reprenant les mots du texte. Une réponse d'un seul mot est souvent insuffisante.

\begin{exemplebox}[Exemples traités]
\textbf{Q1.} \eng{Why are Amadou's harvests smaller every year?}\par
\textbf{Réponse :} \eng{Because the seasons are confused and the rains no longer come when they should.}\par\medskip
\textbf{Q2.} \eng{What solutions does the text mention?}\par
\textbf{Réponse :} \eng{The text mentions planting trees, using cleaner energy and educating young people about the environment.}\par\medskip
\textbf{Q3.} \eng{How do greenhouse gases warm the planet?}\par
\textbf{Réponse :} \eng{They act like a blanket: they keep the sun's heat close to the Earth.}
\end{exemplebox}

\courssub{Les questions de vocabulaire en contexte}

On vous demandera de trouver dans le texte un synonyme ou un équivalent : \eng{Find in the text a word or expression that means...}

\begin{propriete}[Déduire le sens d'un mot inconnu : les 3 indices]
\begin{enumerate}
\item \textbf{Le contexte de la phrase} : les mots autour donnent le sens. \eng{“droughts last longer and the desert is advancing”} → \eng{drought} est lié à la sécheresse.
\item \textbf{La famille du mot} : \eng{warm} → \eng{warming} → \eng{global warming} ; \eng{to flood} → \eng{floods}.
\item \textbf{Le champ lexical voisin} : \eng{harvests, crops, farmers, plant} appartiennent tous au champ de l'agriculture.
\end{enumerate}
\end{propriete}

\begin{exemplebox}[Exemples traités]
\textbf{Q1.} \eng{Find in the text a word that means “récoltes” (paragraph 1).} → \eng{harvests}.\par
\textbf{Q2.} \eng{Find in the text an expression that means “il n'est pas trop tard” (paragraph 4).} → \eng{it is not too late to act}.\par
\textbf{Q3.} \eng{Find in the text a synonym of “to increase, to go up” (paragraph 3).} → \eng{is rising}.
\end{exemplebox}

\courssub{Les questions d'interprétation et de référence}

Deux types de questions exigent de réfléchir au-delà des mots :

\begin{itemize}
\item \textbf{L'attitude de l'auteur} : \eng{What is the writer's attitude?} — Est-il optimiste, pessimiste, neutre, inquiet ? Cherchez les indices : ici, la fin du texte (\eng{“it is not too late to act”}, \eng{“Simple solutions exist”}) montre une attitude \eng{worried but hopeful} (inquiète mais pleine d'espoir).
\item \textbf{Les questions de référence} : \eng{What does “it” refer to in paragraph X?}
\end{itemize}

\begin{methode}[Répondre à une question de référence (it, they, this...)]
\begin{enumerate}
\item Repérez le pronom dans le paragraphe indiqué.
\item Regardez le \textbf{nom qui le précède} dans la même phrase ou dans la phrase précédente.
\item Vérifiez : remplacez le pronom par ce nom ; la phrase doit rester logique.
\end{enumerate}
\textbf{Exemple traité :} \eng{What does “it” refer to in “If we destroy it, we will have nowhere else to go” (paragraph 4)?}\par
\textbf{Réponse :} \eng{“It” refers to the Earth.} Vérification : \eng{“If we destroy the Earth, we will have nowhere else to go”} — la phrase est logique.\par\medskip
\textbf{Autre exemple :} \eng{What does “it” refer to in “it is not too late to act”?} → \eng{It} renvoie à l'idée d'agir contre le réchauffement : \eng{it is not too late to act (against global warming)}.
\end{methode}

\begin{popfigure}
\begin{tikzpicture}[node distance=1.2cm and 1.5cm]
\node[draw=popblue, fill=popblueL, rounded corners, text width=4.6cm, align=center, font=\small\bfseries] (t1) {Type of question};
\node[draw=poporange, fill=poporangeL, rounded corners, text width=6.8cm, align=center, font=\small\bfseries] (s1) [right=of t1] {Strategy};
\node[draw=popblue, fill=white, rounded corners, text width=4.6cm, align=center, font=\small] (t2) [below=of t1] {True or False};
\node[draw=poporange, fill=white, rounded corners, text width=6.8cm, align=center, font=\small] (s2) [right=of t2] {Repérer la phrase exacte dans le texte et la citer entre guillemets};
\node[draw=popblue, fill=white, rounded corners, text width=4.6cm, align=center, font=\small] (t3) [below=of t2] {Reference question (it, they, this)};
\node[draw=poporange, fill=white, rounded corners, text width=6.8cm, align=center, font=\small] (s3) [right=of t3] {Regarder le nom qui précède le pronom, puis vérifier en le remplaçant};
\node[draw=popblue, fill=white, rounded corners, text width=4.6cm, align=center, font=\small] (t4) [below=of t3] {Vocabulary question};
\node[draw=poporange, fill=white, rounded corners, text width=6.8cm, align=center, font=\small] (s4) [right=of t4] {Lire la phrase entière et utiliser le contexte, la famille du mot et le champ lexical};
\node[draw=popblue, fill=white, rounded corners, text width=4.6cm, align=center, font=\small] (t5) [below=of t4] {Wh- question (What, Why, How)};
\node[draw=poporange, fill=white, rounded corners, text width=6.8cm, align=center, font=\small] (s5) [right=of t5] {Répondre par une phrase complète avec les mots du texte};
\end{tikzpicture}
\legende{Chaque type de question a sa stratégie de lecture.}
\end{popfigure}

\courssub{Tableau récapitulatif : les formulations de l'examen}

\begin{center}
\begin{tabularx}{\linewidth}{|X|X|}
\hline
\rowcolor{popblueL} \textbf{Consigne anglaise} & \textbf{Ce qu'on vous demande} \\
\hline
\eng{True or False? Justify your answer.} & Vrai ou faux + citation exacte du texte \\
\hline
\eng{Choose the correct answer (a, b or c).} & QCM : une seule bonne option \\
\hline
\eng{Answer the following questions.} & Phrases complètes en anglais \\
\hline
\eng{Find in the text a word that means...} & Synonyme ou équivalent dans le texte \\
\hline
\eng{What does “it” refer to in paragraph 4?} & Le nom représenté par le pronom \\
\hline
\eng{What is the writer's attitude?} & Le ton de l'auteur (worried, hopeful...) \\
\hline
\end{tabularx}
\end{center}

\courssub{Exercice résolu et correction collective}

\begin{exoresolu}[Compréhension globale du texte]
\textbf{Énoncé.} Répondez aux questions suivantes sur le texte \eng{“The Heat Is On”}.\par
1. \eng{The desert is advancing in the Far North of Cameroon. True or False? Justify.}\par
2. \eng{What do students do in the “green clubs” of Yaoundé?}\par
3. \eng{Find in the text a word that means “inondations” (paragraph 3).}\par
4. \eng{What does “they” refer to in “they keep the sun's heat close to the Earth” (paragraph 2)?}
\tcblower
\textbf{Solution détaillée.}\par
1. \eng{True. The text says “In the North, droughts last longer and the desert is advancing” (paragraph 3).} — La citation exacte justifie la réponse.\par
2. \eng{In the green clubs, students plant trees and recycle plastic.} — Phrase complète construite à partir du paragraphe 4.\par
3. \eng{floods} — Paragraphe 3 : \eng{“floods destroy houses and roads”} ; le contexte (mer qui monte, destructions) confirme le sens.\par
4. \eng{“They” refers to the gases (such as carbon dioxide).} — Le nom qui précède le pronom est \eng{these gases} ; vérification : \eng{“the gases keep the sun's heat close to the Earth”} — logique.
\end{exoresolu}

\begin{attention}[Pièges fréquents des francophones]
\begin{itemize}
\item Ne traduisez pas la question, répondez \textbf{en anglais}.
\item Dans un vrai/faux, méfiez-vous des mots absolus : \eng{only}, \eng{never}, \eng{all}, \eng{always} rendent souvent l'affirmation fausse.
\item Ne recopiez pas trois lignes du texte pour justifier : citez uniquement le passage pertinent.
\item Pour les questions de référence, ne répondez jamais par un verbe ou une phrase entière : un pronom renvoie à un \textbf{nom} (ou à une idée nominale).
\end{itemize}
\end{attention}

\begin{exercice}[À vous de jouer]
1. \eng{Amadou is a fisherman. True or False? Justify your answer with a line from the text.}\par
2. \eng{Why do greenhouse gases make the planet warmer?}\par
3. \eng{Find in the text an expression that means “il n'est pas trop tard”.}\par
4. \eng{What does “it” refer to in “If we destroy it” (paragraph 4)?}\par
5. \eng{What is the writer's attitude at the end of the text: hopeless or hopeful? Justify.}
\end{exercice}

\begin{corrige}[Exercice : À vous de jouer]
1. \eng{False. The text says Amadou is “a sixty-year-old farmer” (paragraph 1).}\par
2. \eng{They act like a blanket and keep the sun's heat close to the Earth.}\par
3. \eng{“It is not too late to act” (paragraph 4).}\par
4. \eng{“It” refers to the Earth.}\par
5. \eng{The writer is hopeful, because the text says “it is not too late to act” and “Simple solutions exist” (paragraph 4).}
\end{corrige}

\begin{aretenir}[L'essentiel de la section]
Comprendre un texte, c'est savoir adapter sa lecture à la question posée : citation exacte pour les vrai/faux, phrase complète pour les questions ouvertes, retour en arrière pour les pronoms de référence, et trois indices (contexte, famille, champ lexical) pour le vocabulaire. La règle d'or du baccalauréat : \textbf{toute réponse doit être ancrée dans le texte}.
\end{aretenir}

\coursec{Lexique thématique : l'environnement et le climat}

Dans la section précédente, nous avons lu et compris un article sur le changement climatique : vous avez répondu aux questions de compréhension et repéré les idées principales. Mais pour parler couramment de l'environnement — à l'oral comme à l'écrit, et notamment dans les compositions du baccalauréat — il ne suffit pas de \emph{comprendre} les mots : il faut savoir les \emph{employer} correctement. Cette section construit méthodiquement votre vocabulaire actif, organisé en quatre champs lexicaux : les problèmes, les causes, les conséquences et les solutions.

\courssub{Vue d'ensemble : la carte mentale du thème}

Avant d'entrer dans le détail, observez cette carte mentale. Elle résume l'organisation logique du vocabulaire de cette leçon : tout discours sur l'environnement suit le schéma \cle{problem} $\rightarrow$ \cle{cause} $\rightarrow$ \cle{effect} $\rightarrow$ \cle{solution}.

\begin{popfigure}
\begin{tikzpicture}[mindmap, grow cyclic, every node/.style={concept, text=white, align=center},
root concept/.append style={concept color=popblue, font=\bfseries},
level 1/.append style={level distance=3.2cm, sibling angle=90},
level 2/.append style={level distance=2.2cm, sibling angle=45, concept color=popmuted, font=\small}]
\node[root concept]{CLIMATE CHANGE}
  child[concept color=poporange]{ node{Problems}
    child{ node{pollution} }
    child{ node{drought} }
    child{ node{floods} }
  }
  child[concept color=popgold]{ node{Causes}
    child{ node{fossil fuels} }
    child{ node{emissions} }
    child{ node{waste} }
  }
  child[concept color=poppink]{ node{Effects}
    child{ node{melting glaciers} }
    child{ node{extinction} }
    child{ node{heatwaves} }
  }
  child[concept color=popgreen]{ node{Solutions}
    child{ node{recycle} }
    child{ node{plant trees} }
    child{ node{solar power} }
  };
\end{tikzpicture}
\legende{Carte mentale du vocabulaire du changement climatique : quatre champs lexicaux autour d'un thème central.}
\end{popfigure}

Retenez cette structure : elle vous servira de plan pour toute production écrite ou orale sur l'environnement. Un bon paragraphe d'opinion commence par le problème, explique ses causes, décrit ses effets, puis propose des solutions.

\courssub{Champ 1 : les problèmes environnementaux}

\begin{definition}[Les problèmes — Problems]
Un \cle{environmental problem} est un phénomène qui dégrade la nature et menace la vie humaine ou animale. Les problèmes majeurs sont : \eng{pollution} (la pollution), \eng{deforestation} (la déforestation), \eng{global warming} (le réchauffement climatique), \eng{drought} (la sécheresse), \eng{flood} (l'inondation) et \eng{erosion} (l'érosion).
\end{definition}

Observez ces phrases, tirées de la presse anglophone :

\begin{exemplebox}[Les problèmes en contexte]
\eng{Air pollution kills thousands of people every year in big cities.} — La pollution de l'air tue des milliers de personnes chaque année dans les grandes villes.

\eng{Deforestation is destroying the Congo Basin rainforest.} — La déforestation détruit la forêt tropicale du bassin du Congo.

\eng{The Far North of Cameroon suffers from severe droughts.} — L'Extrême-Nord du Cameroun souffre de sécheresses sévères.

\eng{Floods destroy crops every year in Douala.} — Les inondations détruisent les récoltes chaque année à Douala.

\eng{Soil erosion makes farming difficult in the West Region.} — L'érosion des sols rend l'agriculture difficile dans la région de l'Ouest.
\end{exemplebox}

\begin{propriete}[Comptables ou incomptables ?]
Attention à la nature grammaticale de ces noms :
\begin{itemize}
\item \eng{pollution}, \eng{erosion}, \eng{global warming} sont \textbf{incomptables} : jamais de pluriel, jamais de \eng{a}. On dit \eng{Pollution is increasing} (et non \emph{pollutions}).
\item \eng{drought} et \eng{flood} sont \textbf{comptables} : \eng{a severe drought}, \eng{frequent floods}.
\end{itemize}
\end{propriete}

\courssub{Champ 2 : les causes}

\begin{definition}[Les causes — Causes]
Les principales causes des problèmes environnementaux sont : \eng{fossil fuels} (les combustibles fossiles : pétrole, charbon, gaz), \eng{emissions} (les émissions de gaz), \eng{waste} (les déchets), \eng{overpopulation} (la surpopulation) et \eng{burning} (les brûlis, la combustion).
\end{definition}

\begin{exemplebox}[Les causes en contexte]
\eng{Burning fossil fuels releases carbon dioxide into the atmosphere.} — Brûler des combustibles fossiles libère du dioxyde de carbone dans l'atmosphère.

\eng{Factories produce dangerous emissions.} — Les usines produisent des émissions dangereuses.

\eng{Plastic waste pollutes our rivers and oceans.} — Les déchets plastiques polluent nos rivières et nos océans.

\eng{Overpopulation increases the demand for wood and land.} — La surpopulation augmente la demande en bois et en terres.
\end{exemplebox}

Notez la structure causale typique de l'anglais : \eng{Burning fossil fuels releases CO2} — sujet en \cle{-ing} + verbe à la 3\textsuperscript{e} personne du singulier. Le gérondif sujet est toujours singulier : \eng{Burning \textbf{releases}}, jamais \emph{burning release}.

\courssub{Champ 3 : les conséquences}

\begin{definition}[Les conséquences — Effects / Consequences]
Les effets du réchauffement climatique sont : \eng{melting glaciers} (la fonte des glaciers), \eng{rising sea levels} (l'élévation du niveau des mers), \eng{extinction} (l'extinction des espèces), \eng{food shortage} (la pénurie alimentaire) et \eng{heatwave} (la canicule, la vague de chaleur).
\end{definition}

\begin{exemplebox}[Les conséquences en contexte]
\eng{Melting glaciers cause rising sea levels.} — La fonte des glaciers provoque l'élévation du niveau des mers.

\eng{Many animal species are threatened with extinction.} — De nombreuses espèces animales sont menacées d'extinction.

\eng{Droughts lead to food shortages in the Sahel.} — Les sécheresses conduisent à des pénuries alimentaires au Sahel.

\eng{Last year, Europe experienced its worst heatwave ever.} — L'année dernière, l'Europe a connu sa pire canicule jamais enregistrée.
\end{exemplebox}

\begin{attention}[Effect ou affect ?]
\eng{Effect} (avec un \textbf{e}) est un \textbf{nom} : \eng{the effects of global warming} (les effets du réchauffement). \eng{Affect} (avec un \textbf{a}) est un \textbf{verbe} : \eng{Global warming affects everyone} (le réchauffement affecte tout le monde). Ne les confondez jamais à l'écrit.
\end{attention}

\courssub{Champ 4 : les solutions}

\begin{definition}[Les solutions — Solutions]
Pour lutter contre les problèmes environnementaux, on peut : \eng{recycle} (recycler), \eng{plant trees} (planter des arbres), \eng{use renewable energy} (utiliser les énergies renouvelables, comme \eng{solar power}, l'énergie solaire), \eng{protect} (protéger), \eng{reduce} (réduire) et \eng{save} (économiser, sauvegarder).
\end{definition}

\begin{exemplebox}[Les solutions en contexte]
\eng{We must reduce carbon emissions.} — Nous devons réduire les émissions de carbone.

\eng{Planting trees helps fight global warming.} — Planter des arbres aide à lutter contre le réchauffement climatique.

\eng{Many schools in Buea now use solar power.} — De nombreuses écoles de Buea utilisent désormais l'énergie solaire.

\eng{Recycling plastic bottles reduces waste.} — Recycler les bouteilles en plastique réduit les déchets.

\eng{Everyone should save water during the dry season.} — Chacun devrait économiser l'eau pendant la saison sèche.
\end{exemplebox}

\courssub{Tableau récapitulatif du vocabulaire}

\begin{center}
\begin{tabularx}{\linewidth}{|X|X|X|}\hline
\rowcolor{popblueL} English & Français & Remarque \\ \hline
pollution & la pollution & incomptable \\ \hline
deforestation & la déforestation & incomptable \\ \hline
drought & la sécheresse & comptable : \eng{a drought} \\ \hline
flood & l'inondation & comptable : \eng{floods} \\ \hline
greenhouse effect & l'effet de serre & \eng{the greenhouse effect} \\ \hline
carbon dioxide (CO2) & le dioxyde de carbone & incomptable \\ \hline
emissions & les émissions & souvent au pluriel \\ \hline
waste & les déchets & incomptable \\ \hline
extinction & l'extinction & \eng{threatened with extinction} \\ \hline
food shortage & la pénurie alimentaire & comptable \\ \hline
renewable energy & l'énergie renouvelable & incomptable \\ \hline
solar power & l'énergie solaire & incomptable \\ \hline
to recycle & recycler & verbe régulier \\ \hline
to plant trees & planter des arbres & \eng{trees} au pluriel \\ \hline
to raise awareness & sensibiliser & expression figée \\ \hline
\end{tabularx}
\end{center}

\courssub{Collocations fréquentes}

En anglais, certains verbes s'associent naturellement à certains noms : ce sont les \cle{collocations}. Les apprendre par blocs, et non mot à mot, est le secret d'un anglais naturel.

\begin{aretenir}[Collocations à mémoriser]
\begin{itemize}
\item \eng{to fight climate change} — lutter contre le changement climatique (et non \emph{to fight against climate change} : \eng{fight} est ici transitif direct).
\item \eng{to reduce emissions} — réduire les émissions.
\item \eng{to protect the environment} — protéger l'environnement.
\item \eng{to raise awareness} — sensibiliser (littéralement « élever la conscience ») : \eng{The club organised a campaign to raise awareness about plastic waste.} — Le club a organisé une campagne pour sensibiliser aux déchets plastiques.
\item \eng{to save energy / water} — économiser l'énergie / l'eau.
\item \eng{to cause damage to} — causer des dégâts à : \eng{Floods cause serious damage to farms.} — Les inondations causent de graves dégâts aux fermes.
\end{itemize}
\end{aretenir}

\courssub{Les familles de mots : construire son vocabulaire}

Un même radical donne souvent un verbe, un nom et un adjectif. Maîtriser ces familles multiplie votre vocabulaire sans effort.

\begin{center}
\begin{tabular}{|l|l|l|}\hline
\rowcolor{popgreenL} Verbe & Nom & Adjectif \\ \hline
to pollute (polluer) & pollution (la pollution) & polluted (pollué) \\ \hline
to warm (réchauffer) & warming (le réchauffement) & warm (chaud) \\ \hline
to protect (protéger) & protection (la protection) & protective (protecteur) \\ \hline
to reduce (réduire) & reduction (la réduction) & reduced (réduit) \\ \hline
\end{tabular}
\end{center}

\begin{exemplebox}[Une famille, trois phrases]
\eng{Factories pollute the river.} — Les usines polluent la rivière. (verbe)

\eng{Water pollution is a serious problem.} — La pollution de l'eau est un problème grave. (nom)

\eng{People cannot drink polluted water.} — On ne peut pas boire de l'eau polluée. (adjectif)
\end{exemplebox}

Notez que \eng{warmth} (la chaleur, la tiédeur) est le nom concret de \eng{warm}, tandis que \eng{warming} désigne le processus : \eng{global warming} = le réchauffement \emph{en cours}.

\courssub{Texte d'application : le vocabulaire en action}

Lisez ce court article. Soulignez mentalement tous les mots des quatre champs lexicaux étudiés.

\begin{textebox}[A Greener Future for Our Cities — original article]
Every morning, the people of Douala wake up to the same sight: plastic waste floating in the drains after the night's rain. Floods are becoming more frequent, and each rainy season brings new damage to homes, roads and crops. Scientists agree on the causes. The burning of fossil fuels in factories and old vehicles releases carbon dioxide and other emissions into the atmosphere. These gases trap the sun's heat — this is the famous greenhouse effect — and the planet slowly becomes warmer.

The effects are already visible. In the Far North, long droughts dry up the fields and cause food shortages. Along the coast, rising sea levels threaten fishing villages. Everywhere, deforestation continues: trees are cut for firewood and farmland, and the soil, no longer protected, suffers from erosion. Many plant and animal species are now threatened with extinction.

But there is hope. Across Cameroon, young people are taking action. School clubs plant trees every year and organise campaigns to raise awareness about waste. Some towns now recycle plastic bottles into paving stones. Solar power is bringing clean, renewable energy to villages far from the national grid. “We cannot wait for others to act,” says Amina, a student in Maroua. “If everyone reduces waste, saves water and plants just one tree, we can fight climate change together.”

The message is clear: protecting the environment is not the job of governments alone. It begins with each of us, today.
\end{textebox}

\textbf{Glossaire} : drains (n., les caniveaux) ; to trap (v., piéger, retenir) ; firewood (n., le bois de chauffe) ; grid (n., le réseau électrique) ; paving stones (n., les pavés).

\textbf{Questions rapides} : (1) Name two causes of global warming mentioned in the text. (2) What effects do droughts have in the Far North? (3) What solutions are young Cameroonians using?

\begin{corrige}[Questions rapides]
(1) The burning of fossil fuels and deforestation. (2) Droughts dry up the fields and cause food shortages. (3) They plant trees, organise awareness campaigns, recycle plastic bottles and use solar power.
\end{corrige}

\courssub{Exercice résolu et pièges à éviter}

\begin{exoresolu}[Complétez avec le mot de la bonne famille]
Énoncé — Complete each sentence with the correct form of the word in brackets: (1) Factories \dots\ the air in big cities. (pollute) (2) Air \dots\ causes many diseases. (pollute) (3) The \dots\ of the planet is a danger for everyone. (warm)
\tcblower
Solution détaillée — (1) Il faut un \textbf{verbe} conjugué, 3\textsuperscript{e} personne du pluriel : \eng{Factories pollute the air in big cities.} (Les usines polluent l'air des grandes villes.) (2) Il faut un \textbf{nom} sujet : \eng{Air pollution causes many diseases.} (La pollution de l'air provoque de nombreuses maladies.) (3) Il faut le nom du \textbf{processus} : \eng{The warming of the planet is a danger for everyone.} (Le réchauffement de la planète est un danger pour tous.) Astuce : demandez-vous toujours si la phrase attend un verbe, un nom ou un adjectif avant de choisir la forme.
\end{exoresolu}

\begin{attention}[Faux amis et orthographe : les pièges des francophones]
\begin{itemize}
\item \textbf{Faux ami} : le français \emph{sensible} se dit \eng{sensitive} en anglais (\eng{She is very sensitive} — elle est très sensible). L'anglais \eng{sensible} signifie \emph{raisonnable} : \eng{It is sensible to save water} — il est raisonnable d'économiser l'eau.
\item \textbf{Orthographe} : \eng{environment} s'écrit avec un seul \textbf{n} avant le \textbf{m} : en-vi-ron-ment. N'écrivez jamais \emph{enviroment}. Prononcez « in-vaï-ron-ment ».
\item \textbf{Préposition} : on dit \eng{to protect the environment \textbf{from} pollution} (protéger l'environnement \emph{de} la pollution), et \eng{to suffer \textbf{from} drought} (souffrir \emph{de} la sécheresse).
\item \textbf{Pas de pluriel} : \eng{waste} et \eng{pollution} sont incomptables ; ne dites jamais \emph{wastes} ou \emph{pollutions}.
\end{itemize}
\end{attention}

\begin{savaistu}[Pourquoi « greenhouse effect » ?]
L'expression \eng{greenhouse effect} (effet de serre) vient d'une comparaison toute simple : dans une \eng{greenhouse} — une serre de jardin, littéralement une « maison verte » —, le verre laisse entrer la lumière du soleil mais retient la chaleur à l'intérieur, ce qui permet aux plantes de pousser même en hiver. Les gaz comme le dioxyde de carbone jouent le rôle de cette vitre autour de la Terre : ils laissent passer les rayons du soleil mais bloquent une partie de la chaleur. Un peu d'effet de serre est naturel et indispensable à la vie ; c'est son \emph{emballement}, dû aux émissions humaines, qui provoque le réchauffement climatique.
\end{savaistu}

\courssub{Méthode de lecture : skimming, scanning et déduction}

\begin{methode}[Trois stratégies pour lire efficacement]
\textbf{Étape 1 — Skimming (lecture rapide pour l'idée générale).} Lisez le titre, les sous-titres, la première et la dernière phrase de chaque paragraphe. Ne vous arrêtez pas sur les mots inconnus. Objectif : identifier le sujet, la thèse de l'auteur et la structure du texte.
\textbf{Étape 2 — Scanning (recherche ciblée).} Repérez les mots-clés de la question (noms propres, chiffres, dates, termes techniques). Balayez le texte du regard pour les localiser. Lisez seulement la phrase concernée.
\textbf{Étape 3 — Déduire le sens (inférence lexicale).} Face à un mot inconnu : (a) Regardez la structure du mot (préfixe, racine, suffixe) : \eng{deforestation} = \eng{de-} + \eng{forest} + \eng{-ation}. (b) Exploitez le contexte : définition (\eng{is a process where...}), exemple (\eng{such as...}), contraste (\eng{but, however}), cause/effet (\eng{because, leads to}). (c) Vérifiez si le mot ressemble à un mot français (cognat) : \eng{pollution}, \eng{environment}, \eng{catastrophe}.
\end{methode}

\courssub{Production guidée : résumé et opinion}

\begin{methode}[Rédiger un résumé structuré + opinion personnelle]
\textbf{Étape 1 — Identifiez les idées forces.} Relisez le texte paragraphe par paragraphe. Notez en marge l'idée principale de chacun en 1 phrase (en anglais).
\textbf{Étape 2 — Assemblez le résumé.} Reliez vos notes avec des connecteurs logiques : \eng{First / Firstly}, \eng{Moreover / Furthermore}, \eng{Consequently / As a result}, \eng{However / On the other hand}, \eng{Finally / In conclusion}. Utilisez le \textbf{present simple} pour les faits généraux et le \textbf{present perfect} pour des situations actuelles résultant du passé (\eng{has become}, \eng{have increased}). Longueur visée : 40–60 mots.
\textbf{Étape 3 — Exprimez votre opinion.} Commencez par : \eng{In my opinion,} / \eng{I strongly believe that} / \eng{It seems to me that}. Appuyez-vous sur le vocabulaire de la leçon (\eng{We must reduce emissions}, \eng{Everyone should plant trees}). Terminez par un appel à l'action : \eng{It is high time we acted.}
\end{methode}

\begin{experience}[Atelier « Résumé \& Opinion » — travail en binôme]
\emph{Consigne :} À partir de l'article « A Greener Future for Our Cities », rédigez ensemble un résumé de 50 mots maximum, puis une opinion personnelle de 30 mots. Utilisez au moins 5 mots du tableau récapitulatif et 2 collocations. Lisez votre production à la classe ; les camarades notent le vocabulaire utilisé.
\end{experience}

\courssub{Exercice d'entraînement : À vous de jouer}

\begin{exercice}[Entraînement complet]
(1) Traduisez en anglais : « Nous devons protéger l'environnement et réduire les déchets. »
(2) Trouvez le nom correspondant : to erode, to extinct (attention au piège !), to emit.
(3) Construisez une phrase complète avec la collocation \eng{to raise awareness}.
(4) Classez ces mots dans les quatre champs lexicaux (Problèmes, Causes, Conséquences, Solutions) : \eng{heatwave, solar power, emissions, flood, to recycle, extinction}.
\end{exercice}

\begin{corrige}[Exercice d'entraînement]
(1) \eng{We must protect the environment and reduce waste.} (Attention : \eng{waste} incomptable, pas de \eng{the} devant \eng{environment} dans un énoncé général).
(2) \eng{to erode} $\rightarrow$ \eng{erosion} ; \eng{to emit} $\rightarrow$ \eng{emission} ; \textbf{piège} : \eng{extinct} est un \textbf{adjectif} (pas un verbe), le nom est \eng{extinction}. Le verbe associé est \eng{to become extinct} ou \eng{to die out}.
(3) Exemple : \eng{Our school organised a march to raise awareness about plastic pollution.}
(4) \textbf{Problèmes} : \eng{flood, heatwave} ; \textbf{Causes} : \eng{emissions} ; \textbf{Conséquences} : \eng{extinction} ; \textbf{Solutions} : \eng{solar power, to recycle}.
\end{corrige}

Vous disposez maintenant d'un vocabulaire structuré et actif sur l'environnement, ainsi que des stratégies de lecture et d'écriture pour réussir vos épreuves du baccalauréat. La prochaine leçon portera sur la grammaire du discours rapporté (\emph{reported speech}), indispensable pour relater les propos des experts climatiques dans vos compositions.

\coursec{Méthode de lecture : skimming, scanning, déduction}

Après avoir acquis le lexique essentiel de l'environnement et du climat dans la section précédente, nous abordons maintenant les stratégies de lecture qui vous permettront d'aborder n'importe quel texte — même long et dense — avec efficacité et sérénité, notamment lors de l'épreuve du Baccalauréat. Ces trois techniques (skimming, scanning, déduction) ne sont pas des « astuces » : ce sont des compétences cognitives que les lecteurs experts mobilisent inconsciemment et que vous allez apprendre à maîtriser consciemment.

\courssub{Le skimming : saisir l'idée générale}

\begin{definition}[Skimming]
Le \textbf{skimming} (ou \emph{lecture diagonale}) consiste à lire rapidement un texte pour en saisir l'idée générale (\eng{the gist}, \eng{the main idea}) sans s'attarder sur les détails. On ne lit pas chaque mot : on repère la structure du texte (titre, sous-titres, première phrase de chaque paragraphe — la \emph{topic sentence} — et mots-clés en gras ou en italique).
\end{definition}

\eng{Skimming helps you decide if a text is relevant to your purpose.} « Le skimming vous aide à décider si un texte est pertinent pour votre objectif. »

\begin{propriete}[Comment pratiquer le skimming]
\begin{enumerate}
\item Lisez le \textbf{titre} et les \textbf{sous-titres} : ils annoncent le thème et la structure.
\item Lisez la \textbf{première phrase} de chaque paragraphe (\eng{topic sentence}) : elle résume l'idée principale du paragraphe.
\item Repérez les \textbf{mots répétés} et les \textbf{mots en gras/italique} : ce sont les concepts clés.
\item Ignorez les exemples détaillés, les chiffres précis, les adjectifs descriptifs.
\item Formulez en une phrase l'idée générale du texte (\eng{the main idea}).
\end{enumerate}
\end{propriete}

\begin{exemplebox}[Skimming en action]
Sur un texte intitulé \eng{« The Hidden Costs of Fast Fashion »}, les \eng{topic sentences} sont :
\begin{itemize}
\item \eng{The fashion industry is the second-largest polluter in the world.}
\item \eng{Water consumption and chemical dyeing devastate local ecosystems.}
\item \eng{Textile waste clogs landfills for centuries.}
\item \eng{Consumers can drive change through conscious choices.}
\end{itemize}
\eng{Main idea :} \eng{The text exposes the environmental damage caused by fast fashion and calls for responsible consumption.}
\end{exemplebox}

\courssub{Le scanning : localiser une information précise}

\begin{definition}[Scanning]
Le \textbf{scanning} (ou \emph{lecture de repérage}) consiste à déplacer rapidement les yeux sur le texte pour trouver une \textbf{information précise} : un chiffre, une date, un nom propre, un terme technique, une définition. On ne lit pas le texte : on le « scanne » comme un code-barres, à la recherche de \textbf{mots-clés} identifiés dans la question.
\end{definition}

\eng{Scanning is useful when you answer specific comprehension questions.} « Le scanning est utile quand vous répondez à des questions de compréhension précises. »

\begin{propriete}[Comment pratiquer le scanning]
\begin{enumerate}
\item Lisez la \textbf{question} et soulignez le ou les \textbf{mots-clés} (noms, verbes, chiffres, mots interrogatifs \eng{who, when, where, how much}).
\item Cherchez \textbf{visuellement} ces mots-clés (ou leurs synonymes) dans le texte.
\item Une fois localisés, lisez \textbf{attentivement la phrase} qui les contient et les phrases voisines.
\item Extrayez la réponse exacte.
\end{enumerate}
\end{propriete}

\begin{center}
\begin{tabularx}{\linewidth}{|X|X|X|}
\hline
\rowcolor{popblueL} \textbf{Question type} & \textbf{Mots-clés à scanner} & \textbf{Information visée} \\
\hline
\eng{When was the Paris Agreement signed?} & \eng{Paris Agreement, signed, date} & Date (chiffre + mois + année) \\
\hline
\eng{What percentage of CO2 comes from transport?} & \eng{percentage, CO2, transport} & Chiffre + \% \\
\hline
\eng{Who warned about rising sea levels?} & \eng{warned, rising sea levels} & Nom de personne / organisme \\
\hline
\eng{Which city hosted the COP27?} & \eng{city, hosted, COP27} & Nom de ville \\
\hline
\end{tabularx}
\end{center}

\courssub{La déduction du sens : devenir autonome face au vocabulaire inconnu}

\begin{propriete}[Trois leviers pour déduire le sens d'un mot inconnu]
\begin{enumerate}
\item \textbf{Le contexte immédiat (co-texte)} : les mots qui entourent le mot inconnu (synonymes, antonymes, énumérations, relations de cause-effet).
\item \textbf{La morphologie (famille de mots, préfixes, suffixes)} : décomposer le mot en unités de sens connues.
\item \textbf{La transparence morphologique (cognats, racines latines/grecques)} : beaucoup de mots anglais académiques partagent des racines avec le français.
\end{enumerate}
\end{propriete}

\begin{exemplebox}[Déduction pas à pas : \eng{heatwave}]
\eng{Context :} \eng{« Extreme weather events such as floods, droughts and \textbf{heatwaves} are becoming more frequent. »}
\begin{enumerate}
\item \textbf{Contexte :} \eng{floods} (inondations), \eng{droughts} (sécheresses) et \eng{heatwaves} forment une énumération d'\eng{extreme weather events} (événements météorologiques extrêmes). \eng{Heatwave} est donc un phénomène météo extrême.
\item \textbf{Morphologie :} \eng{heat} (chaleur) + \eng{wave} (vague) $\rightarrow$ composé nominal.
\item \textbf{Déduction :} \eng{heatwave} = « vague de chaleur » (période de chaleur intense prolongée).
\end{enumerate}
\end{exemplebox}

\begin{exemplebox}[Autres déductions morphologiques utiles pour l'environnement]
\begin{center}
\begin{tabular}{|l|l|l|l|}
\hline
\rowcolor{popgreenL} \textbf{Mot anglais} & \textbf{Décomposition} & \textbf{Sens déduit} & \textbf{Français} \\
\hline
\eng{renewable} & \eng{re-} (de nouveau) + \eng{new} + \eng{-able} & qui peut être renouvelé & renouvelable \\
\hline
\eng{unsustainable} & \eng{un-} (non) + \eng{sustain} (soutenir) + \eng{-able} & qui ne peut pas être maintenu & non durable \\
\hline
\eng{deforestation} & \eng{de-} (privation) + \eng{forest} + \eng{-ation} & action de priver de forêt & déforestation \\
\hline
\eng{overconsumption} & \eng{over-} (excès) + \eng{consumption} & consommation excessive & surconsommation \\
\hline
\eng{biodiversity} & \eng{bio-} (vie) + \eng{diversity} & diversité du vivant & biodiversité \\
\hline
\eng{carbon-neutral} & \eng{carbon} + \eng{neutral} & neutre en carbone & neutre en carbone \\
\hline
\end{tabular}
\end{center}
\end{exemplebox}

\courssub{Exemple traité pas à pas sur un texte support}

Voici un texte original d'environ 250 mots, typique de ce que vous pouvez rencontrer au Bac. Appliquez-y la méthode complète.

\begin{textebox}[Texte support : Cameroon's Green Cities Initiative]
Cameroon's Green Cities Initiative, launched in 2021 by the Ministry of Housing and Urban Development (MINHDU), aims to transform major urban centres like Yaoundé, Douala and Buea into models of climate resilience. Rapid urbanisation has led to chaotic sprawl, inadequate waste management and the destruction of wetlands that once regulated flooding. The programme rests on three pillars: sustainable mobility, urban reforestation and circular waste economy.

First, sustainable mobility seeks to reduce the carbon footprint of daily commutes. In Douala, the “Bus Rapid Transit” (BRT) project will introduce dedicated lanes for high-capacity electric buses, cutting travel time by 30\% and CO2 emissions by an estimated 15,000 tonnes per year. Cycle lanes and pedestrian zones are also planned in city centres.

Second, urban reforestation targets the “heat island” effect. Over 500,000 native trees — including iroko, mahogany and moabi — will be planted by 2030 along streets, in school yards and in newly created urban parks. These species absorb pollutants, lower ambient temperature by up to 3°C and restore biodiversity corridors.

Third, the circular waste economy turns a liability into a resource. Pilot sorting centres in Yaoundé convert organic waste into compost for peri-urban agriculture and recycle plastic into paving stones. Informal waste pickers are integrated into cooperatives, securing livelihoods while boosting recycling rates from 12\% to a target of 45\% by 2027.

Funding comes from the national budget, the African Development Bank and private green bonds. Success depends on citizen participation: the “Eco-Citizen” mobile app allows residents to report illegal dumping and track tree-planting progress. If targets are met, Cameroon's cities could become a blueprint for sustainable urbanisation in Central Africa.
\end{textebox}

\begin{definition}[Glossaire]
\begin{itemize}
\item \eng{sprawl} (n.) : étalement urbain anarchique
\item \eng{wetlands} (n. pl.) : zones humides
\item \eng{commutes} (n. pl.) : trajets domicile-travail
\item \eng{heat island} (n.) : îlot de chaleur urbain
\item \eng{peri-urban} (adj.) : périurbain (autour de la ville)
\item \eng{livelihoods} (n. pl.) : moyens de subsistance
\item \eng{blueprint} (n.) : modèle, plan de référence
\end{itemize}
\end{definition}

\begin{methode}[Démonstration : Skimming $\rightarrow$ Questions $\rightarrow$ Scanning $\rightarrow$ Vérification]
\textbf{Étape 1 — Skimming (30 s)} \\
Titre : \eng{Cameroon's Green Cities Initiative}. \\
Sous-structure : trois piliers annoncés au paragraphe 1. \\
\eng{Topic sentences} : \\
$\bullet$ Paragr. 2 : \eng{First, sustainable mobility seeks to reduce the carbon footprint...} \\
$\bullet$ Paragr. 3 : \eng{Second, urban reforestation targets the ``heat island'' effect.} \\
$\bullet$ Paragr. 4 : \eng{Third, the circular waste economy turns a liability into a resource.} \\
$\bullet$ Paragr. 5 : \eng{Funding comes from... Success depends on citizen participation.} \\
\eng{Main idea} : \eng{The article presents a three-pillar government programme to make Cameroonian cities greener and more resilient.}

\textbf{Étape 2 — Lecture des questions (identification des mots-clés)} \\
Q1 : \eng{What are the three pillars of the initiative?} $\rightarrow$ mots-clés : \eng{three pillars} \\
Q2 : \eng{How much will CO2 emissions drop in Douala thanks to the BRT?} $\rightarrow$ mots-clés : \eng{CO2 emissions, drop, Douala, BRT} \\
Q3 : \eng{Which tree species are mentioned?} $\rightarrow$ mots-clés : \eng{tree species, mentioned} \\
Q4 : \eng{What is the target recycling rate for 2027?} $\rightarrow$ mots-clés : \eng{target, recycling rate, 2027} \\
Q5 : \eng{How can citizens participate?} $\rightarrow$ mots-clés : \eng{citizens, participate}

\textbf{Étape 3 — Scanning (1 min)} \\
Q1 : Paragr. 1, phrase 3 : \eng{sustainable mobility, urban reforestation and circular waste economy}. \\
Q2 : Paragr. 2 : \eng{cutting travel time by 30\% and CO2 emissions by an estimated 15,000 tonnes per year}. \\
Q3 : Paragr. 3 : \eng{iroko, mahogany and moabi}. \\
Q4 : Paragr. 4 : \eng{boosting recycling rates from 12\% to a target of 45\% by 2027}. \\
Q5 : Paragr. 5 : \eng{the ``Eco-Citizen'' mobile app allows residents to report illegal dumping and track tree-planting progress}.

\textbf{Étape 4 — Vérification et justification} : Relire les phrases citées pour confirmer l'exactitude des réponses.
\end{methode}

\courssub{Méthode pour gagner du temps au Bac}

\begin{methode}[Stratégie « Zéro traduction intégrale » — Spéciale Baccalauréat]
\begin{enumerate}
\item \textbf{Ne traduisez jamais le texte en entier.} C'est une perte de temps massive (15–20 min) et source d'erreurs.
\item \textbf{Lisez d'abord les questions} (1 minute). Soulignez les mots-clés de chaque question.
\item \textbf{Faites un skimming ultra-rapide} (30 secondes) : titre + première phrase de chaque paragraphe. Cela vous donne la « carte » du texte.
\item \textbf{Scannez pour chaque question} : cherchez les mots-clés (ou synonymes) repérés à l'étape 2.
\item \textbf{Rédigez la réponse en anglais} en reprenant les mots du texte (citation) ou en reformulant simplement.
\item \textbf{Vérifiez} : la réponse répond-elle exactement à la question ? L'orthographe des mots cités est-elle correcte ?
\end{enumerate}
\end{methode}

\begin{attention}[Pièges fréquents des francophones]
\begin{itemize}
\item \textbf{Traduire mot à mot} : vous perdez du temps et vous ratez la cohérence globale. \eng{Skimming} et \eng{scanning} sont des stratégies de \emph{repérage}, pas de traduction.
\item \textbf{S'arrêter sur chaque mot inconnu} : 80\% des mots difficiles sont soit dans le glossaire, soit déductibles (contexte + morphologie), soit \textbf{inutiles pour répondre aux questions}. Passez outre.
\item \textbf{Confondre \eng{skimming} et \eng{scanning}} : \eng{skimming} = idée générale (haut du texte) ; \eng{scanning} = info précise (n'importe où dans le texte).
\item \textbf{Répondre en français} : au Bac, les réponses aux questions de compréhension doivent être en \textbf{anglais}, sauf consigne contraire explicite.
\item \textbf{Copier des phrases entières} sans les adapter à la question : citez le segment pertinent seulement.
\end{itemize}
\end{attention}

\courssub{Entraînement chronométré}

\begin{experience}[Atelier « Chrono-Reading » — 10 minutes]
\textbf{Consigne :} Utilisez le texte \eng{« Cameroon's Green Cities Initiative »} ci-dessus.
\begin{enumerate}
\item \textbf{Skimming chronométré (30 s) :} Fermez les yeux, ouvrez-les, lisez titre + \eng{topic sentences} uniquement. Écrivez l'idée générale en une phrase anglaise.
\item \textbf{Scanning chronométré (1 min) :} Répondez aux 5 questions Q1–Q5 de la démonstration sans relire le texte en entier. Chronométrez-vous.
\item \textbf{Auto-correction :} Comparez vos réponses à la démonstration. Notez le temps total et le nombre de réponses exactes.
\item \textbf{Défi :} Refaites l'exercice dans 48 h pour battre votre temps tout en gardant 5/5.
\end{enumerate}
\textbf{Objectif :} Skimming $\le$ 30 s, Scanning 5 questions $\le$ 1 min, Score 5/5.
\end{experience}

\courssub{Synthèse visuelle : organigramme de la tâche de lecture}

\begin{popfigure}
\begin{tikzpicture}[
  node distance=1.2cm and 1.5cm,
  box/.style={rectangle, rounded corners, draw=popblue, thick, fill=popblueL, text width=3.2cm, align=center, font=\small},
  arrow/.style={-Stealth, thick, draw=popdark}
]
\node[box, align=center] (start){READING TASK\\Text + Questions};
\node[box, below=of start, align=center] (skim){SKIM\\General idea\\(30 s)};
\node[box, below=of skim, align=center] (readq){READ QUESTIONS\\Underline keywords\\(1 min)};
\node[box, below=of readq, align=center] (scan){SCAN\\Locate answers\\(1 min)};
\node[box, below=of scan, align=center] (check){CHECK \& JUSTIFY\\Quote / Paraphrase\\(30 s)};

\draw[arrow] (start) -- (skim);
\draw[arrow] (skim) -- (readq);
\draw[arrow] (readq) -- (scan);
\draw[arrow] (scan) -- (check);

% Feedback loop
\draw[arrow, poporange] (check.east) -- ++(2,0) |- (readq.east) node[pos=0.25, right, font=\tiny, text=poporange] {If unsure};
\end{tikzpicture}
\legende{Organigramme de la stratégie de lecture efficace : Skimming $\rightarrow$ Questions $\rightarrow$ Scanning $\rightarrow$ Vérification. La boucle de retour (orange) invite à relire les questions si une réponse reste incertaine.}
\end{popfigure}

\courssub{Exercice résolu}

\begin{exoresolu}[Application sur un court extrait]
\textbf{Extrait :} \eng{« The Congo Basin rainforest, often called the ``second lung of the Earth'', stores 30 billion tonnes of carbon. However, illegal logging and agricultural expansion threaten 1.5 million hectares annually. The CAFI initiative (Central African Forest Initiative) has pledged \$500 million to protect these ecosystems by 2030. »}

\textbf{Questions :}
\begin{enumerate}
\item What nickname is given to the Congo Basin rainforest?
\item How much carbon does it store?
\item What are the two main threats mentioned?
\item How much money has CAFI pledged?
\item By when does CAFI aim to protect the forest?
\end{enumerate}

\tcblower

\textbf{Solution détaillée :}

\textbf{Skimming (5 s) :} Sujet = forêt du bassin du Congo / menaces / initiative de protection.

\textbf{Scanning par question :}
\begin{enumerate}
\item Mots-clés : \eng{nickname, called} $\rightarrow$ \eng{``second lung of the Earth''}. \textbf{Réponse :} \eng{It is called the ``second lung of the Earth''.}
\item Mots-clés : \eng{stores, carbon} $\rightarrow$ \eng{30 billion tonnes}. \textbf{Réponse :} \eng{It stores 30 billion tonnes of carbon.}
\item Mots-clés : \eng{threats, threatened} $\rightarrow$ \eng{illegal logging and agricultural expansion}. \textbf{Réponse :} \eng{Illegal logging and agricultural expansion.}
\item Mots-clés : \eng{pledged, money, \$} $\rightarrow$ \eng{\$500 million}. \textbf{Réponse :} \eng{CAFI has pledged \$500 million.}
\item Mots-clés : \eng{by when, 2030} $\rightarrow$ \eng{by 2030}. \textbf{Réponse :} \eng{By 2030.}
\end{enumerate}

\textbf{Vérification :} Toutes les réponses sont des citations directes ou quasi-directes, précises et en anglais correct.
\end{exoresolu}

\begin{aretenir}[Checklist « Lecture efficace » — À coller dans votre cahier]
\begin{itemize}
\item[$\square$] Je lis le titre et les \eng{topic sentences} (skimming) avant de lire les questions.
\item[$\square$] Je souligne les mots-clés des questions.
\item[$\square$] Je scanne le texte pour localiser ces mots-clés (ou synonymes).
\item[$\square$] Je déduis le sens des mots inconnus par le contexte + la morphologie.
\item[$\square$] Je réponds en anglais, en citant le segment pertinent.
\item[$\square$] Je ne traduis jamais le texte en entier.
\end{itemize}
\end{aretenir}

\coursec{Production guidée : résumé et opinion}

Dans la section précédente, tu as appris à lire efficacement grâce à trois stratégies : le \cle{skimming} (lecture rapide pour l'idée générale), le \cle{scanning} (recherche d'informations précises) et la déduction du sens des mots inconnus à partir du contexte. Mais lire, c'est aussi savoir \emph{restituer} : à l'examen comme dans la vie, on te demandera souvent de résumer un texte et de donner ton avis. Cette dernière étape de la leçon te donne tous les outils pour écrire, en cinq à huit lignes, un paragraphe clair qui résume un texte sur l'environnement et exprime ton opinion personnelle.

\courssub{Observer avant d'écrire : un modèle commenté}

Avant toute règle, lis attentivement ce court paragraphe écrit par une élève de Terminale après sa lecture d'un article sur le réchauffement climatique. Observe l'ordre des idées et les mots soulignés par ta propre attention : ce sont les \cle{connecteurs}.

\begin{exemplebox}[Un modèle de résumé + opinion]
\eng{The text deals with global warming. First, the writer explains that human activities release greenhouse gases. Then, he shows the effects: melting glaciers and extreme weather. However, he believes solutions exist, such as renewable energy. In my opinion, every Cameroonian should plant trees because our forests are disappearing.}

« Le texte traite du réchauffement climatique. D'abord, l'auteur explique que les activités humaines rejettent des gaz à effet de serre. Ensuite, il montre les effets : la fonte des glaciers et les phénomènes météorologiques extrêmes. Cependant, il pense que des solutions existent, comme les énergies renouvelables. À mon avis, chaque Camerounais devrait planter des arbres parce que nos forêts disparaissent. »
\end{exemplebox}

Que remarques-tu ? Le paragraphe suit un chemin précis : \textbf{problème} (\eng{global warming}) $\rightarrow$ \textbf{causes} (\eng{human activities release greenhouse gases}) $\rightarrow$ \textbf{conséquences} (\eng{melting glaciers, extreme weather}) $\rightarrow$ \textbf{solutions} (\eng{renewable energy}) $\rightarrow$ \textbf{opinion personnelle} (\eng{In my opinion...}). C'est exactement l'architecture que tu vas reproduire.

\begin{propriete}[Les règles d'or du résumé]
Un \cle{résumé} (\eng{summary}) obéit à trois règles strictes :
\begin{enumerate}
\item Il ne contient \textbf{ni exemples personnels, ni citations longues} du texte : on ne recopie pas, on \emph{reformule} les idées principales avec ses propres mots.
\item Il est \textbf{court} : une phrase d'introduction et trois à quatre phrases pour le développement suffisent ; on ajoute ensuite une ou deux phrases d'opinion.
\item Il suit l'ordre logique \textbf{problème $\rightarrow$ causes $\rightarrow$ conséquences $\rightarrow$ solutions}, annoncé par des connecteurs.
\end{enumerate}
\end{propriete}

\courssub{Les connecteurs : la charpente de ton paragraphe}

Les connecteurs (\cle{linking words}) sont les panneaux indicateurs de ton texte. Sans eux, tes phrases se succèdent sans lien logique. Voici les six connecteurs essentiels de cette leçon, à connaître par cœur.

\begin{center}
\begin{tabularx}{\linewidth}{|l|X|X|}
\hline
\rowcolor{popblueL} \textbf{Connecteur} & \textbf{Fonction} & \textbf{Exemple} \\
\hline
\eng{first} & introduire la première idée & \eng{First, the writer presents the problem.} \\
\hline
\eng{then} & enchaîner l'idée suivante & \eng{Then, he explains the causes.} \\
\hline
\eng{moreover} & ajouter une information & \eng{Moreover, floods destroy crops every year.} \\
\hline
\eng{however} & exprimer un contraste & \eng{However, solutions exist.} \\
\hline
\eng{therefore} & exprimer une conséquence & \eng{Therefore, we must act now.} \\
\hline
\eng{in conclusion} & conclure & \eng{In conclusion, the article is a call to action.} \\
\hline
\end{tabularx}
\end{center}

\begin{attention}[Ponctuation des connecteurs]
En anglais, un connecteur placé en début de phrase est suivi d'une \textbf{virgule} : \eng{However, the situation is serious.} (« Cependant, la situation est grave. »). N'oublie jamais cette virgule : son absence est une faute classique dans les copies. De même, \eng{first} s'écrit sans majuscule sauf en tout début de phrase, et on évite de répéter \eng{then} à chaque phrase : varie avec \eng{moreover} ou \eng{also}.
\end{attention}

\courssub{Exprimer son opinion, son accord et son désaccord}

Après le résumé, tu donnes ton avis. L'anglais dispose d'expressions figées qu'il faut employer telles quelles, sans les traduire mot à mot depuis le français.

\begin{center}
\begin{tabularx}{\linewidth}{|X|X|}
\hline
\rowcolor{popgreenL} \textbf{English} & \textbf{Français} \\
\hline
\eng{I think that climate change is real.} & Je pense que le changement climatique est réel. \\
\hline
\eng{In my opinion, the government should invest in solar energy.} & À mon avis, le gouvernement devrait investir dans l'énergie solaire. \\
\hline
\eng{I believe that young people can change things.} & Je crois que les jeunes peuvent changer les choses. \\
\hline
\eng{To my mind, planting trees is the best solution.} & Selon moi, planter des arbres est la meilleure solution. \\
\hline
\eng{I agree that every small action counts.} & Je suis d'accord que chaque petit geste compte. \\
\hline
\eng{I disagree because the problem is too big for individuals alone.} & Je ne suis pas d'accord parce que le problème est trop vaste pour les individus seuls. \\
\hline
\end{tabularx}
\end{center}

\begin{propriete}[Forme et emploi des expressions d'opinion]
\begin{itemize}
\item \eng{I think that} + proposition complète (sujet + verbe). Le mot \eng{that} est facultatif : \eng{I think (that) we should recycle.}
\item \eng{In my opinion} et \eng{to my mind} se placent en tête de phrase, suivis d'une virgule.
\item Pour le désaccord : \eng{I disagree because...} — jamais \eng{I am disagree} : \eng{disagree} est un verbe ordinaire, sans auxiliaire \eng{be}.
\item Après \eng{should}, le verbe est toujours à la base verbale : \eng{the government should invest} (et non \eng{should to invest} ni \eng{should invests}).
\end{itemize}
\end{propriete}

\begin{attention}[Le piège du « Moi, je pense »]
Ne commence \textbf{jamais} une opinion par \eng{Me, I think...} : c'est un calque du français « Moi, je pense ». En anglais, le pronom sujet suffit : \eng{I think that...} De même, évite \eng{For me, ...} au sens de « à mon avis » ; dis \eng{In my opinion, ...} Autre erreur fréquente : \eng{I am agree with you} est incorrect — dis \eng{I agree with you} (« Je suis d'accord avec toi »), car \eng{agree} est un verbe d'action, pas un adjectif.
\end{attention}

\courssub{Le plan en escalier : visualiser avant de rédiger}

La figure suivante résume la progression de ton paragraphe : quatre marches pour le résumé, puis la marche de l'opinion, justifiée par \eng{because}.

\begin{popfigure}
\begin{center}
\begin{tikzpicture}[font=\small]
\node[draw=popblue, fill=popblueL, rounded corners, minimum width=2.6cm, minimum height=0.85cm, align=center] (p) at (0,0) {\textbf{Problem}\\ \eng{deals with...}};
\node[draw=poporange, fill=poporangeL, rounded corners, minimum width=2.6cm, minimum height=0.85cm, align=center] (c) at (3.2,0.9) {\textbf{Causes}\\ \eng{first, ...}};
\node[draw=poppurple, fill=poppurpleL, rounded corners, minimum width=2.6cm, minimum height=0.85cm, align=center] (e) at (6.4,1.8) {\textbf{Effects}\\ \eng{then, moreover...}};
\node[draw=popgreen, fill=popgreenL, rounded corners, minimum width=2.6cm, minimum height=0.85cm, align=center] (s) at (9.6,2.7) {\textbf{Solutions}\\ \eng{however, ...}};
\node[draw=popgold, fill=popgoldL, rounded corners, minimum width=3.4cm, minimum height=0.95cm, align=center] (o) at (12.6,3.8) {\textbf{Opinion}\\ \eng{I think... because...}};
\draw[-{Stealth}, very thick, popdark] (p) -- (c);
\draw[-{Stealth}, very thick, popdark] (c) -- (e);
\draw[-{Stealth}, very thick, popdark] (e) -- (s);
\draw[-{Stealth}, very thick, popdark] (s) -- (o);
\node[align=center, text=popmuted] at (6.4,-0.9) {Summary (3--4 sentences) \quad $\rightarrow$ \quad Opinion (1--2 sentences)};
\end{tikzpicture}
\end{center}
\legende{L'escalier du paragraphe : problème, causes, effets, solutions, puis opinion justifiée.}
\end{popfigure}

\courssub{Texte d'entraînement : lire, résumer, réagir}

Pour t'entraîner, voici un court article original. Lis-le d'abord en \cle{skimming} (idée générale), puis en \cle{scanning} (causes, effets, solutions), avant de passer à la production.

\begin{textebox}[A river in danger: the Wouri under pressure]
The Wouri river, which flows through Douala, is one of the most important rivers in Cameroon. Fishermen depend on it, traders cross it every day, and thousands of families use its water. However, the river is now in serious danger.

First, factories near the port release chemical waste directly into the water. Moreover, plastic bags and bottles, thrown away in the streets of Douala, are carried into the river by the rain. As a result, the water is becoming more polluted every year.

The consequences are already visible. Fish are disappearing, and many fishermen earn less money than before. In addition, people who use the river water sometimes fall ill. During the rainy season, plastic waste blocks the drains, and this causes floods in several neighbourhoods.

Fortunately, solutions exist. The city council has started a cleaning campaign, and some schools organise “green days” when pupils collect plastic around the river. Environmental associations also ask factories to treat their waste before releasing it.

In conclusion, the Wouri can be saved, but only if everybody takes part: the authorities, the factories and the citizens of Douala. As one fisherman says, “The river feeds us; now it is our turn to feed the river.”
\end{textebox}

\textbf{Glossaire.} \emph{to flow} (v.) : couler ; \emph{waste} (n.) : déchets ; \emph{to throw away} (v.) : jeter ; \emph{drain} (n.) : égout, canal d'évacuation ; \emph{flood} (n.) : inondation ; \emph{city council} (n.) : conseil municipal ; \emph{to treat} (v.) : traiter ; \emph{to take part} (v.) : participer.

\textbf{Questions de vérification rapide.} (1) \eng{What are the two main causes of pollution in the Wouri?} (2) \eng{What happens during the rainy season?} (3) \eng{Who organises “green days”?}

\begin{exoresolu}[Rédiger le résumé + opinion du texte sur le Wouri]
Énoncé : écris un paragraphe de 5 à 8 lignes qui résume l'article \emph{A river in danger} et donne ton opinion personnelle. Utilise au moins quatre connecteurs et une expression d'opinion.
\tcblower
Solution détaillée. On suit l'escalier : problème $\rightarrow$ causes $\rightarrow$ effets $\rightarrow$ solutions $\rightarrow$ opinion.

\eng{The text deals with the pollution of the Wouri river in Douala. First, the writer explains that factories release chemical waste into the water. Then, he shows that plastic waste from the streets makes the situation worse. Moreover, the effects are serious: fish are disappearing and floods hit the neighbourhoods. However, solutions exist, such as cleaning campaigns and “green days” in schools. In my opinion, every citizen of Douala should stop throwing plastic in the streets, because a clean river means a healthy city.}

Analyse : le résumé tient en quatre phrases (une d'introduction et trois de développement), sans recopier de phrase entière du texte ; les connecteurs \eng{first}, \eng{then}, \eng{moreover}, \eng{however} structurent le propos ; l'opinion est introduite par \eng{in my opinion} et justifiée par \eng{because}. Le paragraphe fait six lignes : la consigne est respectée.
\end{exoresolu}

\courssub{Grille d'auto-évaluation}

Après avoir rédigé ton propre paragraphe, évalue-toi honnêtement avec cette grille. Coche chaque critère : c'est exactement ce que ton correcteur regardera.

\begin{center}
\begin{tabularx}{\linewidth}{|X|c|c|c|}
\hline
\rowcolor{popblueL} \textbf{Critère} & \textbf{Oui} & \textbf{Un peu} & \textbf{Non} \\
\hline
\textbf{Contenu} : mon résumé suit l'ordre problème $\rightarrow$ causes $\rightarrow$ effets $\rightarrow$ solutions & & & \\
\hline
\textbf{Contenu} : je n'ai recopié aucune phrase du texte & & & \\
\hline
\textbf{Lexique} : j'ai utilisé au moins trois mots du thème (\eng{pollution, waste, flood, solution...}) & & & \\
\hline
\textbf{Connecteurs} : j'ai utilisé au moins quatre connecteurs variés & & & \\
\hline
\textbf{Opinion} : j'ai utilisé \eng{in my opinion} ou \eng{I think that} + une justification avec \eng{because} & & & \\
\hline
\textbf{Grammaire} : verbes conjugués, virgule après les connecteurs, pas de \eng{Me, I think} & & & \\
\hline
\end{tabularx}
\end{center}

\begin{methode}[Rédiger ton paragraphe en cinq étapes]
\begin{enumerate}
\item \textbf{Relis} le texte en skimming et note en un mot l'idée de chaque paragraphe.
\item \textbf{Classe} tes notes dans les quatre cases de l'escalier : problème, causes, effets, solutions.
\item \textbf{Rédige} le résumé en reformulant avec tes propres mots, une phrase par marche, avec un connecteur par phrase.
\item \textbf{Ajoute} ton opinion : une expression d'opinion + une justification introduite par \eng{because}.
\item \textbf{Relis-toi} avec la grille d'auto-évaluation et corrige avant de rendre.
\end{enumerate}
\end{methode}

\begin{exercice}[À toi de produire]
Choisis l'un des deux sujets suivants et rédige un paragraphe de 5 à 8 lignes (résumé + opinion) :
\begin{enumerate}
\item Résume l'article \emph{A river in danger} et explique si, selon toi, les écoles camerounaises devraient organiser des “green days” chaque mois.
\item Imagine que tu as lu un article sur la déforestation dans la région de l'Est du Cameroun. Rédige directement un paragraphe complet : présente le problème, une cause, deux effets, une solution, puis ton opinion.
\end{enumerate}
Contraintes obligatoires : quatre connecteurs différents, une expression d'opinion, une justification avec \eng{because}, aucune phrase recopiée.
\end{exercice}

\begin{corrige}[Exercice — éléments de réponse attendus]
Il n'existe pas une seule bonne réponse, mais tout bon paragraphe doit contenir : (a) une phrase d'introduction avec \eng{the text deals with...} ; (b) les quatre étapes de l'escalier reliées par \eng{first}, \eng{then} ou \eng{moreover}, \eng{however} ; (c) une opinion correctement formulée (\eng{I think that...} ou \eng{In my opinion,...}) suivie de \eng{because} ; (d) aucun calque du français (\eng{Me, I think}, \eng{I am agree}). Exemple pour le sujet 2 : \eng{The text deals with deforestation in the East region of Cameroon. First, it explains that farmers cut trees to create fields. Then, it shows two effects: the soil becomes poor and many animals lose their homes. However, the writer believes that tree-planting programmes can help. I agree that every small action counts, because if each pupil plants one tree every year, our forests will survive.}
\end{corrige}

\begin{aretenir}[L'essentiel de la section]
Un bon paragraphe « résumé + opinion » tient en 5 à 8 lignes, suit l'escalier \textbf{problème $\rightarrow$ causes $\rightarrow$ effets $\rightarrow$ solutions $\rightarrow$ opinion}, reformule sans recopier, s'appuie sur les connecteurs \eng{first, then, moreover, however, therefore, in conclusion}, et exprime l'opinion avec \eng{I think that...}, \eng{in my opinion...} ou \eng{I agree/disagree because...} — jamais avec un calque du français.
\end{aretenir}

\newpage

\coursec{Activité d'intégration}

\coursec{Chapter 3: Environment, Well-being and Health}
\courssub{Lesson 31 — Reading: Texts about environmental issues like climate change and global warming}

\courssub{Objectifs de l'activité}

À l'issue de cette activité d'intégration, l'élève sera capable de~:
\begin{itemize}
\item \cle{réutiliser} le lexique de l'environnement et du changement climatique étudié dans la leçon (\eng{air pollution}, \eng{deforestation}, \eng{global warming}, \eng{waste}, \eng{recycle}, \eng{reduce}...) dans une production authentique~;
\item \cle{produire} un message de sensibilisation court, clair et correct en anglais~: slogan, problèmes, solutions, phrase d'opinion~;
\item \cle{mobiliser} les structures grammaticales utiles~: l'impératif pour les appels à l'action (\eng{Plant trees!}), le présent simple pour les faits, \eng{We believe that...} pour l'opinion, les modaux \eng{must} / \eng{should} pour le devoir~;
\item \cle{présenter} oralement un document en une minute, avec une prononciation claire et un débit adapté~;
\item \cle{travailler en groupe}~: répartir les tâches, négocier les choix de langue, voter et justifier un choix.
\end{itemize}

\courssub{Situation}

Le club environnemental de ton lycée, le \eng{Green Club}, prépare la \eng{Green Week} (semaine verte). Le proviseur a accepté que le tableau d'affichage principal du lycée, près de la cantine, soit réservé aux affiches de sensibilisation rédigées par les élèves de Terminale. Voici l'annonce officielle publiée par le club~:

\begin{textebox}[Green Club — Official Announcement]
\textbf{GREEN WEEK POSTER COMPETITION!}

Dear students,

Our school is organising its first Green Week next month. We invite every class to create an awareness poster in English entitled “Save Our Planet!”

Your poster must include:
a catchy slogan; three environmental problems you can observe in your town or neighbourhood; three concrete solutions with action verbs; and one sentence giving your opinion.

The best posters will be displayed on the main notice board near the canteen, and the winning team will receive a certificate during the closing ceremony.

Deadline: Friday, 4 p.m. — Good luck, and remember: there is no Planet B!

The Green Club Committee
\end{textebox}

\begin{definition}[Glossaire de l'annonce]
\begin{itemize}
\item \eng{catchy} (adj.) : accrocheur, facile à retenir.
\item \eng{awareness} (n.) : sensibilisation ; \eng{an awareness poster} : une affiche de sensibilisation.
\item \eng{notice board} (n.) : tableau d'affichage.
\item \eng{canteen} (n.) : cantine.
\item \eng{deadline} (n.) : date limite.
\end{itemize}
\end{definition}

\courssub{Consignes}

Travaillez par groupes de quatre. Suivez les étapes dans l'ordre.

\textbf{Étape 1 — Brainstorming (en anglais) (10 min).} Dans votre cahier, listez au moins cinq problèmes environnementaux que vous observez dans votre ville ou votre quartier (par exemple à Yaoundé, Douala, Bafoussam ou Bamenda). Utilisez le lexique de la leçon~: \eng{air pollution}, \eng{plastic waste}, \eng{deforestation}, \eng{flooding}, \eng{water contamination}, \eng{global warming}. Choisissez ensuite ensemble les \cle{trois} problèmes les plus pertinents.

\textbf{Étape 2 — Les solutions (10 min).} Pour chaque problème retenu, proposez une solution concrète formulée avec un \cle{verbe d'action à l'impératif}~: \eng{plant}, \eng{recycle}, \eng{save}, \eng{reduce}, \eng{reuse}, \eng{clean}, \eng{protect}. Exemple~: \eng{Recycle plastic bottles to keep our streets clean.}

\textbf{Étape 3 — Le slogan et l'opinion (5 min).} Rédigez un slogan court et rythmé (maximum huit mots), puis une phrase d'opinion commençant par \eng{We believe that...}.

\textbf{Étape 4 — Fabrication de l'affiche (15 min).} Sur une feuille A3 ou une page de cahier, disposez~: le titre \eng{Save Our Planet!}, le slogan, les trois problèmes, les trois solutions et la phrase d'opinion. Soignez l'orthographe~: relisez-vous mutuellement.

\textbf{Étape 5 — Présentation orale (1 min par groupe).} Chaque membre du groupe prend la parole au moins une fois. Présentez le slogan, un problème, une solution et l'opinion. Parlez lentement et articulez.

\textbf{Étape 6 — Le vote (5 min).} La classe vote pour la meilleure affiche selon trois critères~: richesse du \cle{lexique}, \cle{correction grammaticale}, \cle{clarté du message}. Justifiez votre vote en anglais~: \eng{I vote for Group 3 because their slogan is catchy and their solutions are realistic.}

\courssub{Ressources à mobiliser}

\begin{aretenir}[Lexique utile de la leçon]
\begin{itemize}
\item Problèmes~: \eng{climate change}, \eng{global warming}, \eng{air pollution}, \eng{plastic waste}, \eng{deforestation}, \eng{flooding}, \eng{drought}, \eng{water contamination}, \eng{erosion}.
\item Solutions et verbes d'action~: \eng{plant trees}, \eng{recycle plastic}, \eng{save water and energy}, \eng{reduce waste}, \eng{reuse bags}, \eng{protect forests}, \eng{clean up the neighbourhood}.
\item Opinion~: \eng{We believe that...} / \eng{In our opinion,...} / \eng{We are convinced that...}.
\end{itemize}
\end{aretenir}

\begin{propriete}[L'impératif pour les appels à l'action]
L'impératif se forme avec la base verbale, sans sujet~: \eng{Plant trees!} (« Plantez des arbres ! »). À la négative~: \eng{Don't throw plastic on the ground!} (« Ne jetez pas de plastique par terre ! »). C'est la structure typique des slogans et des affiches de sensibilisation.
\end{propriete}

\begin{propriete}[Exprimer le devoir : \eng{must} et \eng{should}]
\eng{We must act now.} (« Nous devons agir maintenant » — obligation forte). \eng{Everyone should recycle.} (« Tout le monde devrait recycler » — conseil). Le modal est suivi de la base verbale, sans \eng{to}.
\end{propriete}

\begin{attention}[Erreurs fréquentes des francophones]
\begin{itemize}
\item Ne dites pas \eng{*the nature} pour « la nature » ; dites \eng{nature} sans article~: \eng{We must protect nature.}
\item \eng{environment} s'écrit avec un \textbf{n} avant le \textbf{m}~: en-vi-ron-\textbf{n}-ment. Prononcez « in-vaï-ron-ment ».
\item Après \eng{must} et \eng{should}, jamais de \eng{to}~: \eng{We must save water} (et non \eng{*must to save}).
\item \eng{advice} est indénombrable~: pas de \eng{*advices}.
\end{itemize}
\end{attention}

\courssub{Proposition de production et corrigé commenté}

\begin{exoresolu}[Modèle d'affiche et de présentation orale]
Produisez l'affiche complète du groupe et le script de sa présentation orale d'une minute.
\tcblower
\textbf{The poster (modèle)~:}

SAVE OUR PLANET!

Slogan~: “Green today, alive tomorrow!”

Problems we see in Yaoundé:
1. Plastic waste blocks our gutters and causes flooding.
2. Deforestation around the city increases global warming.
3. Open burning of rubbish pollutes the air we breathe.

Our solutions:
1. Recycle plastic bottles and reuse bags!
2. Plant a tree every month in your neighbourhood!
3. Don't burn rubbish — sort it and reduce waste!

Our opinion: We believe that young people can change the future if we act together now.

\textbf{The oral presentation (script)~:}

“Good morning, everyone. This is our poster, ‘Save Our Planet!’. Our slogan is: ‘Green today, alive tomorrow!’ In Yaoundé, plastic waste blocks the gutters and causes flooding during the rainy season. Forests are disappearing, and burning rubbish pollutes our air. Our solutions are simple: recycle plastic, plant trees every month, and stop burning rubbish. We believe that young people can change the future if we act together now. Thank you!”

\textbf{Commentaire des choix de langue~:}
\begin{itemize}
\item Le slogan joue sur la rime \eng{today / tomorrow} et l'antithèse \eng{green / alive}~: c'est ce qui le rend \eng{catchy}.
\item Les problèmes sont formulés au présent simple, temps des faits généraux~: \eng{plastic waste blocks... and causes...}.
\item Les solutions utilisent l'impératif, structure de l'affiche~: \eng{Recycle...! Plant...! Don't burn...!}.
\item L'opinion s'appuie sur \eng{We believe that...} suivi d'une subordonnée, avec le modal \eng{can} (capacité) et la structure de condition \eng{if we act together}.
\item Le lexique de la leçon est réemployé avec précision~: \eng{plastic waste}, \eng{flooding}, \eng{deforestation}, \eng{global warming}, \eng{pollute}, \eng{recycle}, \eng{plant}, \eng{reduce}.
\end{itemize}
\end{exoresolu}

\courssub{Bilan de l'activité}

Cette activité a permis de transformer une compétence de réception (lire un texte sur le changement climatique) en une compétence de production~: écrire et présenter un message de sensibilisation. Vous avez réinvesti le lexique de l'environnement dans un contexte réel et local, pratiqué l'impératif, les modaux \eng{must} et \eng{should}, ainsi que la structure d'opinion \eng{We believe that...}. Vous avez également travaillé l'expression orale en public et l'argumentation lors du vote justifié.

\begin{methode}[Grille d'auto-évaluation]
\begin{enumerate}
\item Mon affiche contient-elle les cinq éléments exigés (titre, slogan, trois problèmes, trois solutions, opinion)~?
\item Ai-je utilisé au moins six mots du lexique de la leçon, correctement orthographiés~?
\item Mes impératifs et mes modaux sont-ils corrects (base verbale, pas de \eng{to} après \eng{must/should})~?
\item Ai-je parlé au moins une fois pendant la présentation, en articulant~?
\item Ai-je justifié mon vote en anglais avec \eng{because}~?
\end{enumerate}
\end{methode}

\begin{savaistu}[“There is no Planet B”]
Ce slogan, devenu célèbre lors des marches pour le climat dans le monde entier, rappelle que la Terre est notre seule maison. Au Cameroun, des initiatives comme le reboisement autour du mont Cameroun ou les campagnes de nettoyage des plages de Kribi et Limbe montrent que les jeunes peuvent passer du slogan à l'action — exactement ce que votre affiche propose.
\end{savaistu}

\newpage

\coursec{Méthodes et exercices résolus}

\coursec{Lesson 31 — Reading: Texts about environmental issues like climate change and global warming}

\courssub{Warm-up : Découverte du thème}

\begin{experience}[Discussion guidée : le climat au Cameroun]
En binôme, observez la carte mentale ci-dessous et répondez oralement aux questions. Utilisez le vocabulaire de la boîte \textbf{Key Words}.
\begin{center}
\begin{popfigure}
\begin{tikzpicture}[
  mindmap,
  concept color=popblue,
  level 1/.style={level distance=3.5cm, sibling angle=90},
  level 2/.style={level distance=2.5cm, sibling angle=45},
  every node/.style={font=\small, align=center},
  grow cyclic
]
\node[concept, align=center]{Climate Change\\ (Changement climatique)}
  child[concept color=poporange] { node[concept, align=center]{Causes\\ (Causes)}
    child { node[align=center]{Burning fossil fuels\\ (Combustibles fossiles)} }
    child { node[align=center]{Deforestation\\ (Déforestation)} }
    child { node[align=center]{Industry \& transport\\ (Industrie, transport)} }
  }
  child[concept color=popgreen] { node[concept, align=center]{Effects\\ (Effets)}
    child { node[align=center]{Rising temperatures\\ (Hausse températures)} }
    child { node[align=center]{Unpredictable rainfall\\ (Pluies imprévisibles)} }
    child { node[align=center]{Floods \& droughts\\ (Inondations, sécheresses)} }
    child { node[align=center]{Sea level rise\\ (Montée des eaux)} }
  }
  child[concept color=poppurple] { node[concept, align=center]{Solutions\\ (Solutions)}
    child { node[align=center]{Renewable energy\\ (Énergies renouvelables)} }
    child { node[align=center]{Reforestation\\ (Reboisement)} }
    child { node[align=center]{Adaptation\\ (Adaptation)} }
  }
  child[concept color=poppink] { node[concept, align=center]{Cameroon Context\\ (Contexte camerounais)}
    child { node[align=center]{Lake Chad shrinking\\ (Lac Tchad)} }
    child { node[align=center]{Cocoa/coffee at risk\\ (Cacao/café menacés)} }
    child { node[align=center]{Coastal cities threatened\\ (Villes côtières)} }
  };
\end{tikzpicture}
\legende{Carte mentale : causes, effets, solutions et contexte camerounais du changement climatique}
\end{popfigure}
\end{center}

\textbf{Key Words} : \eng{climate change / global warming / greenhouse gases / carbon dioxide / emissions / deforestation / drought / flood / rising sea levels / adaptation / mitigation / renewable energy / reforestation / sustainable development}.

\begin{enumerate}
\item \eng{What do you notice about the rainy seasons in your region compared to ten years ago?}
\item \eng{Which human activities in Cameroon contribute most to \cle{greenhouse gas emissions}?}
\item \eng{How does \cle{deforestation} affect the climate and agriculture in the Centre and South regions?}
\end{enumerate}
\textbf{Consigne :} Chaque binôme présente une idée à la classe. Le professeur note le vocabulaire clé au tableau.
\end{experience}

\courssub{Texte support : Article de lecture}

\begin{textebox}[Article original — Climate Change Hits Cameroon’s Agriculture Hard]
\eng{For the third year in a row, farmers in the Centre Region of Cameroon have watched the rains arrive late. According to climate experts, average temperatures in Central Africa have risen steadily over the past decades, and rainfall patterns have become unpredictable. Cocoa and coffee producers, who depend on regular seasons, are among the most affected. “We no longer know when to plant,” says one farmer near Obala. Scientists warn that without adaptation — such as drought-resistant seeds and better water management — harvests could fall sharply, threatening both food supplies and family incomes.}

\eng{Global warming is mainly caused by greenhouse gases, such as carbon dioxide, which trap heat in the atmosphere. The burning of fossil fuels — coal, oil and gas — for transport and industry is responsible for most of these emissions. Deforestation makes the problem worse, because trees absorb carbon dioxide. As a result of rising temperatures, glaciers are melting and sea levels are rising, threatening coastal cities like Douala and Limbe.}

\eng{Climate change is already affecting Africa. Temperatures are rising faster than the global average, and extreme weather events such as floods and droughts are becoming more frequent. For example, Lake Chad has shrunk dramatically over the last fifty years, depriving millions of people of water and fish. Experts say that African countries, although they produce very little carbon dioxide, suffer the most from its effects. They are calling on rich nations to finance adaptation projects, such as solar energy and reforestation programmes.}
\end{textebox}

\courssub{Méthode 1 : Lire et comprendre un texte sur l'environnement (skimming)}

\begin{methode}[Comprendre le sens général d'un texte sur le climat]
Le \cle{skimming} consiste à lire rapidement pour saisir l'idée générale (\eng{gist}) sans s'arrêter sur les mots inconnus.
\begin{enumerate}
\item Lire le \textbf{titre}, le \textbf{sous-titre} et la \textbf{source} : ils annoncent le sujet (ex. \eng{climate change}, \eng{global warming}, \eng{deforestation}).
\item Lire la \textbf{première phrase de chaque paragraphe} (\eng{topic sentence}) : elle contient l'idée principale du paragraphe.
\item Repérer les \textbf{mots répétés} : ce sont les mots-clés du thème (\eng{emissions}, \eng{temperatures}, \eng{floods}, \eng{drought}).
\item Formuler le sujet en une phrase : \eng{This text is about... / The writer explains that...}
\item Ne pas paniquer devant un mot inconnu : au skimming, on l'ignore ; on y reviendra au scanning.
\end{enumerate}
\textbf{Réflexe examen :} une question de type \eng{What is the text about?} ou \eng{Choose the best title} se résout uniquement par skimming, en 2 minutes maximum.
\end{methode}

\begin{exoresolu}[Skimming : trouver l'idée générale]
\textbf{Énoncé.} Read the following text quickly and answer: (a) What is the text about? (b) Choose the best title: 1. \eng{The History of Cameroon} 2. \eng{Rising Temperatures Threaten Cameroon's Farms} 3. \eng{How to Plant Cocoa Trees}.

\eng{For the third year in a row, farmers in the Centre Region of Cameroon have watched the rains arrive late. According to climate experts, average temperatures in Central Africa have risen steadily over the past decades, and rainfall patterns have become unpredictable. Cocoa and coffee producers, who depend on regular seasons, are among the most affected. “We no longer know when to plant,” says one farmer near Obala. Scientists warn that without adaptation — such as drought-resistant seeds and better water management — harvests could fall sharply, threatening both food supplies and family incomes.}

\tcblower
\textbf{Solution détaillée.}
\begin{itemize}
\item \textbf{Étape 1 — indices du titre et des mots répétés :} les mots \eng{rains}, \eng{temperatures}, \eng{farmers}, \eng{harvests}, \eng{drought} reviennent plusieurs fois : le thème est le climat et l'agriculture.
\item \textbf{Étape 2 — premières phrases :} la première phrase parle des pluies tardives ; la dernière parle des conséquences sur les récoltes. Le texte relie donc réchauffement climatique et agriculture camerounaise.
\item \textbf{(a) Réponse rédigée :} \eng{The text is about the effects of climate change — rising temperatures and unpredictable rainfall — on Cameroonian farmers and their harvests.} (« Le texte traite des effets du changement climatique — hausse des températures et pluies imprévisibles — sur les agriculteurs camerounais et leurs récoltes. »)
\item \textbf{(b) Meilleur titre :} titre 2, \eng{Rising Temperatures Threaten Cameroon's Farms}. Le titre 1 est trop large (le texte n'est pas historique) ; le titre 3 est faux (le texte ne donne pas de conseils de plantation).
\end{itemize}
\textbf{Conclusion :} au skimming, on ne traduit pas tout : on croise mots répétés + premières phrases pour dégager le thème en une phrase.
\end{exoresolu}

\courssub{Méthode 2 : Repérer une information précise (scanning)}

\begin{methode}[Répondre aux questions de compréhension par scanning]
Le \cle{scanning} consiste à chercher une information précise (chiffre, date, nom, cause, conséquence) sans relire tout le texte.
\begin{enumerate}
\item Lire la \textbf{question d'abord} et souligner le mot-clé (\eng{when}, \eng{why}, \eng{who}, \eng{how much}).
\item Repérer dans le texte le \textbf{mot-clé ou son synonyme} : les examinateurs reformulent (\eng{increase} = \eng{rise} ; \eng{decrease} = \eng{fall} ; \eng{harmful} = \eng{dangerous}).
\item Lire la phrase entière autour du mot repéré, puis la phrase précédente et la suivante.
\item Rédiger la réponse avec \textbf{ses propres mots}, en phrase complète, sans recopier un long passage.
\item Vérifier que la réponse répond bien à la question posée (cause ? conséquence ? opinion de l'auteur ?).
\end{enumerate}
\textbf{Réflexe :} pour une question \eng{Why...?}, chercher dans le texte \eng{because}, \eng{due to}, \eng{as a result of}, \eng{that is why}.
\end{methode}

\begin{exoresolu}[Scanning : questions de compréhension]
\textbf{Énoncé.} Read the text and answer the questions in complete sentences.

\eng{Global warming is mainly caused by greenhouse gases, such as carbon dioxide, which trap heat in the atmosphere. The burning of fossil fuels — coal, oil and gas — for transport and industry is responsible for most of these emissions. Deforestation makes the problem worse, because trees absorb carbon dioxide. As a result of rising temperatures, glaciers are melting and sea levels are rising, threatening coastal cities like Douala and Limbe.}

(a) \eng{What causes global warming?} (b) \eng{Why does deforestation make the problem worse?} (c) \eng{Which two Cameroonian cities are threatened by rising sea levels?}

\tcblower
\textbf{Solution détaillée.}
\begin{itemize}
\item \textbf{(a)} Mot-clé de la question : \eng{causes}. Dans le texte : \eng{is mainly caused by greenhouse gases... The burning of fossil fuels...}. Réponse : \eng{Global warming is mainly caused by greenhouse gases, especially carbon dioxide produced by the burning of fossil fuels.} (« Le réchauffement climatique est principalement causé par les gaz à effet de serre, notamment le dioxyde de carbone issu de la combustion des énergies fossiles. »)
\item \textbf{(b)} Question \eng{Why} : on repère \eng{because} dans le texte : \eng{because trees absorb carbon dioxide}. Réponse : \eng{Deforestation makes the problem worse because trees absorb carbon dioxide, so cutting them down leaves more of this gas in the atmosphere.} (« ...parce que les arbres absorbent le dioxyde de carbone ; les abattre laisse donc davantage de ce gaz dans l'atmosphère. »)
\item \textbf{(c)} Repérage de noms propres (majuscules faciles à scanner) : \eng{Douala and Limbe}. Réponse : \eng{Douala and Limbe are threatened by rising sea levels.}
\end{itemize}
\textbf{Conclusion :} chaque réponse est une phrase complète, reformulée, et justifiée par un indice précis du texte (\eng{caused by}, \eng{because}, noms propres).
\end{exoresolu}

\courssub{Lexique thématique : L'environnement et le climat}

\begin{definition}[Vocabulaire essentiel — Noms, verbes, adjectifs et collocations]
Le tableau ci-dessous regroupe le lexique indispensable pour lire, comprendre et produire un texte sur l'environnement au niveau Terminale. Apprenez les \textbf{collocations} (associations de mots naturelles) : elles font la différence à l'écrit et à l'oral.
\begin{center}
\begin{tabularx}{\linewidth}{|X|X|X|}
\hline
\rowcolor{popblueL}
\textbf{English} & \textbf{Français} & \textbf{Collocations / Exemples} \\
\hline
\eng{climate change} & changement climatique & \eng{fight climate change, climate change mitigation} \\
\eng{global warming} & réchauffement climatique & \eng{limit global warming to 1.5°C} \\
\eng{greenhouse gases} & gaz à effet de serre & \eng{greenhouse gas emissions, reduce emissions} \\
\eng{carbon dioxide (CO2)} & dioxyde de carbone & \eng{carbon footprint, carbon-neutral} \\
\eng{deforestation} & déforestation & \eng{combat deforestation, illegal logging} \\
\eng{drought} & sécheresse & \eng{prolonged drought, drought-resistant crops} \\
\eng{flood} & inondation & \eng{flash flood, flood-prone areas} \\
\eng{rising sea levels} & montée du niveau de la mer & \eng{threaten coastal communities} \\
\eng{renewable energy} & énergie renouvelable & \eng{solar power, wind energy, hydroelectricity} \\
\eng{reforestation} & reboisement & \eng{reforestation programme, plant trees} \\
\eng{adaptation} & adaptation & \eng{adaptation strategies, adapt to new conditions} \\
\eng{sustainable development} & développement durable & \eng{sustainable agriculture, sustainable future} \\
\hline
\end{tabularx}
\end{center}

\textbf{Verbes clés :} \eng{emit (émettre)}, \eng{trap (piéger)}, \eng{absorb (absorber)}, \eng{melt (fondre)}, \eng{shrink (rétrécir)}, \eng{deprive sb of (priver qqn de)}, \eng{finance (financer)}, \eng{raise awareness (sensibiliser)}.

\textbf{Adjectifs clés :} \eng{unpredictable (imprévisible)}, \eng{detrimental / harmful (nuisible)}, \eng{resilient (résilient)}, \eng{sustainable (durable)}, \eng{extreme (extrême)}.
\end{definition}

\begin{popfigure}
\begin{tikzpicture}[
  mindmap,
  concept color=popgreen,
  level 1/.style={level distance=3cm, sibling angle=90},
  level 2/.style={level distance=2.2cm, sibling angle=45},
  every node/.style={font=\footnotesize, align=center},
  grow cyclic
]
\node[concept, align=center]{Environment\\ Vocabulary}
  child[concept color=popblue] { node[concept] {Nouns}
    child { node {emissions} }
    child { node {pollution} }
    child { node {conservation} }
    child { node {biodiversity} }
  }
  child[concept color=poporange] { node[concept] {Verbs}
    child { node {pollute} }
    child { node {conserve} }
    child { node {recycle} }
    child { node {sustain} }
  }
  child[concept color=poppurple] { node[concept] {Adjectives}
    child { node {eco-friendly} }
    child { node {biodegradable} }
    child { node {toxic} }
    child { node {green} }
  }
  child[concept color=poppink] { node[concept] {Expressions}
    child { node {carbon footprint} }
    child { node {go green} }
    child { node {climate action} }
    child { node {net zero} }
  };
\end{tikzpicture}
\legende{Réseau lexical : environnement (noms, verbes, adjectifs, expressions)}
\end{popfigure}

\courssub{Méthode 3 : Déduire le sens des mots inconnus}

\begin{methode}[Deviner le vocabulaire sans dictionnaire]
\begin{enumerate}
\item Identifier la \textbf{nature du mot} (nom, verbe, adjectif) grâce à sa place et à sa terminaison : \eng{-tion}, \eng{-ness}, \eng{-ity} = nom ; \eng{-ful}, \eng{-able}, \eng{-ive} = adjectif ; \eng{-ise/-ize}, \eng{-ate} = verbe.
\item Exploiter le \textbf{contexte} : synonymes, contraires (\eng{but}, \eng{whereas}, \eng{unlike}, \eng{however}), exemples (\eng{such as}, \eng{for example}), définitions implicites (appositions, relatives).
\item Décomposer le mot : préfixe (\eng{re-} = encore ; \eng{un-}, \eng{in-}, \eng{im-} = négation ; \eng{over-} = excès), racine connue (\eng{forest} dans \eng{deforestation}), suffixe.
\item Vérifier l'hypothèse en remplaçant le mot par sa traduction supposée : la phrase doit rester logique.
\item Se méfier des \textbf{faux amis} : \eng{actually} = « en fait » (pas « actuellement » \eng{currently}) ; \eng{sensible} = « raisonnable » (pas « sensible » \eng{sensitive}) ; \eng{issue} = « problème / question » (pas seulement « numéro »).
\end{enumerate}
\end{methode}

\begin{exoresolu}[Déduction lexicale en contexte]
\textbf{Énoncé.} Find the meaning of the underlined words from the context. Choose from: \eng{very dry periods / floods / to adjust / harmful}.

\eng{1. Prolonged droughts have destroyed crops in the Far North Region, where no rain has fallen for months. 2. Farmers must adapt to the new climate conditions if they want to survive. 3. Air pollution is detrimental to children's health, whereas clean air helps them grow strong.}

\tcblower
\textbf{Solution détaillée.}
\begin{itemize}
\item \textbf{1.} \eng{droughts} : le contexte donne un indice explicite — \eng{no rain has fallen for months} — et la conséquence \eng{destroyed crops}. Nature : nom pluriel (terminaison \eng{-s}). Sens : \eng{very dry periods} (« sécheresses »).
\item \textbf{2.} \eng{adapt} : place après modal \eng{must} = verbe à l'infinitif ; complément \eng{to the new climate conditions} suggère un ajustement. Sens : \eng{to adjust} (« s'adapter »). Famille : \eng{adaptation} (nom), \eng{adaptable} (adj.).
\item \textbf{3.} \eng{detrimental} : adjectif (après \eng{is}) ; l'opposition marquée par \eng{whereas clean air helps them} indique un contraire de « bénéfique ». Sens : \eng{harmful} (« nuisible »). Synonyme : \eng{damaging}.
\end{itemize}
\textbf{Conclusion :} nature du mot + indice de contexte (définition, opposition, exemple) permettent de répondre sans dictionnaire — compétence directement évaluée au Bac.
\end{exoresolu}

\courssub{Méthode 4 : Résumer un texte en quelques phrases}

\begin{methode}[Rédiger un résumé court (summary)]
\begin{enumerate}
\item Repérer les \textbf{3 ou 4 idées principales} (une par paragraphe) ; ignorer exemples, chiffres et détails secondaires.
\item Utiliser des \textbf{verbes de reformulation} : \eng{The text explains that... / The writer shows that... / According to the author... / The article highlights...}
\item Relier les idées avec des connecteurs logiques : \eng{first / firstly}, \eng{moreover / furthermore}, \eng{as a result / consequently}, \eng{finally / in conclusion}, \eng{although / while} (concession).
\item Respecter la \textbf{limite de mots} (souvent 30 à 50 mots) et ne jamais recopier de phrases entières du texte (pas de \eng{copy-paste}).
\item Ne pas donner son opinion dans le résumé : elle vient après, dans une partie séparée (\eng{Personal response}).
\end{enumerate}
\end{methode}

\begin{exoresolu}[Summary : résumer un article]
\textbf{Énoncé.} Summarize the following text in about 40 words.

\eng{Climate change is already affecting Africa. Temperatures are rising faster than the global average, and extreme weather events such as floods and droughts are becoming more frequent. For example, Lake Chad has shrunk dramatically over the last fifty years, depriving millions of people of water and fish. Experts say that African countries, although they produce very little carbon dioxide, suffer the most from its effects. They are calling on rich nations to finance adaptation projects, such as solar energy and reforestation programmes.}

\tcblower
\textbf{Solution détaillée.}
\begin{itemize}
\item \textbf{Idées principales :} (1) l'Afrique subit déjà le changement climatique (chaleur, événements extrêmes) ; (2) elle pollue peu mais souffre le plus (injustice) ; (3) les experts demandent aux pays riches de financer l'adaptation. Le détail du lac Tchad est un exemple : on l'écarte.
\item \textbf{Résumé rédigé (environ 40 mots) :} \eng{The text explains that Africa is already suffering from climate change, with rising temperatures and extreme weather. Although Africans produce little pollution, they are the most affected, so experts want rich countries to finance adaptation projects.} (« Le texte explique que l'Afrique souffre déjà du changement climatique... ; bien que les Africains polluent peu, ils sont les plus touchés, c'est pourquoi les experts veulent que les pays riches financent des projets d'adaptation. »)
\item \textbf{Vérification :} aucune phrase copiée, connecteurs utilisés (\eng{although}, \eng{so}), opinion absente, limite respectée (39 mots).
\end{itemize}
\end{exoresolu}

\courssub{Méthode 5 : Donner son opinion en quelques phrases}

\begin{methode}[Exprimer et justifier son opinion (giving your opinion)]
\begin{enumerate}
\item Annoncer clairement sa position : \eng{In my opinion... / I believe that... / As far as I am concerned... / From my point of view...}
\item Justifier avec \eng{because}, \eng{since}, \eng{as}, \eng{that is why} : une opinion sans justification ne vaut rien à l'examen.
\item Nuancer éventuellement : \eng{I agree that..., but I think...} ou \eng{To some extent... / While it is true that..., I would argue that...}
\item Conclure par une proposition concrète : \eng{That is why we should... / The government ought to... / It is essential to...}
\item Employer le présent simple pour les vérités générales et \eng{should / must / ought to / have to} pour les recommandations. Utiliser le conditionnel (\eng{would}) pour les hypothèses.
\end{enumerate}
\textbf{Structures à retenir :} \eng{I strongly believe that everyone should protect the environment.} / \eng{In my view, planting trees is the cheapest solution.} / \eng{If we act now, we will save the planet.}
\end{methode}

\begin{exoresolu}[Opinion : réagir à un texte sur le réchauffement]
\textbf{Énoncé.} \eng{Do you think young people can help fight climate change? Give your opinion in four or five sentences.}

\tcblower
\textbf{Solution détaillée.}
\begin{itemize}
\item \textbf{Plan :} position (oui) → deux arguments (gestes quotidiens, sensibilisation) → conclusion avec recommandation.
\item \textbf{Réponse rédigée :} \eng{In my opinion, young people can play an important role in the fight against climate change. First, they can adopt simple habits, such as planting trees, avoiding plastic bags and saving water at home and at school. Moreover, young people are very active on social media, so they can raise awareness among their friends and families. I believe that if every student in Cameroon planted just one tree a year, the results would be enormous. That is why schools should create environmental clubs in every town.} (« À mon avis, les jeunes peuvent jouer un rôle important... D'abord, ils peuvent adopter des habitudes simples... De plus, ils sont très actifs sur les réseaux sociaux et peuvent sensibiliser leur entourage... C'est pourquoi les écoles devraient créer des clubs environnementaux dans chaque ville. »)
\item \textbf{Justifications grammaticales :} \eng{can} + base verbale (capacité) ; \eng{should} + base verbale (recommandation) ; conditionnel type 2 \eng{if every student planted..., the results would be...} (hypothèse) ; connecteurs \eng{first}, \eng{moreover}, \eng{that is why}.
\end{itemize}
\textbf{Conclusion :} opinion claire + deux arguments + conclusion concrète = réponse complète et bien notée.
\end{exoresolu}

\begin{savaistu}[Le saviez-vous ? — Le Lac Tchad et la Grande Muraille Verte]
Le \textbf{Lac Tchad} a perdu environ 90\% de sa surface depuis les années 1960 à cause du changement climatique et de la surexploitation de l'eau. Cela affecte 30 millions de personnes au Cameroun, Tchad, Niger, Nigeria. Pour lutter contre la désertification, l'Union Africaine a lancé la \textbf{Grande Muraille Verte} (\eng{Great Green Wall}) : une bande d'arbres de 8 000 km à travers le Sahel, du Sénégal à Djibouti. Le Cameroun y participe activement dans l'Extrême-Nord. C'est un exemple concret d'\eng{adaptation} et de \eng{reforestation} à grande échelle.
\end{savaistu}

\begin{attention}[Les 5 erreurs qui coûtent des points au Bac]
\begin{enumerate}
\item \textbf{Faux amis :} \eng{actually} = « en fait » (pas « actuellement » $\rightarrow$ \eng{currently}) ; \eng{to assist} = « aider » (pas « assister à » $\rightarrow$ \eng{to attend}) ; \eng{an issue} = « un problème / une question » (pas « un numéro » $\rightarrow$ \eng{a number}) ; \eng{eventually} = « finalement » (pas « éventuellement » $\rightarrow$ \eng{possibly}).
\item \textbf{Orthographe du lexique :} \eng{environment} (un \textbf{n} avant le \textbf{m}), \eng{government} (un \textbf{n} après \eng{govern-}), \eng{pollution} (double \textbf{ll} : p-o-\textbf{ll}-u-t-i-o-n), \eng{climate} (pas de \textbf{h}), \eng{agriculture} (pas \emph{agricuture}), \eng{temperature} (pas \emph{temperatur}).
\item \textbf{Noms indénombrables (uncountable) :} \eng{information}, \eng{advice}, \eng{pollution}, \eng{progress}, \eng{evidence}, \eng{equipment}, \eng{furniture}, \eng{homework} sont \textbf{singuliers uniquement} : jamais de \eng{-s} ni de \eng{a/an}. \eng{The information is useful} (pas \emph{are}), \eng{a piece of advice} (pas \emph{an advice}).
\item \textbf{Calques du français :} « le réchauffement de la planète » = \eng{global warming} (pas \emph{the reheating of the planet}) ; « protéger la nature » = \eng{protect the environment} (pas \emph{protect the nature}) ; « changer d'avis » = \eng{change one's mind} (pas \emph{change opinion}) ; « prendre conscience » = \eng{become aware} (pas \emph{take conscience}).
\item \textbf{Accords et temps dans les réponses :} répondre au temps de la question ; \eng{The text says that temperatures are rising} (présent continu pour évolution en cours) ; \textbf{-s} obligatoire à la 3\up{e} personne du singulier au présent simple : \eng{Climate change affects farmers} (pas \emph{affect}) ; \eng{The government has decided} (pas \emph{have}).
\end{enumerate}
\end{attention}

\newpage

\coursec{Exercices}

\courssub{Je vérifie mes connaissances}

\begin{exercice}[QCM : Compréhension globale]
\textbf{Read the following statements and choose the correct answer (A, B, C or D).}
\begin{textebox}[Extrait d'article sur le réchauffement climatique]
Global warming refers to the long-term rise in the average temperature of the Earth's climate system. It is mainly caused by increased concentrations of greenhouse gases in the atmosphere, resulting from human activities such as burning fossil fuels, deforestation and industrial farming. The consequences are already visible: melting ice caps, rising sea levels and more frequent extreme weather events.
\end{textebox}
\begin{enumerate}
\item What is the main cause of global warming mentioned in the text?
      \begin{enumerate}
      \item Natural cycles of the Earth.
      \item Increased concentrations of greenhouse gases.
      \item Volcanic eruptions.
      \item Changes in solar radiation.
      \end{enumerate}
\item Which human activity is \textbf{NOT} listed as a cause in the text?
      \begin{enumerate}
      \item Burning fossil fuels.
      \item Deforestation.
      \item Industrial farming.
      \item Using renewable energy.
      \end{enumerate}
\item According to the text, what is a visible consequence of global warming?
      \begin{enumerate}
      \item Decreasing sea levels.
      \item Melting ice caps.
      \item Fewer extreme weather events.
      \item Lower average temperatures.
      \end{enumerate}
\item The phrase \eng{“long-term rise”} suggests that the process is:
      \begin{enumerate}
      \item Sudden and temporary.
      \item Gradual and continuing over time.
      \item Immediate and reversible.
      \item Seasonal and predictable.
      \end{enumerate}
\end{enumerate}
\end{exercice}

\begin{exercice}[Vrai ou Faux justifié]
\textbf{Read the text below. Say whether the statements are True or False. Justify your answer by quoting a specific phrase from the text.}
\begin{textebox}[Plastic pollution in oceans]
Every year, millions of tonnes of plastic waste end up in the oceans. This pollution poses a serious threat to marine life. Sea turtles often mistake plastic bags for jellyfish and eat them, which can be fatal. Microplastics, tiny fragments of degraded plastic, have entered the food chain and are now found in the fish we eat. Governments and organisations are trying to ban single-use plastics and promote recycling, but the problem remains huge.
\end{textebox}
\begin{enumerate}
\item Plastic waste in the oceans only affects large marine animals like whales.
\item Microplastics are created when larger plastic items break down over time.
\item Single-use plastics have already been completely banned worldwide.
\item The text suggests that humans consume plastic through the food chain.
\end{enumerate}
\end{exercice}

\begin{exercice}[Vocabulaire : Mots-clés de l'environnement]
\textbf{Match each English word (1--8) with its French definition (a--h). Then, use four words of your choice in personal sentences.}
\begin{center}
\begin{tabular}{|c|l|c|l|}
\hline
\textbf{English} & & \textbf{Français} & \\
\hline
1. Drought & & a. Gaz à effet de serre \\
2. Emissions & & b. Sécheresse \\
3. Greenhouse gas & & c. Déforestation \\
4. Deforestation & & d. Émissions \\
5. Sustainable & & e. Énergies renouvelables \\
6. Renewable energy & & f. Durable / Soutenable \\
7. Carbon footprint & & g. Empreinte carbone \\
8. Biodiversity & & h. Biodiversité \\
\hline
\end{tabular}
\end{center}
\textbf{Personal sentences:}
\begin{enumerate}
\item \dots
\item \dots
\item \dots
\item \dots
\end{enumerate}
\end{exercice}

\begin{exercice}[Grammaire : Passif et modaux d'obligation/nécessité]
\textbf{Complete the sentences with the correct form of the verb in brackets (Passive voice or Modal + Passive Infinitive).}
\begin{enumerate}
\item The forest \eng{\_\_\_\_\_\_\_\_\_\_\_\_} (destroy) by fire last summer.
\item New laws \eng{\_\_\_\_\_\_\_\_\_\_\_\_} (must / implement) to protect endangered species immediately.
\item Plastic bottles \eng{\_\_\_\_\_\_\_\_\_\_\_\_} (can / recycle) into fleece jackets.
\item The agreement \eng{\_\_\_\_\_\_\_\_\_\_\_\_} (sign) by all member countries next week.
\item Something \eng{\_\_\_\_\_\_\_\_\_\_\_\_} (has to / do) to reduce carbon emissions before it is too late.
\item The river \eng{\_\_\_\_\_\_\_\_\_\_\_\_} (pollute) by chemical waste for years before the factory closed.
\end{enumerate}
\end{exercice}

\courssub{Je m'entraîne}

\begin{exercice}[Traduction : Thème (Français $\rightarrow$ Anglais)]
\textbf{Translate the following sentences into English. Pay attention to vocabulary (environment) and grammar (passive, modals, connectors).}
\begin{enumerate}
\item La déforestation en Amazonie menace la biodiversité et accélère le changement climatique.
\item Les gouvernements doivent interdire les plastiques à usage unique pour réduire la pollution des océans.
\item Si nous n'agissons pas maintenant, les générations futures subiront les conséquences du réchauffement climatique.
\item Des énergies renouvelables comme le solaire et l'éolien sont en train d'être développées dans de nombreux pays africains.
\item Il est essentiel que chacun réduise son empreinte carbone en changeant ses habitudes de consommation.
\end{enumerate}
\end{exercice}

\begin{exercice}[Transformation de phrases]
\textbf{Rewrite the sentences as indicated. Do not change the meaning.}
\begin{enumerate}
\item \eng{People burn fossil fuels to produce energy.} (Turn into \textbf{Passive voice})
\item \eng{“We must protect the rainforest,” the scientist said.} (Turn into \textbf{Reported speech})
\item \eng{The government will ban petrol cars in 2035.} (Turn into \textbf{Passive voice, future})
\item \eng{Climate change causes rising sea levels.} (Rewrite using: \eng{Rising sea levels \_\_\_\_\_\_\_\_\_\_\_\_})
\item \eng{It is necessary to plant more trees.} (Rewrite using: \eng{More trees \_\_\_\_\_\_\_\_\_\_\_\_})
\item \eng{Although it was raining, they planted trees.} (Rewrite using: \eng{In spite of \_\_\_\_\_\_\_\_\_\_\_\_})
\end{enumerate}
\end{exercice}

\begin{exercice}[Texte à trous : Cloze test]
\textbf{Complete the text with the words from the box. There are two extra words.}
\begin{center}
\emph{threat \quad sustainable \quad emissions \quad drought \quad recycle \quad greenhouse \quad waste \quad ban \quad rainfall \quad harmful}
\end{center}
\begin{textebox}[Climate Action in Cameroon]
Cameroon faces several environmental challenges. Irregular \eng{(1)\_\_\_\_\_\_\_\_\_\_\_\_} patterns lead to severe \eng{(2)\_\_\_\_\_\_\_\_\_\_\_\_} in the North, affecting agriculture. In cities like Douala, plastic \eng{(3)\_\_\_\_\_\_\_\_\_\_\_\_} clogs drainage systems, causing floods. The government plans to \eng{(4)\_\_\_\_\_\_\_\_\_\_\_\_} non-biodegradable plastic bags. Industries must reduce their carbon \eng{(5)\_\_\_\_\_\_\_\_\_\_\_\_} to limit the \eng{(6)\_\_\_\_\_\_\_\_\_\_\_\_} effect. Promoting \eng{(7)\_\_\_\_\_\_\_\_\_\_\_\_} farming methods is crucial for food security. Citizens are encouraged to \eng{(8)\_\_\_\_\_\_\_\_\_\_\_\_} paper, glass and metal. Every small action helps fight this global \eng{(9)\_\_\_\_\_\_\_\_\_\_\_\_}.
\end{textebox}
\end{exercice}

\begin{exercice}[Appariement : Causes, Conséquences, Solutions]
\textbf{Match each cause (1--4) with its likely consequence (A--D) and a possible solution (I--IV). Write your answers as triplets (e.g., 1 -- A -- I).}
\begin{center}
\begin{tabular}{|c|l|c|l|c|l|}
\hline
\textbf{Causes} & & \textbf{Conséquences} & & \textbf{Solutions} & \\
\hline
1. Burning fossil fuels & & A. Rising sea levels & & I. Reforestation programmes \\
2. Cutting down trees (deforestation) & & B. Air pollution / Smog & & II. Switch to solar/wind energy \\
3. Excessive use of chemical fertilizers & & C. Soil erosion / Desertification & & III. Organic farming / Composting \\
4. Dumping industrial waste in rivers & & D. Water contamination & & IV. Strict fines / Water treatment plants \\
\hline
\end{tabular}
\end{center}
\textbf{Answers:}
\begin{enumerate}
\item 1 -- \dots -- \dots
\item 2 -- \dots -- \dots
\item 3 -- \dots -- \dots
\item 4 -- \dots -- \dots
\end{enumerate}
\end{exercice}

\courssub{Je me prépare au Bac}

\begin{exercice}[Compréhension écrite type Baccalauréat (Texte + Questions)]
\textbf{Read the following text carefully and answer the questions that follow.}
\begin{textebox}[The Green Belt Movement: A Local Answer to a Global Crisis]
Wangari Maathai, the first African woman to win the Nobel Peace Prize, understood the link between environment, democracy and peace long before it became a global slogan. In 1977, she founded the Green Belt Movement in Kenya. The idea was simple but powerful: pay rural women a small stipend to plant trees. 

At the time, Kenya faced massive deforestation. Forests were cleared for timber and commercial agriculture. The consequences were immediate: soil erosion dried up rivers, firewood became scarce, and women had to walk longer distances to collect it. Maathai realized that environmental degradation hurt the poorest first, especially women.

The movement grew rapidly. Over 51 million trees have been planted across Kenya. But the Green Belt Movement did more than reforest land. It empowered women. By earning income and managing nurseries, they gained confidence and a voice in their communities. They also learned about civic education, voting rights and sustainable farming techniques.

Maathai often said: “When we plant trees, we plant the seeds of peace and seeds of hope.” She fought against land grabbing by powerful politicians, was arrested and beaten, but never gave up. Today, the movement is a model for community-based conservation across Africa. It proves that local action, driven by ordinary citizens, can mitigate the effects of climate change and build resilience.

\textit{(Approx. 220 words)}
\end{textebox}

\textbf{A. Vocabulary in Context (2.5 marks)}
\textbf{Find in the text words or phrases that mean the same as:}
\begin{enumerate}
\item A fixed regular payment (para 1)
\item Cutting down trees over a large area (para 2)
\item Wearing away of topsoil by wind or rain (para 2)
\item Gave power/authority to someone (para 3)
\item Ability to recover quickly from difficulties (para 5)
\end{enumerate}

\textbf{B. True / False with Justification (5 marks)}
\textbf{For each statement, write T (True) or F (False) and quote the sentence or phrase from the text that justifies your choice.}
\begin{enumerate}
\item Wangari Maathai won the Nobel Prize for Literature.
\item The Green Belt Movement only focused on planting trees.
\item Women benefited economically and socially from the movement.
\item Maathai faced opposition from political figures.
\item The movement has remained limited to Kenya.
\end{enumerate}

\textbf{C. Multiple Choice Questions (2.5 marks)}
\textbf{Choose the best answer (A, B, C or D).}
\begin{enumerate}
\item The main goal of the Green Belt Movement was initially to:
      \begin{enumerate}
      \item Fight for democracy.
      \item Combat deforestation and its effects on rural women.
      \item Oppose the government.
      \item Promote commercial agriculture.
      \end{enumerate}
\item The phrase \eng{“plant the seeds of peace and hope”} is a:
      \begin{enumerate}
      \item Literal description of farming.
      \item Metaphor for long-term positive change.
      \item Political slogan for an election.
      \item Scientific term for reforestation.
      \end{enumerate}
\item According to the text, who suffers most from environmental degradation?
      \begin{enumerate}
      \item Wealthy businessmen.
      \item Politicians.
      \item The poorest, especially women.
      \item Urban men.
      \end{enumerate}
\end{enumerate}

\textbf{D. Reference Questions (2.5 marks)}
\textbf{What do the underlined words in the text refer to?}
\begin{enumerate}
\item \eng{“\textbf{It} empowered women.”} (para 3, line 1)
\item \eng{“By earning income and managing nurseries, \textbf{they} gained confidence…”} (para 3)
\item \eng{“\textbf{She} fought against land grabbing…”} (para 4)
\item \eng{“\textbf{It} proves that local action…”} (para 5, line 1)
\end{enumerate}

\textbf{E. Short Answer Questions (5 marks)}
\textbf{Answer the following questions in your own words as far as possible.}
\begin{enumerate}
\item Why did women have to walk longer distances according to paragraph 2?
\item Besides planting trees, what other skills did the women acquire?
\item What does the author mean by \eng{“community-based conservation”} in the last paragraph?
\item Give two reasons why Wangari Maathai is considered a heroine.
\item What lesson can Cameroon learn from the Green Belt Movement?
\end{enumerate}
\end{exercice}

\begin{exercice}[Traduction : Version (Anglais $\rightarrow$ Français)]
\textbf{Translate the following passage into French.}
\begin{textebox}[Excerpt]
“Climate change is no longer a distant threat; it is a present reality affecting every corner of the globe. In Cameroon, the shrinking of Lake Chad has devastated local economies and forced thousands to migrate. Meanwhile, coastal cities like Limbe and Kribi face rising tides and erosion. The solution lies not only in international agreements but in local adaptation: planting mangroves to protect shores, investing in drought-resistant crops, and educating communities on waste management. Every citizen has a role to play in building a sustainable future.”
\end{textebox}
\end{exercice}

\begin{exercice}[Expression écrite type Baccalauréat]
\textbf{Choose \textbf{ONE} of the following topics and write a composition of about \textbf{150--180 words}.}

\textbf{Topic 1:}
Your school wants to celebrate World Environment Day (June 5th). Write a short article for the school magazine presenting \textbf{one major environmental problem in your town or city} (e.g., plastic waste, erosion, air pollution, deforestation) and proposing \textbf{two concrete solutions} that students and the community can implement. 
\textit{Hints: Describe the problem (causes, effects). Solution 1: Action + Expected result. Solution 2: Action + Expected result. Conclusion: Call to action.}

\textbf{Topic 2:}
“\textbf{We do not inherit the Earth from our ancestors; we borrow it from our children.}” (Native American Proverb). 
Write an argumentative essay explaining what this proverb means in the context of climate change today. Give \textbf{three arguments} to support the idea that the current generation has a moral duty to protect the environment for future generations.
\textit{Hints: Introduction (explain proverb + thesis). Argument 1: Resource depletion. Argument 2: Irreversible damage (biodiversity, climate). Argument 3: Intergenerational justice/equity. Conclusion.}
\end{exercice}

\newpage

\coursec{Corrigés des exercices}

\courssub{Corrigés des exercices}

\begin{corrige}[Exercice 1 — QCM : Compréhension globale]
\textbf{Answers:}
\begin{enumerate}
\item \textbf{B} — Increased concentrations of greenhouse gases. 
      \textit{Justification:} The text states: \eng{“It is mainly caused by increased concentrations of greenhouse gases in the atmosphere…”}
\item \textbf{D} — Using renewable energy. 
      \textit{Justification:} The text lists \eng{“burning fossil fuels, deforestation and industrial farming”}. Renewable energy is a solution, not a cause.
\item \textbf{B} — Melting ice caps. 
      \textit{Justification:} The text explicitly mentions: \eng{“The consequences are already visible: melting ice caps, rising sea levels…”}
\item \textbf{B} — Gradual and continuing over time. 
      \textit{Justification:} \eng{“Long-term”} means over a long period; \eng{“rise”} indicates a continuous increase. It is not sudden (A), immediate (C), or seasonal (D).
\end{enumerate}
\end{corrige}

\begin{corrige}[Exercice 2 — Vrai ou Faux justifié]
\textbf{Answers:}
\begin{enumerate}
\item \textbf{False.} The text says: \eng{“This pollution poses a serious threat to \textbf{marine life}”} (general) and gives the example of sea turtles, not only whales. Quote: \eng{“This pollution poses a serious threat to marine life.”}
\item \textbf{True.} Quote: \eng{“Microplastics, tiny fragments of degraded plastic…”} (\eng{“degraded plastic”} means broken down over time).
\item \textbf{False.} The text says: \eng{“Governments and organisations \textbf{are trying to ban} single-use plastics…”} (present continuous = action in progress, not completed). Quote: \eng{“are trying to ban single-use plastics”}
\item \textbf{True.} Quote: \eng{“Microplastics… have entered the food chain and are now found in the fish we eat.”}
\end{enumerate}
\end{corrige}

\begin{corrige}[Exercice 3 — Vocabulaire : Mots-clés de l'environnement]
\textbf{Matching:}
\begin{center}
\begin{tabularx}{\linewidth}{|X|X|X|X|}
\hline
\rowcolor{popblueL}
\textbf{English} & \textbf{Answer} & \textbf{Français} & \textbf{Letter} \\
\hline
1. Drought & 1--b & Sécheresse & b \\
2. Emissions & 2--d & Émissions & d \\
3. Greenhouse gas & 3--a & Gaz à effet de serre & a \\
4. Deforestation & 4--c & Déforestation & c \\
5. Sustainable & 5--f & Durable / Soutenable & f \\
6. Renewable energy & 6--e & Énergies renouvelables & e \\
7. Carbon footprint & 7--g & Empreinte carbone & g \\
8. Biodiversity & 8--h & Biodiversité & h \\
\hline
\end{tabularx}
\end{center}

\textbf{Personal sentences (examples):}
\begin{enumerate}
\item \eng{The prolonged \textbf{drought} in the North Region has destroyed the maize harvest.} (La sécheresse prolongée dans la région du Nord a détruit la récolte de maïs.)
\item \eng{Factories must reduce their carbon \textbf{emissions} to meet the new standards.} (Les usines doivent réduire leurs émissions de carbone pour respecter les nouvelles normes.)
\item \eng{Using \textbf{sustainable} farming methods protects the soil for future generations.} (L'utilisation de méthodes agricoles durables protège le sol pour les générations futures.)
\item \eng{Cameroon has a rich \textbf{biodiversity} with many endemic species in the Korup National Park.} (Le Cameroun possède une riche biodiversité avec de nombreuses espèces endémiques dans le parc national de Korup.)
\end{enumerate}
\end{corrige}

\begin{corrige}[Exercice 4 — Grammaire : Passif et modaux d'obligation/nécessité]
\textbf{Règle rappelée :} Passif = \eng{be} (conjugué) + \eng{V3 (participe passé)}. Modal + Passif Infinitif = \eng{Modal + be + V3}.
\begin{enumerate}
\item \eng{The forest \textbf{was destroyed} by fire last summer.} (Passé simple passif : action terminée dans le passé).
\item \eng{New laws \textbf{must be implemented} to protect endangered species immediately.} (Modal \eng{must} + passif infinitif : obligation forte).
\item \eng{Plastic bottles \textbf{can be recycled} into fleece jackets.} (Modal \eng{can} + passif infinitif : possibilité/capacité).
\item \eng{The agreement \textbf{will be signed} by all member countries next week.} (Futur passif : \eng{will be + V3}).
\item \eng{Something \textbf{has to be done} to reduce carbon emissions before it is too late.} (Semi-modal \eng{have to} + passif infinitif : nécessité externe).
\item \eng{The river \textbf{had been polluted} by chemical waste for years before the factory closed.} (Plus-que-parfait passif : \eng{had been + V3} — antériorité par rapport à un passé \eng{closed}).
\end{enumerate}
\end{corrige}

\begin{corrige}[Exercice 5 — Traduction : Thème (Français $\rightarrow$ Anglais)]
\textbf{Proposed translations:}
\begin{enumerate}
\item \eng{\textbf{Deforestation in the Amazon threatens biodiversity and accelerates climate change.}} 
      (\textit{Note:} \eng{deforestation} (nom), \eng{threatens} (verbe + s), \eng{accelerates} (verbe + s).)
\item \eng{\textbf{Governments must ban single-use plastics to reduce ocean pollution.}} 
      (\textit{Note:} \eng{single-use plastics} = plastiques à usage unique ; \eng{must ban} = obligation.)
\item \eng{\textbf{If we do not act now, future generations will suffer the consequences of global warming.}} 
      (\textit{Note:} Conditionnelle 1\textsuperscript{re} forme : \eng{If + present, will + BV}. \eng{suffer the consequences} collocation.)
\item \eng{\textbf{Renewable energies such as solar and wind power are being developed in many African countries.}} 
      (\textit{Note:} Passif présent BE+ING : \eng{are being developed} = “sont en train d'être développées”. \eng{such as} = comme.)
\item \eng{\textbf{It is essential that everyone reduce their carbon footprint by changing their consumption habits.}} 
      (\textit{Note:} Subjonctif présent (forme base verbale sans \eng{s}) après \eng{It is essential that}. \eng{reduce} pas \eng{reduces}. \eng{consumption habits} = habitudes de consommation.)
\end{enumerate}
\end{corrige}

\begin{corrige}[Exercice 6 — Transformation de phrases]
\begin{enumerate}
\item \eng{\textbf{Fossil fuels are burned to produce energy.}} (Passif présent simple : \eng{are + V3}. Agent \eng{by people} omis car évident.)
\item \eng{\textbf{The scientist said (that) we must / had to protect the rainforest.}} (Discours rapporté : \eng{must} $\rightarrow$ \eng{had to} ou \eng{must} inchangé si l'obligation persiste. \eng{said that} introduit la subordonnée.)
\item \eng{\textbf{Petrol cars will be banned in 2035 (by the government).}} (Futur passif : \eng{will be + V3}.)
\item \eng{\textbf{Rising sea levels are caused by climate change.}} (Passif présent : \eng{are caused by}.)
\item \eng{\textbf{More trees must / have to be planted.}} (Modal + passif infinitif : \eng{must/have to be + V3}. \eng{It is necessary to} $\rightarrow$ \eng{More trees must be planted}.)
\item \eng{\textbf{In spite of the rain, they planted trees.}} (\eng{In spite of} + nom/Gérondif. \eng{Although it was raining} $\rightarrow$ \eng{In spite of the rain} ou \eng{In spite of raining}.)
\end{enumerate}
\end{corrige}

\begin{corrige}[Exercice 7 — Texte à trous : Cloze test]
\textbf{Completed text:}
\begin{quote}
Cameroon faces several environmental challenges. Irregular \textbf{rainfall} patterns lead to severe \textbf{drought} in the North, affecting agriculture. In cities like Douala, plastic \textbf{waste} clogs drainage systems, causing floods. The government plans to \textbf{ban} non-biodegradable plastic bags. Industries must reduce their carbon \textbf{emissions} to limit the \textbf{greenhouse} effect. Promoting \textbf{sustainable} farming methods is crucial for food security. Citizens are encouraged to \textbf{recycle} paper, glass and metal. Every small action helps fight this global \textbf{threat}.
\end{quote}

\textbf{Word box:} \eng{threat, sustainable, emissions, drought, recycle, greenhouse, waste, ban, rainfall, harmful, pollution}. \\
\textbf{Extra words (not used):} \eng{harmful, pollution}. (Il y a 11 mots dans la boîte pour 9 trous. Les deux mots non utilisés sont \eng{harmful} et \eng{pollution}.)
\end{corrige}

\begin{corrige}[Exercice 8 — Appariement : Causes, Conséquences, Solutions]
\textbf{Answers:}
\begin{enumerate}
\item \textbf{1 -- B -- II} (Burning fossil fuels $\rightarrow$ Air pollution/Smog $\rightarrow$ Switch to solar/wind energy)
\item \textbf{2 -- C -- I} (Deforestation $\rightarrow$ Soil erosion/Desertification $\rightarrow$ Reforestation programmes)
\item \textbf{3 -- D -- III} (Chemical fertilizers $\rightarrow$ Water contamination $\rightarrow$ Organic farming/Composting) 
      \textit{Note:} Chemical fertilizers runoff contaminates water (eutrophication). Solution: Organic farming.
\item \textbf{4 -- D -- IV} (Dumping industrial waste $\rightarrow$ Water contamination $\rightarrow$ Strict fines/Water treatment plants)
      \textit{Note:} Both 3 and 4 lead to Water contamination (D), but 4 matches best with Strict fines/Water treatment (IV) as direct industrial regulation. 3 matches best with Organic farming (III) as agricultural alternative.
\end{enumerate}
\textbf{Corrected Triplets:}
1 -- B -- II \\
2 -- C -- I \\
3 -- D -- III \\
4 -- D -- IV \quad (\textit{Conséquence D partagée, mais solutions distinctes})
\end{corrige}

\begin{corrige}[Exercice 9 — Compréhension écrite type Baccalauréat]
\textbf{Barème indicatif : Total 17.5 marks}

\textbf{A. Vocabulary in Context (2.5 marks — 0.5 each)}
\begin{enumerate}
\item \eng{\textbf{stipend}} (para 1: \eng{“pay rural women a small \textbf{stipend}”})
\item \eng{\textbf{deforestation}} (para 2: \eng{“Kenya faced massive \textbf{deforestation}”})
\item \eng{\textbf{soil erosion}} (para 2: \eng{“\textbf{Soil erosion} dried up rivers”})
\item \eng{\textbf{empowered}} (para 3: \eng{“It \textbf{empowered} women”})
\item \eng{\textbf{resilience}} (para 5: \eng{“build \textbf{resilience}”})
\end{enumerate}

\textbf{B. True / False with Justification (5 marks — 1 each: 0.5 T/F + 0.5 Quote)}
\begin{enumerate}
\item \textbf{F.} Quote: \eng{“Wangari Maathai, the first African woman to win the \textbf{Nobel Peace Prize}…”}
\item \textbf{F.} Quote: \eng{“But the Green Belt Movement did \textbf{more than reforest land}. It empowered women.”}
\item \textbf{T.} Quote: \eng{“By earning income and managing nurseries, they gained confidence and a voice in their communities.”}
\item \textbf{T.} Quote: \eng{“She fought against land grabbing by \textbf{powerful politicians}, was arrested and beaten…”}
\item \textbf{F.} Quote: \eng{“Today, the movement is a model for community-based conservation \textbf{across Africa}.”}
\end{enumerate}

\textbf{C. Multiple Choice Questions (2.5 marks — ~0.83 each)}
\begin{enumerate}
\item \textbf{B} (Combat deforestation and its effects on rural women). Text: \eng{“The idea was simple… pay rural women… to plant trees”} to counter deforestation effects (firewood scarcity).
\item \textbf{B} (Metaphor for long-term positive change). \eng{“Seeds of peace and hope”} symbolises future benefits, not literal seeds only.
\item \textbf{C} (The poorest, especially women). Text: \eng{“Maathai realized that environmental degradation hurt \textbf{the poorest first, especially women}.”}
\end{enumerate}

\textbf{D. Reference Questions (2.5 marks — ~0.625 each)}
\begin{enumerate}
\item \eng{\textbf{It}} = \eng{The movement} (or \eng{The Green Belt Movement}).
\item \eng{\textbf{They}} = \eng{rural women} (or \eng{women}).
\item \eng{\textbf{She}} = \eng{Wangari Maathai}.
\item \eng{\textbf{It}} = \eng{The movement} (or \eng{The Green Belt Movement}).
\end{enumerate}

\textbf{E. Short Answer Questions (5 marks — 1 each)}
\textit{Answers in own words (key ideas required):}
\begin{enumerate}
\item Because \eng{soil erosion dried up rivers} and \eng{firewood became scarce}, so they had to walk further to find it.
\item They acquired \eng{civic education, voting rights knowledge, and sustainable farming techniques} (also: income management, nursery management).
\item It means \eng{conservation efforts led and managed by local communities themselves} (bottom-up approach), not imposed by outside governments or NGOs.
\item (Any two) She founded the Green Belt Movement; she linked environment/democracy/peace; she empowered rural women; she resisted political oppression (arrested/beaten); she won the Nobel Peace Prize; her model spreads across Africa.
\item Cameroon can learn that \eng{involving local communities (especially women) in tree planting} combats desertification (North) and erosion, provides income, and builds climate resilience \eng{through simple, low-cost, participatory actions}.
\end{enumerate}

\textbf{Points de vigilance (Bac) :}
\begin{itemize}
\item Pour le Vocabulaire : recopier \textbf{exactement} le mot du texte (orthographe, majuscule).
\item Pour V/F : \textbf{La citation est obligatoire} sans elle, 0 point même si T/F est bon.
\item Pour Référence : identifier l'antécédent \textbf{dans le texte} (nom ou groupe nominal).
\item Pour Questions ouvertes : répondre \textbf{en anglais}, phrases complètes, \textbf{propres mots} (pas de copier-coller long).
\end{itemize}
\end{corrige}

\begin{corrige}[Exercice 10 — Traduction : Version (Anglais $\rightarrow$ Français)]
\textbf{Proposed translation:}
\begin{quote}
« Le changement climatique n'est plus une menace lointaine ; c'est une réalité présente qui affecte tous les coins du globe. Au Cameroun, le rétrécissement du lac Tchad a dévasté les économies locales et forcé des milliers de personnes à migrer. Pendant ce temps, les villes côtières comme Limbé et Kribi font face à la montée des eaux et à l'érosion. La solution ne réside pas seulement dans les accords internationaux, mais dans l'adaptation locale : planter des mangroves pour protéger les rives, investir dans des cultures résistantes à la sécheresse, et éduquer les communautés sur la gestion des déchets. Chaque citoyen a un rôle à jouer pour construire un avenir durable. »
\end{quote}

\textbf{Notes de traduction :}
\begin{itemize}
\item \eng{distant threat} : menace lointaine / future.
\item \eng{present reality} : réalité actuelle / présente.
\item \eng{shrinking of Lake Chad} : rétrécissement / assèchement du lac Tchad.
\item \eng{devastated local economies} : dévasté / ruiné les économies locales.
\item \eng{forced thousands to migrate} : forcé des milliers (de personnes) à migrer / à fuir.
\item \eng{rising tides and erosion} : montée des marées / des eaux et érosion.
\item \eng{lies not only in… but in} : ne réside pas seulement dans… mais dans.
\item \eng{mangroves} : mangroves / palétuviers.
\item \eng{drought-resistant crops} : cultures résistantes à la sécheresse.
\item \eng{waste management} : gestion des déchets.
\item \eng{sustainable future} : avenir durable / soutenable.
\end{itemize}
\end{corrige}

\begin{corrige}[Exercice 11 — Expression écrite type Baccalauréat]
\textbf{Barème indicatif (Total : 10 marks) :}
\begin{itemize}
\item \textbf{Content / Task Achievement (4 pts) :} Respect du sujet, idées pertinentes, développement des arguments/solutions.
\item \textbf{Organisation \& Coherence (2 pts) :} Structure (Intro/Body/Conclusion), paragraphing, connecteurs logiques (\eng{Firstly, Moreover, However, In conclusion}).
\item \textbf{Grammar \& Vocabulary (3 pts) :} Correction syntaxique (passifs, modaux, conditionnels), lexique environnement précis (\eng{deforestation, emissions, sustainable, resilience, stakeholder}), orthographe.
\item \textbf{Mechanics (1 pt) :} Ponctuation, majuscules, lisibilité, longueur (150--180 mots).
\end{itemize}

\textbf{Modèle de rédaction — Topic 1 (Article for School Magazine)}

\begin{quote}
\textbf{Turning the Tide on Plastic Waste in Douala}

By [Student Name], Form 5

As we celebrate World Environment Day this June 5th, our school community must confront a visible crisis choking our city: plastic waste. In Douala, thousands of tonnes of non-biodegradable bags and bottles are discarded daily. Poor waste collection and the habit of dumping rubbish in streams cause blocked drains, leading to devastating floods during the rainy season. Moreover, this plastic ends up in the Wouri River and the Atlantic, threatening fish stocks and the livelihoods of fishermen in Youpwe.

However, students can drive concrete change. \textbf{Firstly}, we should launch a “Zero Plastic Campus” campaign: ban single-use bottles in school, install water fountains, and distribute reusable bottles. This would drastically reduce our daily plastic footprint and educate peers by example. \textbf{Secondly}, we can organise monthly “Eco-Walks” to clean neighbourhood gutters and riverbanks, partnering with local NGOs like “Game Changers” for recycling. This restores drainage flow and turns waste into resources.

In conclusion, the fight for a cleaner Douala starts with us. Let us not wait for the government alone; every refused plastic bag, every bottle recycled, is a vote for a sustainable future. \textbf{Act today, breathe tomorrow.}

(\textit{Word count: ~175 words})
\end{quote}

\textbf{Modèle de rédaction — Topic 2 (Argumentative Essay)}

\begin{quote}
\textbf{A Debt to the Future: Our Moral Duty on Climate Change}

The Native American proverb, “We do not inherit the Earth from our ancestors; we borrow it from our children,” encapsulates the essence of intergenerational justice. It reminds us that we are not owners of the planet, but temporary stewards. In the context of today’s climate emergency, this moral duty translates into three imperatives.

\textbf{Firstly}, we are depleting finite resources at an unsustainable rate. Fossil fuels, minerals, and ancient aquifers take millennia to form, yet we consume them in decades. If we exhaust them, future generations will inherit a bankrupt planet, forced into energy poverty and resource wars. We must transition to a circular economy now.

\textbf{Secondly}, we are causing irreversible damage to biodiversity and the climate system. Species extinction is forever; melted glaciers do not refreeze easily; crossed tipping points (like permafrost thaw) trigger cascading feedback loops. Our children will face a biologically impoverished, climatically unstable world — a direct theft of their natural heritage.

\textbf{Thirdly}, climate change deepens inequality across time. The wealthy nations and generations that emitted most greenhouse gases will not suffer the worst consequences; the poor and the unborn will. This violates the fundamental principle of equity: we cannot enjoy prosperity today by mortgaging the survival of tomorrow.

In conclusion, honouring this “loan” demands immediate, radical action: net-zero emissions, massive reforestation, and climate finance for the vulnerable. We must return the Earth not as we found it, but better — fertile, stable, and just. That is the only legacy worthy of the name.

(\textit{Word count: ~210 words — slightly over, but acceptable for depth; can be trimmed to 180 by condensing intro/conclusion.})
\end{quote}

\textbf{Points de vigilance (Rédaction Bac) :}
\begin{itemize}
\item \textbf{Respecter le format :} Article (titre, auteur, date, ton engageant) vs Essai (thèse, arguments numérotés, ton formel).
\item \textbf{Connecteurs :} \eng{Furthermore, Consequently, For instance, On the other hand, Ultimately}.
\item \textbf{Grammaire attendue :} Passifs (\eng{is caused by}), Modaux (\eng{must, should, have to}), Conditionnelles (\eng{If we act...}), Relatives (\eng{which, who, where}).
\item \textbf{Vocabulaire précis :} Éviter \eng{good/bad thing} $\rightarrow$ \eng{beneficial/detrimental, sustainable/unsustainable, mitigation/adaptation}.
\item \textbf{Relecture :} Accords (s du pluriel, 3\textsuperscript{e} pers. sing.), temps verbaux, mots-outils (\eng{the, a, an}).
\end{itemize}
\end{corrige}

\newpage

\coursec{Fiche bilan}

\begin{popfigure}
\begin{tikzpicture}[
  mindmap,
  grow cyclic,
  every node/.style={concept, execute at begin node=\hskip0pt},
  concept color=popblue,
  level 1/.append style={level distance=4.5cm, sibling angle=360/7},
  level 2/.append style={level distance=3cm, sibling angle=45},
  text=white,
  font=\small\bfseries,
  popblue/.style={concept color=popblue},
  popgreen/.style={concept color=popgreen},
  poporange/.style={concept color=poporange},
  poppurple/.style={concept color=poppurple},
  poppink/.style={concept color=poppink},
  popgold/.style={concept color=popgold},
  popdark/.style={concept color=popdark},
]
\node[popblue, minimum size=2.2cm, align=center]{Environmental Issues\\Climate Change}
  child[popgreen] { node[minimum size=1.6cm, align=center]{Key Vocabulary\\Global warming, footprint, renewables} }
  child[poporange] { node[minimum size=1.6cm, align=center]{Reading Strategies\\Skimming, Scanning, Guessing} }
  child[poppurple] { node[minimum size=1.6cm, align=center]{Grammar Focus\\Passive Voice, Modals, Connectors} }
  child[poppink] { node[minimum size=1.6cm, align=center]{Text Structure\\Cause/Effect, Problem/Solution} }
  child[popgold] { node[minimum size=1.6cm, align=center]{Writing Summary\\Main ideas, Synthesis, Own words} }
  child[popdark] { node[minimum size=1.6cm, align=center]{Speaking Opinion\\Arguments, Solutions, Cameroon context} }
  child[popblue!70!popgreen] { node[minimum size=1.6cm, align=center]{Exam Skills\\Time mgmt, Question types, Keywords} };
\end{tikzpicture}
\legende{Carte mentale récapitulative : Lecture — Enjeux environnementaux et changement climatique}
\end{popfigure}

\begin{bilan}

\end{bilan}

\courssub{Lexique clé — Environment \& Climate Change (Anglais / Français)}

\begin{center}
\begin{tabularx}{\linewidth}{|>{\bfseries}X|X|X|}
\hline
\rowcolor{popblueL} \textbf{English Term} & \textbf{Français} & \textbf{Collocations / Contexte} \\
\hline
Global warming & Réchauffement climatique & \eng{global warming effects}, \eng{to combat global warming} \\
\hline
Climate change & Changement climatique & \eng{climate change mitigation}, \eng{climate change denial} \\
\hline
Greenhouse effect & Effet de serre & \eng{greenhouse gases (GHG)}, \eng{enhanced greenhouse effect} \\
\hline
Carbon footprint & Empreinte carbone & \eng{reduce one's carbon footprint}, \eng{carbon footprint calculator} \\
\hline
Renewable energy & Énergie renouvelable & \eng{solar/wind/hydro power}, \eng{transition to renewables} \\
\hline
Fossil fuels & Combustibles fossiles & \eng{burning fossil fuels}, \eng{phase out fossil fuels} \\
\hline
Deforestation & Déforestation & \eng{illegal logging}, \eng{deforestation in the Congo Basin} \\
\hline
Biodiversity & Biodiversité & \eng{loss of biodiversity}, \eng{protect biodiversity} \\
\hline
Sustainable development & Développement durable & \eng{UN Sustainable Development Goals (SDGs)}, \eng{sustainable agriculture} \\
\hline
Emissions & Émissions (gaz) & \eng{carbon emissions}, \eng{emission reduction targets} \\
\hline
Heatwave & Canicule & \eng{severe heatwave}, \eng{heatwave warning} \\
\hline
Drought & Sécheresse & \eng{prolonged drought}, \eng{drought-resistant crops} \\
\hline
Flooding & Inondation & \eng{flash flooding}, \eng{flood defences} \\
\hline
Melting ice caps & Fonte des calottes glaciaires & \eng{rising sea levels}, \eng{Arctic ice melt} \\
\hline
Eco-friendly / Green & Écologique, vert & \eng{eco-friendly products}, \eng{green technology} \\
\hline
\end{tabularx}
\end{center}

\vspace{0.5cm}

\courssub{Points de grammaire indispensables (à connaître par cœur)}

\begin{center}
\begin{tabularx}{\linewidth}{|>{\bfseries}p{3.5cm}|X|X|}
\hline
\rowcolor{popblueL} \textbf{Notion} & \textbf{Règle / Forme} & \textbf{Exemple (Environment)} \\
\hline
\textbf{Passive Voice}\newline(Voix passive) & \eng{be (conjugué) + V-en (participe passé)}\newline Focus on action, not agent. Agent (\eng{by...}) often omitted in scientific texts. & \eng{CO2 \textbf{is released} by factories.}\newline\eng{Forests \textbf{are being destroyed} at an alarming rate.}\newline\eng{New laws \textbf{have been passed} to protect wildlife.} \\
\hline
\textbf{Modals of Obligation}\newline\& Advice & \eng{must / have to} (obligation forte, loi)\newline\eng{should / ought to} (conseil, recommandation)\newline\eng{had better} (avertissement urgent) & \eng{Governments \textbf{must reduce} emissions.}\newline\eng{We \textbf{should invest} in solar energy.}\newline\eng{You \textbf{had better save} water now.} \\
\hline
\textbf{Modals of Probability}\newline (Prediction) & \eng{will / won't} (certitude)\newline\eng{may / might / could} (possibilité)\newline\eng{must / can't} (déduction logique) & \eng{Temperatures \textbf{will rise} by 2°C.}\newline\eng{It \textbf{might rain} later (uncertain).}\newline\eng{The glaciers \textbf{must be melting} fast (evidence).} \\
\hline
\textbf{Connectors}\newline(Logique texte) & \textbf{Cause:} \eng{because (of), due to, since, as}\newline\textbf{Consequence:} \eng{therefore, consequently, as a result, thus}\newline\textbf{Contrast:} \eng{however, whereas, although, despite/in spite of}\newline\textbf{Addition:} \eng{furthermore, moreover, in addition} & \eng{\textbf{Due to} deforestation, \textbf{consequently} soil erodes.}\newline\eng{\textbf{Although} renewables are cheap, \textbf{however} we use coal.}\newline\eng{\textbf{Moreover}, the IPCC warns of irreversible damage.} \\
\hline
\textbf{Relative Clauses}\newline(Propositions relatives) & \eng{who} (personnes), \eng{which/that} (choses), \eng{where} (lieu), \eng{whose} (possession).\newline Defining (no comma) vs Non-defining (commas). & \eng{Scientists \textbf{who} study climate agree.}\newline\eng{The Amazon, \textbf{which} is a carbon sink, burns.}\newline\eng{Cameroon, \textbf{where} biodiversity is high, acts.} \\
\hline
\end{tabularx}
\end{center}

\vspace{0.5cm}

\courssub{Méthodes express de lecture (Reading Strategies)}

\begin{methode}[Skimming — Lecture rapide pour l'idée générale]
\begin{enumerate}
\item Lire le \textbf{titre}, le \textbf{sous-titre}, les \textbf{images/légendes}.
\item Lire la \textbf{1ère phrase} de chaque paragraphe (topic sentence).
\item Repérer les \textbf{mots-clés} répétés (ex. \eng{climate, warming, CO2, rise}).
\item Répondre : \textit{« De quoi parle le texte en une phrase ? »}
\end{enumerate}
\end{methode}

\begin{methode}[Scanning — Recherche d'informations précises]
\begin{enumerate}
\item Identifier le \textbf{mot-clé} de la question (chiffre, nom propre, date, terme technique).
\item Balayer le texte \textbf{verticalement} (yeux en diagonale) sans lire chaque mot.
\item Localiser le mot-clé (ou synonyme) dans le texte.
\item Lire \textbf{autour} (1 phrase avant/après) pour extraire la réponse exacte.
\end{enumerate}
\end{methode}

\begin{methode}[Déduire le sens des mots inconnus (Guessing Meaning)]
\begin{enumerate}
\item \textbf{Contexte immédiat :} Regarder la phrase avant/après (explications, exemples, contrastes).
\item \textbf{Formation du mot :} Préfixes (\eng{re-, de-, un-, non-, over-, under-}), suffixes (\eng{-tion, -ment, -able, -ity, -ize}), racines latines/grecques.
\item \textbf{Classe grammaticale :} Nom, verbe, adjectif, adverbe (position dans la phrase).
\item \textbf{Cognats / Faux-amis :} Vérifier si le mot ressemble au français (\eng{environment $\rightarrow$ environnement} vs \eng{actually $\neq$ actuellement}).
\item \textbf{Substitution :} Remplacer par un synonyme probable et vérifier si la phrase a du sens.
\end{enumerate}
\end{methode}

\vspace{0.3cm}

\courssub{Structure type d'un résumé (Summary Writing)}
\begin{itemize}
\item \textbf{Introduction (1 phrase) :} Auteur / Source + Sujet principal + Thèse centrale. \eng{In the article ``...'', X argues that...}
\item \textbf{Corps (3-4 phrases max) :} Idées principales \textbf{dans l'ordre du texte}, regroupées. Utiliser \textbf{connecteurs} (\eng{Firstly, Furthermore, However, Finally}). \textbf{Passif} et \textbf{discours indirect} privilégiés. \textbf{Propre vocabulaire} (pas de copier-coller).
\item \textbf{Conclusion (1 phrase) :} Message final / Appel à l'action / Conséquence future.
\end{itemize}

\courssub{Expressions pour donner son avis (Giving Opinion)}
\begin{center}
\begin{tabularx}{\linewidth}{|X|X|}
\hline
\rowcolor{popblueL} \textbf{Introduire} & \textbf{Argumenter / Nuancer} \\
\hline
\eng{In my opinion / view,...} & \eng{Firstly,... / Moreover,... / Additionally,...} \\
\eng{I strongly believe that...} & \eng{For instance,... / A striking example is...} \\
\eng{As far as I am concerned,...} & \eng{However,... / On the other hand,...} \\
\eng{It seems to me that...} & \eng{Despite this,... / Although..., ...} \\
\hline
\end{tabularx}
\end{center}



\begin{aretenir}[Checklist avant le Bac — Lesson 31 : Environmental Reading]
$\square$ Je sais faire du \textbf{skimming} pour trouver l'idée générale en moins de 2 minutes. \\
$\square$ Je sais faire du \textbf{scanning} pour localiser une date, un chiffre, un nom propre ou un terme technique. \\
$\square$ Je connais par cœur les \textbf{20 mots de vocabulaire clé} (anglais $\leftrightarrow$ français + collocations). \\
$\square$ Je reconnais et je sais former la \textbf{voix passive} (présent, passé, présent perfect, modaux) pour décrire des processus scientifiques. \\
$\square$ J'utilise correctement \textbf{must/have to} (obligation) et \textbf{should/ought to} (recommandation) pour proposer des solutions. \\
$\square$ J'utilise \textbf{will/may/might/could} pour exprimer des prédictions sur le climat avec le bon degré de certitude. \\
$\square$ Je maîtrise les \textbf{connecteurs logiques} (cause, conséquence, concession, addition) pour structurer mon résumé. \\
$\square$ Je sais rédiger un \textbf{résumé en 5-6 phrases} (intro, idées clés ordonnées, conclusion) \textbf{sans copier le texte}. \\
$\square$ Je sais exprimer mon \textbf{opinion personnelle} avec 2 arguments structurés et un exemple concret (Cameroun ou international). \\
$\square$ Je gère mon temps : \textbf{10 min lecture}, \textbf{15 min questions}, \textbf{10 min résumé/opinion}, \textbf{5 min relecture}. \\
$\square$ Je repère les \textbf{faux-amis} fréquents (ex. \eng{actually, eventually, sensible, pollution vs pollution}) dans le contexte environnemental.
\end{aretenir}

\findecours

\end{document}
